\documentclass{fkt-notes}
\usepackage{macros}
\usepackage{multicol}
\usepackage{array}
\usepackage{etoolbox}
\usepackage{needspace}
\usepackage{colortbl}
\usepackage{makeidx}
\makeindex
\usepackage[most]{tcolorbox}

\course{Diskret støddfrøði}
\coursecode{005019.23}
\semester{Fall 2026}
\author{Oddmar Dam}
\title{Discrete Mathematics: Exam Companion}
\date{\today}
\DeclareMathOperator{\lcm}{lcm}

% Exam companion for Diskret støddfrøði (005019.23): the lecture material
% (slides, Fall 2026), the textbook and my own notes merged
% into one printable document, ordered by Kolman chapter, with one index.
% Meant to travel to the MC exams as permitted notes (notes/books allowed,
% no electronics).
%
%   * Yellow shading = rules, laws and theorems.
%   * "Exam method" subsections come from my own notes: step-by-step ways to
%     attack the typical MC question types, with worked examples.
%   * "Not on the slides" = taken from Kolman; "Slide erratum" = a slide
%     mistake, with the textbook's correct version; "(answer not on
%     slides)" = a worked answer the slides left blank.
%   * The multiplication table and the list of primes are the last page
%     before the index.

% A bit more on the page than the class default -- this is for printing.
\newgeometry{margin=2cm,headheight=14pt}
\savegeometry{refbase}
\setlist{itemsep=1pt,topsep=3pt}

\newtcolorbox{keybox}[1][]{colback=white,colframe=black,boxrule=0.6pt,
  arc=2pt,left=4pt,right=4pt,top=3pt,bottom=3pt,fonttitle=\bfseries,
  coltitle=black,colbacktitle=black!10,#1}
% Rules, laws and theorems are highlighted in yellow: rulebox = a framed
% box (the old keybox, shaded); rulepar = shading only, for rules that sit
% in running text or in a table.
\definecolor{rulefill}{RGB}{255,244,150}
\newtcolorbox{rulebox}[1][]{enhanced,colback=rulefill,colframe=black,boxrule=0.6pt,
  arc=2pt,left=4pt,right=4pt,top=3pt,bottom=3pt,fonttitle=\bfseries,
  coltitle=black,colbacktitle=rulefill!60!orange!45,#1}
\newtcolorbox{rulepar}{enhanced,breakable,colback=rulefill,colframe=rulefill,
  boxrule=0pt,arc=1pt,left=3pt,right=3pt,top=2pt,bottom=2pt,
  before skip=4pt,after skip=4pt}
\newcommand{\erratum}[1]{\par\noindent\textbf{Slide erratum.} #1\par}
\newcommand{\notslides}[1]{\par\noindent\textbf{Not on the slides.} #1\par}
\newcommand{\src}[1]{\par\noindent{\small\itshape Source: #1}\par\smallskip}

% Numbering follows Kolman: each \part is a textbook chapter, and each
% section carries its textbook section number, set by hand with \kolman
% because the slides do not cover every section (e.g.\ 4.6, 4.8). So
% "6.3 Lattices" has subsections 6.3.1, 6.3.2, ..., like the exercises.
\renewcommand{\partname}{Chapter}
\renewcommand{\thepart}{\arabic{part}}
\renewcommand{\thesection}{\arabic{part}.\arabic{section}}
\renewcommand{\theHsection}{\arabic{part}.\arabic{section}}
\newcommand{\kolman}[1]{\setcounter{section}{\numexpr#1-1\relax}}
% Article's TOC leaves 1.5em for a section number -- too narrow for "6.3".
\makeatletter
\patchcmd{\l@section}{1.5em}{2.3em}{}{}
\renewcommand*\l@subsection{\@dottedtocline{2}{2.3em}{3.2em}}
\makeatother
% Kolman doesn't number definitions, and "Definition 6.3.1" would read as
% subsection 6.3.1, so they are unnumbered here.
\theoremstyle{definition}
\newtheorem*{defn}{Definition}
% keep every heading together with the text that follows it
\pretocmd{\section}{\Needspace{9\baselineskip}}{}{}
\pretocmd{\subsection}{\Needspace{6\baselineskip}}{}{}

% Hasse diagrams: pt = a plain dot (label it with label=...), dv = a
% circled number for the divisor lattices.
\tikzset{
  pt/.style={circle,fill,inner sep=1.2pt},
  dv/.style={circle,draw,minimum size=6mm,inner sep=0pt,font=\small},
  hasse/.style={x=0.8cm,y=0.8cm,every label/.style={font=\small}},
  % labels that may sit across an edge (the cube-shaped B_3 diagrams)
  boxed/.style={every label/.style={font=\small,fill=white,inner sep=0.8pt}},
}
% Logic diagrams (6.5): US-style gate symbols.
\usetikzlibrary{circuits.logic.US}
\tikzset{gates/.style={circuit logic US,font=\small,
  every circuit symbol/.style={thick},logic gate input sep=0.32cm}}
% Hasse-style diagrams drawn with plain text labels (from my own notes).
\tikzset{nhasse/.style={every node/.style={inner sep=1.5pt}, line width=0.5pt}}
% Trees (7.1-7.3): tn = a circled vertex; ex = a bare label (expression trees).
% Undirected graphs (7.4-7.5): pt dots with labels; w = an edge weight on
% the edge; mst = an edge picked for the spanning tree.
\tikzset{
  tn/.style={circle,draw,minimum size=6mm,inner sep=0pt,font=\small},
  ex/.style={inner sep=1.5pt,font=\small},
  ug/.style={every label/.style={font=\small},line width=0.5pt},
  w/.style={midway,fill=white,inner sep=1pt,font=\scriptsize},
  mst/.style={line width=1.8pt},
}

\begin{document}
\maketitle

\begin{keybox}[title={Notation bridge: slides vs.\ Kolman vs.\ past MC papers}]
  The slides write $\neg p$, $p\to q$, $p\leftrightarrow q$. Kolman and the
  past MC papers write $\sim p$, $p\Rightarrow q$, $p\Leftrightarrow q$ ---
  same meaning. On the slides $\N=\{1,2,3,\dots\}$; the MC papers and some
  books include $0$. If a question hinges on it, check which convention it
  uses.
\end{keybox}

\begin{rulepar}\noindent\textbf{Yellow shading} marks the rules, laws and theorems.\end{rulepar}

\noindent\textbf{Exam method} subsections are step-by-step recipes for the typical MC question types (own notes). \textbf{Not on the slides} = from Kolman; \textbf{Slide erratum} = a mistake on a slide. The multiplication table and a list of primes are on the last page before the index.

\tableofcontents

% =====================================================================
\newpage
\part{Fundamentals}

\kolman{1}
\section{Sets}\src{Kolman 1.1}
\index{set}\index{subset}\index{power set}\index{empty set}\index{universal set}\index{cardinality}\index{number sets}\index{natural numbers}\index{integers}%

\begin{defn}
  A \textbf{set} is a well-defined collection of objects: for any object
  you can decide whether it is in the set. Order and repetition don't
  matter: $\{1,2,3\}=\{3,2,1\}$. We write $x\in A$ or $x\notin A$.
\end{defn}

\begin{center}
  \begin{tabular}{@{}ll@{}}
    \toprule
    $\Z^{+}$         & positive integers $\{1,2,3,\dots\}$                    \\
    $\N$             & natural numbers $\{1,2,3,\dots\}$ (slides' convention) \\
    $\Z$             & integers $\{\dots,-2,-1,0,1,2,\dots\}$                 \\
    $\Q$             & rationals $\{a/b \mid a,b\in\Z,\ b\neq0\}$             \\
    $\R$             & real numbers                                           \\
    $\emptyset$, $U$ & empty set, universal set                               \\
    \bottomrule
  \end{tabular}
  \qquad
  $\Z^{+}\subseteq\N\subset\Z\subset\Q\subset\R$
\end{center}

\textbf{Terms.} $A\subseteq B$: every element of $A$ is in $B$.
$A\nsubseteq B$: at least one element of $A$ is not in $B$.
$A\cap B=\emptyset$: disjoint. $|A|$ is the \textbf{cardinality} (number
of elements). The \textbf{power set} $P(A)$ is the set of all subsets,
including $\emptyset$ and $A$, and $|P(A)|=2^{|A|}$.

\begin{example}
  $A=\{1,2,3\}$: $|A|=3$,
  $P(A)=\{\emptyset,\{1\},\{2\},\{3\},\{1,2\},\{1,3\},\{2,3\},\{1,2,3\}\}$,
  $|P(A)|=2^3=8$.
\end{example}
\textbf{Set-builder notation.} $P(x)$ is a sentence about $x$ (a
\textbf{predicate}); $\{x\mid P(x)\}$ is read ``the set of all $x$ such that
$P(x)$''. E.g.\ $\{x\mid x\in\Z\text{ and }x^2<12\}=\{-3,-2,\dots,3\}$.
\index{set-builder notation}\index{predicate}\index{element of ($\in$)}\index{Venn diagram}%

\textbf{Venn diagram:} the universal set $U$ is a rectangle, a set $A$ a
circle inside it.
\begin{center}
  \begin{tikzpicture}
    \draw (-2,-1.2) rectangle (2,1.2);
    \node at (-1.7,0.9) {$U$};
    \draw (0,0) circle (0.8cm);
    \node at (0,0) {$A$};
  \end{tikzpicture}
\end{center}


\kolman{2}
\section{Operations on sets}\src{Kolman 1.2}
\index{union}\index{intersection}\index{set difference}\index{symmetric difference}\index{complement!of a set}%
\begin{align*}
  A\cup B      & = \{x : x\in A \vee x\in B\}      &  & \text{union (OR)}                 \\
  A\cap B      & = \{x : x\in A \wedge x\in B\}    &  & \text{intersection (AND)}         \\
  A-B          & = \{x : x\in A \wedge x\notin B\} &  & \text{difference}                 \\
  A\oplus B    & = (A-B)\cup(B-A)                  &  & \text{symmetric difference (XOR)} \\
  \overline{A} & = U-A                             &  & \text{complement}
\end{align*}

\begin{example}
  $A=\{a,b,c,e,f\}$, $B=\{b,d,r,s\}$: $A\cup B=\{a,b,c,d,e,f,r,s\}$.
  $A=\{a,b,c,e,f\}$, $B=\{b,e,f,r,s\}$: $A\cap B=\{b,e,f\}$.
\end{example}

\Needspace{9\baselineskip}
\subsection{Set laws (Theorem 1.2.1)}
\index{set laws}\index{De Morgan's laws!sets}\index{distributive law!sets}\index{commutative law}\index{associative law}\index{idempotent law}\index{domination law}\index{identity law}%
\begin{rulepar}
  \begin{center}
    \renewcommand{\arraystretch}{1.15}
    \begin{tabular}{@{}lll@{}}
      \toprule
      Commutative  & $A\cup B=B\cup A$                                 & $A\cap B=B\cap A$                                 \\
      Associative  & $A\cup(B\cup C)=(A\cup B)\cup C$                  & $A\cap(B\cap C)=(A\cap B)\cap C$                  \\
      Distributive & $A\cap(B\cup C)=(A\cap B)\cup(A\cap C)$           & $A\cup(B\cap C)=(A\cup B)\cap(A\cup C)$           \\
      Identity     & $A\cup\emptyset=A$                                & $A\cap U=A$                                       \\
      Domination   & $A\cup U=U$                                       & $A\cap\emptyset=\emptyset$                        \\
      Idempotent   & $A\cup A=A$                                       & $A\cap A=A$                                       \\
      Complement   & $A\cup\overline{A}=U$                             & $A\cap\overline{A}=\emptyset$                     \\
                   & $\overline{\emptyset}=U$                          & $\overline{U}=\emptyset$                          \\
      De Morgan    & $\overline{A\cup B}=\overline{A}\cap\overline{B}$ & $\overline{A\cap B}=\overline{A}\cup\overline{B}$ \\
      \bottomrule
    \end{tabular}
  \end{center}
  De Morgan: when the complement moves inside, $\cup\leftrightarrow\cap$
  swap.
\end{rulepar}

\subsection{Inclusion--exclusion}
\index{inclusion--exclusion}\index{counting!size of a union}%
\begin{rulebox}
  \vspace{-0.8em}
  \begin{align*}
    |A\cup B|       & = |A|+|B|-|A\cap B|                                             \\
    |A\cup B\cup C| & = |A|+|B|+|C| - |A\cap B|-|A\cap C|-|B\cap C| + |A\cap B\cap C|
  \end{align*}
  Add singles, subtract pairwise overlaps, add back the triple overlap.
\end{rulebox}

\begin{example}[1.2 Ex.~7]
  $A=\{a,b,c,d,e\}$, $B=\{c,e,f,h,k,m\}$: $|A\cup B|=5+6-2=9$.
\end{example}
\begin{example}[1.2 Ex.~8]
  $A=\{a,b,c,d,e\}$, $B=\{a,b,e,g,h\}$, $C=\{b,d,e,g,h,k,m,n\}$:
  $|A\cup B\cup C|=5+5+8-3-2-4+2=11$.
\end{example}
\begin{example}[1.2 Ex.~10, commuter survey]
  $|B|=30$, $|T|=35$, $|A|=100$, $|B\cap T|=15$, $|B\cap A|=15$,
  $|T\cap A|=20$, $|B\cap T\cap A|=5$. Respondents
  $=30+35+100-15-15-20+5=120$ (answer not on slides).
\end{example}

% =====================================================================
\kolman{3}
\section{Sequences and characteristic functions}\src{Kolman 1.3}
\index{sequence}\index{characteristic function}\index{explicit formula}%

A \textbf{sequence} is a list in a definite order; unlike a set, order
and repeats matter. $(1,2,3,2,2,3,1)$ is a sequence; swapping two
elements gives a different sequence. It can be finite or infinite.

\begin{itemize}
  \item \textbf{Explicit formula}: $a_n$ computed directly from $n$, e.g.\
        $a_n=2n$ gives $2,4,6,8,\dots$
  \item \textbf{Recursive formula}: $a_n$ from earlier terms, e.g.\
        $a_n=a_{n-1}+a_{n-2}$ gives $1,1,2,3,5,8,\dots$
\end{itemize}

\begin{defn}
  The \textbf{characteristic function} of $A\subseteq U$ is
  $f_A(x)=1$ if $x\in A$, $0$ if $x\notin A$.
\end{defn}
\begin{rulebox}[title=Theorem 1.3.1]
  $f_{A\cap B}=f_Af_B$ \qquad
  $f_{A\cup B}=f_A+f_B-f_Af_B$ \qquad
  $f_{A\oplus B}=f_A+f_B-2f_Af_B$
\end{rulebox}
\begin{example}[1.3 Ex.~12]
  $U=\{1,\dots,6\}$, $A=\{1,2\}$, $B=\{2,4,6\}$, $C=\{4,5,6\}$ as bit
  strings: $f_A=(1,1,0,0,0,0)$, $f_B=(0,1,0,1,0,1)$,
  $f_C=(0,0,0,1,1,1)$.
\end{example}

% =====================================================================
\kolman{4}
\section{Properties of integers}\src{Kolman 1.4}
\index{integers!properties}\index{divisibility}\index{prime}%

\begin{rulebox}[title=Theorem 1.4.1 (division with remainder)]
  \index{division with remainder}\index{quotient}\index{remainder}\index{dividend and divisor}%
  For integers $m$, $n$ with $n>0$ there are unique $q,r$ with
  $m=qn+r$, $0\le r<n$. ($m$ dividend, $n$ divisor, $q$ quotient,
  $r$ remainder.) \quad E.g.\ $17=5\cdot3+2$.
\end{rulebox}

\index{modulo}%
The remainder is \emph{never} negative --- careful with negative $n$:

\begin{example}
  $47 = 6 \cdot 7 + 5$, so $q = 6$, $r = 5$.\\
  $-47 = (-7) \cdot 7 + 2$, so $q = -7$, $r = 2$ (not $q = -6$, $r = -5$).\\
  $5 = 0 \cdot 7 + 5$: when $n < m$ the quotient is $0$ and $r = n$.
\end{example}

\noindent $m \mid n$ exactly when the remainder is $0$. In code this is
\texttt{n \% m} --- Python's \texttt{\%} follows the same rule
(\texttt{-47 \% 7 == 2}), C\# and C do not (\texttt{-47 \% 7 == -5}).

\textbf{Divisibility.} $n\mid m$ if $m=qn$ for some integer $q$.
$4\mid24$, $5\mid135$, $4\nmid25$.

\begin{rulepar}
  \textbf{Theorem 1.4.2.} For integers $a,b,c$:
  $a\mid b \wedge a\mid c\Rightarrow a\mid(b+c)$ and $a\mid(b-c)$;\;
  $a\mid b\Rightarrow a\mid bc$;\;
  $a\mid b\wedge b\mid c\Rightarrow a\mid c$.
\end{rulepar}

\textbf{Primes.} $p>1$ whose only positive divisors are $1$ and $p$:
$2,3,5,7,11,13,17,19,23,29,\dots$ (note that $1$ is \emph{not} prime).

\begin{rulepar}
  \textbf{Theorem 1.4.3 (unique factorisation).} Every $n>1$ is uniquely
  \index{unique factorisation}\index{fundamental theorem of arithmetic}\index{prime factorization}%
  $p_1^{k_1}p_2^{k_2}\cdots p_s^{k_s}$ with primes $p_1<\dots<p_s$, $k_i\ge1$.
  \quad $30=2\cdot3\cdot5$, \quad $1890=2\cdot3^3\cdot5\cdot7$.
\end{rulepar}
\subsection{Divisibility rules (no calculator)}
\index{divisibility rules}\index{divides ($\mid$)}$a \mid b$ means ``$a$ divides $b$'': $b = a \cdot k$ for some integer $k$.

\begin{rulepar}
  \begin{center}
    \renewcommand{\arraystretch}{1.25}
    \begin{tabular}{@{}c p{5.2cm} p{6.4cm}@{}}
      \toprule
      $d$ & $d \mid n$ when \ldots                                           & Why it works \\
      \midrule
      2   & last digit is even
          & $10$ is a multiple of $2$, so only the last digit matters                       \\
      3   & digit sum is a multiple of $3$
          & $10, 100, 1000, \ldots$ are each $1$ more than a multiple of $3$
      ($9, 99, 999, \ldots$)                                                                \\
      4   & last two digits form a multiple of $4$
          & $100$ is a multiple of $4$                                                      \\
      5   & last digit is $0$ or $5$
          & $10$ is a multiple of $5$                                                       \\
      6   & divisible by both $2$ and $3$
          & $6 = 2 \cdot 3$ and $\gcd(2,3) = 1$                                             \\
      7   & chop off the last digit $b$, subtract $2b$ from what is left;
      the result is a multiple of $7$
          & $n = 10a + b$ and $n - 21b = 10(a - 2b)$, with $7 \mid 21$ and
      $7 \nmid 10$                                                                          \\
      8   & last three digits form a multiple of $8$
          & $1000$ is a multiple of $8$                                                     \\
      9   & digit sum is a multiple of $9$
          & same reason as $3$                                                              \\
      10  & last digit is $0$
          &                                                                                 \\
      11  & alternating digit sum (from the right: $+ - + \cdots$) is
      $0$ or a multiple of $11$
          & $10, 1000, \ldots$ are $1$ \emph{less} than a multiple of $11$,
      $100, 10000, \ldots$ are $1$ more, so the signs alternate                             \\
      \bottomrule
    \end{tabular}
  \end{center}

  \noindent \textbf{Combining rules:} if $\gcd(m, k) = 1$ then
  $mk \mid n \iff m \mid n \text{ and } k \mid n$. So $12 \mid n$ iff
  $3 \mid n$ and $4 \mid n$ --- but $2$ and $6$ do \emph{not} give $12$,
  since $\gcd(2, 6) \neq 1$.
\end{rulepar}

\subsection{Prime factorization by trial division (exam method)}
\index{trial division}\index{prime factorization!by trial division}\begin{enumerate}[label=\arabic*.]
  \item Try the primes in increasing order: $2, 3, 5, 7, 11, \ldots$
  \item When a prime $p$ divides the current number, write $p$ down,
        divide, and try the \emph{same} $p$ again on the quotient.
  \item When $p$ no longer divides, move on to the next prime.
  \item Stop when the quotient is $1$, or when $p^2$ exceeds the
        quotient --- the quotient is then itself prime.
\end{enumerate}

\noindent \textbf{Why stop at $\sqrt{n}$:} if $n = ab$ with
$1 < a \leq b$, then $a \cdot a \leq a \cdot b = n$, so $a \leq \sqrt{n}$.
A composite number always has a prime factor $\leq \sqrt{n}$; if no
prime up to $\sqrt{n}$ divides $n$, then $n$ is prime. Trying primes
in increasing order also means every factor found is automatically
prime --- any smaller factor would already have been divided out.

\noindent \textbf{Layout:} write it as a ``division ladder'', the number
on the left and the prime that divides it on the right.

\begin{example}[$495$]
  \leavevmode\par\noindent
  \begin{minipage}[t]{0.25\textwidth}
    \vspace{0pt}
    \begin{tabular}{r|l}
      495 & 3  \\
      165 & 3  \\
      55  & 5  \\
      11  & 11 \\
      1   &
    \end{tabular}
  \end{minipage}
  \begin{minipage}[t]{0.7\textwidth}
    \vspace{0pt}
    \begin{itemize}[nosep]
      \item $2$: last digit $5$ is odd --- no.
      \item $3$: $4 + 9 + 5 = 18$ --- yes, $495 = 3 \cdot 165$.
      \item $3$ again: $1 + 6 + 5 = 12$ --- yes, $165 = 3 \cdot 55$.
      \item $3$ again: $5 + 5 = 10$ --- no, move on.
      \item $5$: last digit $5$ --- yes, $55 = 5 \cdot 11$.
      \item $11$: next prime $7$ has $7^2 = 49 > 11$, so $11$ is prime. Stop.
    \end{itemize}
  \end{minipage}
  \[
    495 = 3^2 \cdot 5 \cdot 11
  \]
  Check by multiplying back: $9 \cdot 5 = 45$, $45 \cdot 11 = 495$ \checkmark
\end{example}

\begin{example}[$1001$]
  $2$: odd. $3$: digit sum $2$. $5$: last digit $1$. $7$: $100 - 2\cdot 1
    = 98 = 7 \cdot 14$ --- yes, $1001 = 7 \cdot 143$. On $143$: $7$ again
  gives $14 - 2 \cdot 3 = 8$, no. $11$: $3 - 4 + 1 = 0$ --- yes,
  $143 = 11 \cdot 13$, and $13$ is prime.
  \[
    1001 = 7 \cdot 11 \cdot 13
  \]
\end{example}

\subsection{\texorpdfstring{Is $n$ prime? (exam method)}{Is n prime?}}
\index{prime!testing (square-root rule)}Test only the primes $p \leq \sqrt{n}$. If none divides $n$, then $n$ is
prime.

\begin{example}
  $211$: $\sqrt{211} \approx 14.5$ (since $14^2 = 196$, $15^2 = 225$), so
  test $2, 3, 5, 7, 11, 13$. Odd; digit sum $4$; last digit $1$;
  $21 - 2 = 19$; $1 - 1 + 2 = 2$; $13 \cdot 16 = 208$, remainder $3$.
  None divide, so $211$ is \textbf{prime}.\\
  $221$: same primes to test, and $13 \cdot 17 = 221$, so it is
  \textbf{not} prime --- and $13 \leq \sqrt{221} \approx 14.9$ as promised.
\end{example}

\subsection{Number of divisors}
\index{number of divisors}\begin{rulepar}
  If $n = p_1^{e_1} \cdots p_k^{e_k}$, every divisor of $n$ is
  $p_1^{f_1} \cdots p_k^{f_k}$ with $0 \leq f_i \leq e_i$: $e_i + 1$
  choices for each exponent, so by the multiplication principle
  \[
    \text{number of positive divisors of } n = (e_1 + 1)(e_2 + 1) \cdots (e_k + 1).
  \]
\end{rulepar}

\begin{example}
  $495 = 3^2 \cdot 5^1 \cdot 11^1$ has $(2+1)(1+1)(1+1) = 12$ divisors:
  $1$, $3$, $5$, $9$, $11$, $15$, $33$, $45$, $55$, $99$, $165$, $495$.
\end{example}


\subsection{GCD and LCM}
\index{greatest common divisor (gcd)}\index{least common multiple (lcm)}%
$\gcd(a,b)$ is the largest $k$ with $k\mid a$ and $k\mid b$.

\begin{rulepar}
  \textbf{Theorem 1.4.4.} If $d=\gcd(a,b)$ then (1) $d=sa+tb$ for some
  integers $s,t$; (2) $d\mid a$ and $d\mid b$; (3) any common divisor $c$ of
  $a$ and $b$ satisfies $c\mid d$.
\end{rulepar}

\begin{keybox}[title=Euclidean algorithm]
  \index{Euclidean algorithm}\index{greatest common divisor (gcd)!Euclidean algorithm}%
  \begin{minipage}[t]{0.52\linewidth}
    \begin{tabular}[t]{@{}cccc@{}}
      \toprule
      $a$ & $b$ & $q$ & $r$ \\\midrule
      190 & 34  & 5   & 20  \\
      34  & 20  & 1   & 14  \\
      20  & 14  & 1   & 6   \\
      14  & 6   & 2   & 2   \\
      6   & 2   & 3   & 0   \\\bottomrule
    \end{tabular}
    \quad $\gcd(190,34)=2$
  \end{minipage}\hfill
  \begin{minipage}[t]{0.40\linewidth}
    \begin{lstlisting}[numbers=none,frame=none]
gcd(a, b):
  while b != 0:
    r = a % b
    a = b
    b = r
  return a
\end{lstlisting}
  \end{minipage}
\end{keybox}
Also $\gcd(12,30)=6$ (answer not on slides).
If $\gcd(a,b)=1$, then $a$ and $b$ are \textbf{relatively prime} (no common prime factor).\index{relatively prime}%

\paragraph{Method 1 --- prime factorizations.} Line up the factorizations
over the same primes (exponent $0$ where a prime is missing) and take the
\textbf{smallest} exponent of each prime:
\[
  \gcd(a, b) = \prod p_i^{\min(e_i, f_i)}
  \qquad \text{where } a = \prod p_i^{e_i},\ b = \prod p_i^{f_i}.
\]

\begin{example}
  $495 = 3^2 \cdot 5^1 \cdot 7^0 \cdot 11^1 \cdot 13^0$ and
  $1001 = 3^0 \cdot 5^0 \cdot 7^1 \cdot 11^1 \cdot 13^1$.
  Only $11$ appears in both, so $\gcd(495, 1001) = 11$.
\end{example}

\paragraph{Method 2 --- Euclidean algorithm.}\index{Euclidean algorithm}\index{greatest common divisor (gcd)!Euclidean algorithm} No factoring needed, so it
is the fast method for large numbers. It rests on
\[
  \gcd(a, b) = \gcd(b, r) \quad \text{where } a = qb + r,
  \qquad \text{and} \qquad \gcd(a, 0) = a.
\]
Divide the larger by the smaller, replace the pair by (smaller,
remainder), and repeat. The last \emph{nonzero} remainder is the gcd. It
is a loop:
\texttt{while b != 0: a, b = b, a \% b}, then the answer is \texttt{a}.

\begin{example}[$\gcd(120, 84)$]
  \begin{align*}
    120 & = 1 \cdot 84 + 36 \\
    84  & = 2 \cdot 36 + 12 \\
    36  & = 3 \cdot 12 + 0
  \end{align*}
  The last nonzero remainder is $12$, so $\gcd(120, 84) = 12$.
\end{example}

\paragraph{Writing the gcd as $sa + tb$.}\index{greatest common divisor (gcd)!written as $sa+tb$} There are always integers $s, t$
with $\gcd(a, b) = sa + tb$. To find them, run Euclid, then solve each
line for its remainder and substitute backwards, starting from the line
that produced the gcd:

\begin{example}[continued]
  From the lines above: $12 = 84 - 2 \cdot 36$ and $36 = 120 - 1 \cdot 84$.
  \begin{align*}
    12 & = 84 - 2 \cdot 36           \\
       & = 84 - 2\,(120 - 84)        \\
       & = 3 \cdot 84 - 2 \cdot 120.
  \end{align*}
  So $s = -2$, $t = 3$ with $a = 120$, $b = 84$. Check:
  $252 - 240 = 12$ \checkmark. (Keep $84$ and $120$ as symbols while
  substituting --- only collect terms at the end.)
\end{example}

\begin{rulepar}
  \paragraph{Useful gcd rules.}
  \begin{itemize}[nosep]
    \item $\gcd(a, b) = \gcd(b, a)$ and $\gcd(a, 0) = a$.
    \item $\gcd(a, b) = \gcd(a, b - ka)$ for any integer $k$ --- the
          Euclidean step.
    \item If $p$ is prime: $\gcd(p, n) = p$ if $p \mid n$, otherwise $1$.
    \item Consecutive integers are always relatively prime:
          $\gcd(n, n+1) = 1$.
  \end{itemize}
\end{rulepar}


$\lcm(a,b)$ is the smallest positive $k$ with $a\mid k$ and $b\mid k$.
\begin{rulebox}[title=Theorem 1.4.6]
  \index{gcd times lcm equals product@$\gcd\cdot\lcm=ab$}%
  $\gcd(a,b)\cdot\lcm(a,b)=ab$.\quad
  Proof idea: $a=dx$, $b=dy$ with $\gcd(x,y)=1$, so $\lcm(a,b)=dxy$ and
  $d\cdot dxy=(dx)(dy)=ab$.
\end{rulebox}
\paragraph{Method 1 --- prime factorizations.} Take the \textbf{largest}
exponent of each prime:
\[
  \lcm(a, b) = \prod p_i^{\max(e_i, f_i)}.
\]

\begin{example}
  $120 = 2^3 \cdot 3 \cdot 5$ and $84 = 2^2 \cdot 3 \cdot 7$:
  $\lcm = 2^3 \cdot 3 \cdot 5 \cdot 7 = 840$ and
  $\gcd = 2^2 \cdot 3 = 12$, matching Euclid.
\end{example}

\begin{rulepar}
  \paragraph{Method 2 --- from the gcd.}
  \[
    \gcd(a, b) \cdot \lcm(a, b) = a \cdot b
    \qquad \Longrightarrow \qquad
    \lcm(a, b) = \frac{a}{\gcd(a, b)} \cdot b.
  \]
  This holds because $\min(e, f) + \max(e, f) = e + f$ for every prime's
  exponents. Divide \emph{before} multiplying to keep the numbers small.
\end{rulepar}

\begin{example}
  $\lcm(495, 1001) = \dfrac{495}{11} \cdot 1001 = 45 \cdot 1001 = 45045$.
\end{example}

\noindent \textbf{Special cases:} if $\gcd(a, b) = 1$ then
$\lcm(a, b) = ab$; if $a \mid b$ then $\gcd(a, b) = a$ and
$\lcm(a, b) = b$.


\subsection{\texorpdfstring{Mod-$n$}{Mod-n} and number bases}
\index{modulo}\index{mod function}\index{number bases}\index{base conversion}\index{binary}\index{hexadecimal}%
$f_n(z)=r$ where $z=qn+r$, $0\le r<n$ (i.e.\ $z \bmod n$).
$f_3(14)=2$ since $14=3\cdot4+2$; \; $f_7(153)=6$ since $153=7\cdot21+6$.

\begin{rulepar}
  \textbf{Theorem 1.4.7.} For any base $b>1$ every positive $n$ is uniquely
  $n=d_kb^k+d_{k-1}b^{k-1}+\dots+d_1b+d_0$ with $0\le d_i<b$.
  Convert by repeated division by $b$; the remainders, read bottom-up, are
  the digits.
\end{rulepar}

\begin{example}[1.4 Ex.~11, answers not on slides]
  $158=2\cdot64+1\cdot16+3\cdot4+2=2132_4$.\quad
  $215=1\cdot125+3\cdot25+3\cdot5+0=1330_5$.
\end{example}
\paragraph{Base $10 \to$ base $b$.} Divide by $b$ repeatedly, writing down
the remainders, until the quotient is $0$. The remainders read
\textbf{bottom-up} are the digits (the first remainder is the last digit).

\begin{example}[$495$ to hex]
  \begin{align*}
    495 & = 30 \cdot 16 + 15 \quad (\mathrm{F}) \\
    30  & = 1 \cdot 16 + 14 \quad (\mathrm{E})  \\
    1   & = 0 \cdot 16 + 1
  \end{align*}
  Bottom-up: $495 = (\mathrm{1EF})_{16}$.
\end{example}

\paragraph{Base $b \to$ base 10.} Either sum $d_i b^i$, or (fewer
multiplications) go left to right: start at $0$, and for each digit
multiply by $b$ and add the digit --- Horner's rule\index{Horner's rule}, a single loop.

\begin{example}
  $(\mathrm{1EF})_{16}$: $\;0 \cdot 16 + 1 = 1$, $\;1 \cdot 16 + 14 = 30$,
  $\;30 \cdot 16 + 15 = 495$ \checkmark\\
  $(2101)_3 = 2 \cdot 27 + 1 \cdot 9 + 0 \cdot 3 + 1 = 64$.
\end{example}

\paragraph{Binary $\leftrightarrow$ hex shortcut.} $16 = 2^4$, so each hex
digit is exactly four bits. Group the bits from the \emph{right} in
fours:
\[
  (\mathrm{1EF})_{16} = \underbrace{0001}_{1}\;\underbrace{1110}_{\mathrm{E}}\;
  \underbrace{1111}_{\mathrm{F}} = (111101111)_2.
\]
The same trick works for octal with groups of three ($8 = 2^3$).

\begin{example}[MC practice: $(123)_{16}$ in base 8]
  \index{base conversion!hex to octal}%
  \textbf{Via binary} (no arithmetic): $1\to0001$, $2\to0010$, $3\to0011$,
  so $(123)_{16}=(000100100011)_2=(100100011)_2$. Regroup from the right in
  threes: $100\,|\,100\,|\,011 = 4,4,3$, i.e.\ $(443)_8$.\\
  \textbf{Check via decimal:} $(123)_{16}=256+32+3=291$, and
  $291=4\cdot64+4\cdot8+3$ \checkmark. The distractors are the same digits
  shuffled, so a reversed grouping (starting from the left) lands on one.
\end{example}

\paragraph{Binary addition.} Add column by column from the right, carrying
when a column reaches $2$ (as $1+1=10_2$):\index{binary addition}
\[
  \begin{array}{r}
    1011 \\ {+}\,1100\\\hline 10111
  \end{array}
  \qquad(11+12=23)
\]


% =====================================================================
\kolman{5}
\section{Matrices}\src{Kolman 1.5}
\index{matrix}\index{matrix multiplication}\index{identity matrix}\index{transpose of a matrix}%

An $m\times n$ matrix has $m$ rows and $n$ columns. \textbf{Square}:
$m=n$. \textbf{Diagonal}: all off-diagonal entries $0$. \textbf{Identity}
$I_n$: diagonal with $1$s. \textbf{Transpose}: $A^{T}=[a_{ji}]$ (swap
rows and columns).

\begin{center}
  \begin{tabular}{@{}ll@{}}
    \toprule
    Addition (Thm 1)       & $A+B=B+A$, \ $(A+B)+C=A+(B+C)$, \ $A+0=A$           \\
    Multiplication (Thm 2) & $A(BC)=(AB)C$, \ $A(B+C)=AB+AC$, \ $(A+B)C=AC+BC$   \\
                           & \textbf{but generally $AB\neq BA$}                  \\
    Transpose (Thm 3)      & $(A^T)^T=A$, \ $(A+B)^T=A^T+B^T$, \ $(AB)^T=B^TA^T$ \\
    \bottomrule
  \end{tabular}
\end{center}

\begin{example}[1.5 Ex.~7, answer not on slides]
  \[
    \begin{bmatrix}2&3&-4\\1&2&3\end{bmatrix}
    \begin{bmatrix}3&1\\-2&2\\5&-3\end{bmatrix}
    =\begin{bmatrix}-20&20\\14&-4\end{bmatrix}
  \]
  Counterexample to commutativity: $A=\begin{bmatrix}1&0\\0&0\end{bmatrix}$,
  $B=\begin{bmatrix}0&1\\0&0\end{bmatrix}$ give $AB=B$ but $BA=0$.
\end{example}

\begin{rulebox}[title=Inverse (Theorem 1.5.4)]
  \index{inverse!of a matrix}\index{determinant formula ($2\times2$ inverse)}%
  $A^{-1}$ satisfies $AA^{-1}=A^{-1}A=I$. For
  $A=\begin{bmatrix}a&b\\c&d\end{bmatrix}$ with $ad-bc\neq0$:\quad
  $A^{-1}=\dfrac{1}{ad-bc}\begin{bmatrix}d&-b\\-c&a\end{bmatrix}$.
\end{rulebox}

\subsection{Boolean (0/1) matrices}
\index{Boolean matrix}\index{join!of matrices}\index{meet!of matrices}\index{Boolean product of matrices}%
\begin{itemize}
  \item \textbf{Join} $A\vee B$: element-wise OR.
  \item \textbf{Meet} $A\wedge B$: element-wise AND.
  \item \textbf{Boolean product} $A\odot B$: like ordinary multiplication,
        with AND for $\times$ and OR for $+$:
        $c_{ij}=(a_{i1}\wedge b_{1j})\vee\dots\vee(a_{ik}\wedge b_{kj})$.
        Needs columns of $A$ = rows of $B$.
\end{itemize}
\begin{rulepar}
  Theorem 1.5.5: $\vee$ and $\wedge$ are commutative and associative;
  $\odot$ is associative (not commutative).
\end{rulepar}


\begin{example}
  \[A=\begin{bmatrix}1&0\\0&1\end{bmatrix},\ B=\begin{bmatrix}0&1\\1&0\end{bmatrix}:\qquad
    A\vee B=\begin{bmatrix}1&1\\1&1\end{bmatrix},\
    A\wedge B=\begin{bmatrix}0&0\\0&0\end{bmatrix},\
    A\odot B=\begin{bmatrix}0&1\\1&0\end{bmatrix}.\]
\end{example}
\begin{example}[1.5 Ex.~12, answers not on slides]
  \[
    A=\begin{bmatrix}1&0&1\\0&1&1\\1&1&0\\0&0&0\end{bmatrix},\quad
    B=\begin{bmatrix}1&1&0\\1&0&1\\0&0&1\\1&1&0\end{bmatrix}:\qquad
    A\vee B=\begin{bmatrix}1&1&1\\1&1&1\\1&1&1\\1&1&0\end{bmatrix},\quad
    A\wedge B=\begin{bmatrix}1&0&0\\0&0&1\\0&0&0\\0&0&0\end{bmatrix}.
  \]
\end{example}
\begin{example}
  \[
    \begin{bmatrix}1&0&1\\0&1&1\end{bmatrix}\odot
    \begin{bmatrix}0&1\\1&1\\1&0\end{bmatrix}
    =\begin{bmatrix}1&1\\1&1\end{bmatrix},
    \qquad c_{11}=(1\wedge0)\vee(0\wedge1)\vee(1\wedge1)=1\text{, likewise the rest.}
  \]
\end{example}

% =====================================================================
\newpage
\part{Logic}

\kolman{1}
\section{Propositions and logical operations}\src{Kolman 2.1}
\index{proposition}\index{truth table}\index{negation}\index{conjunction}\index{disjunction}%

A \textbf{proposition} is a declarative sentence that is either true or
false. ``$2+3=5$'' (T) and ``7 is even'' (F) are propositions; ``Close the
door'' (command), ``What time is it?'' (question) and ``$x+2=5$'' (open
sentence, $x$ unknown) are not.

\begin{center}
  \begin{tabular}{@{}cc|c|c|c|c|c@{}}
    $p$ & $q$ & $\neg p$ & $p\wedge q$ & $p\vee q$ & $p\to q$ & $p\leftrightarrow q$ \\\hline
    T   & T   & F        & T           & T         & T        & T                    \\
    T   & F   & F        & F           & T         & F        & F                    \\
    F   & T   & T        & F           & T         & T        & F                    \\
    F   & F   & T        & F           & F         & T        & T
  \end{tabular}
\end{center}
$p\to q$ is false \emph{only} when a true $p$ gives a false $q$.

\begin{example}[2.1 Ex.~5]\leavevmode\\[2pt]
  \begin{tabular}[t]{@{}cc|cc|c@{}}
    $p$ & $q$ & $p\wedge q$ & $\neg p$ & $(p\wedge q)\vee\neg p$ \\\hline
    T   & T   & T           & F        & T                       \\ T&F&F&F&F\\ F&T&F&T&T\\ F&F&F&T&T
  \end{tabular}
  \\[2pt]
  Similarly $\neg(p\vee q)\wedge r$ is true only in the row $p=F$, $q=F$, $r=T$.
\end{example}

% =====================================================================
\kolman{2}
\section{Conditional statements and equivalence}\src{Kolman 2.2}
\index{conditional statement}\index{converse}\index{inverse!of a conditional}\index{contrapositive}\index{biconditional}\index{tautology}\index{contradiction}\index{contingency}\index{logical equivalence}%

\begin{keybox}[title={The conditional family ($p$: it is raining, $q$: I get wet)}]
  \begin{tabular}{@{}lll@{}}
    Conditional    & $p\to q$             & If it rains, I get wet.                \\
    Converse       & $q\to p$             & If I get wet, it is raining.           \\
    Inverse        & $\neg p\to\neg q$    & If it is not raining, I don't get wet. \\
    Contrapositive & $\neg q\to\neg p$    & If I don't get wet, it is not raining. \\
    Biconditional  & $p\leftrightarrow q$ & It rains if and only if I get wet.
  \end{tabular}\\[3pt]
  Only the \textbf{contrapositive} is equivalent to the original:
  $p\to q\equiv\neg q\to\neg p$. (Converse $\equiv$ inverse, but neither
  $\equiv$ the original.)
\end{keybox}

\begin{example}[MC practice: contrapositive of $A\ne B\Rightarrow A\cap B=\emptyset$]
  \index{contrapositive!of a statement about sets}%
  Recipe: \textbf{swap the two sides, negate both}. Here $p$ is
  ``$A\ne B$'' and $q$ is ``$A\cap B=\emptyset$'', so
  $\neg q$ is ``$A\cap B\ne\emptyset$'' and $\neg p$ is ``$A=B$'':
  \[A\cap B\ne\emptyset\;\Rightarrow\;A=B.\]
  The distractors are the converse ($A\cap B=\emptyset\Rightarrow A\ne B$:
  swapped, not negated), the inverse ($A=B\Rightarrow A\cap B\ne\emptyset$:
  negated, not swapped), and $\wedge$/$\vee$ forms, which are not
  conditionals at all. (Negating $=$ gives $\ne$ and vice versa; it is the
  sets' relation that flips, never the $\cap$ or $\cup$.)
\end{example}

\textbf{Tautology}: true in every row ($p\vee\neg p$).
\textbf{Contradiction}: false in every row ($p\wedge\neg p$).
\textbf{Contingency}: sometimes true, sometimes false ($p\to q$).

\begin{defn}
  $P$ and $Q$ are \textbf{logically equivalent}, $P\equiv Q$, if
  $P\leftrightarrow Q$ is a tautology --- i.e.\ their truth-table columns
  are identical in every row.
\end{defn}

\Needspace{9\baselineskip}
\subsection{Theorem 2.2.1 (fundamental equivalences)}
\index{logical equivalences}\index{De Morgan's laws!logic}\index{distributive law!logic}\index{double negation}%
\begin{rulepar}
  \begin{center}
    \renewcommand{\arraystretch}{1.15}
    \begin{tabular}{@{}lll@{}}
      \toprule
      Commutative     & $p\vee q\equiv q\vee p$                          & $p\wedge q\equiv q\wedge p$                        \\
      Associative     & $p\vee(q\vee r)\equiv(p\vee q)\vee r$            & $p\wedge(q\wedge r)\equiv(p\wedge q)\wedge r$      \\
      Distributive    & $p\vee(q\wedge r)\equiv(p\vee q)\wedge(p\vee r)$ & $p\wedge(q\vee r)\equiv(p\wedge q)\vee(p\wedge r)$ \\
      Idempotent      & $p\vee p\equiv p$                                & $p\wedge p\equiv p$                                \\
      Double negation & $\neg(\neg p)\equiv p$                           &                                                    \\
      De Morgan       & $\neg(p\vee q)\equiv\neg p\wedge\neg q$          & $\neg(p\wedge q)\equiv\neg p\vee\neg q$            \\
      \bottomrule
    \end{tabular}
  \end{center}
\end{rulepar}

\Needspace{9\baselineskip}
\subsection{Theorem 2.2.2 (conditionals)}
\index{conditional!equivalent forms}\index{negation of a conditional}%
\begin{rulebox}
  \begin{tabular}{@{}ll@{}}
    (a) $p\to q\equiv\neg p\vee q$                                         &
    (b) $p\to q\equiv\neg q\to\neg p$                                        \\
    (c) $p\leftrightarrow q\equiv(p\to q)\wedge(q\to p)$                   &
    (d) $\neg(p\to q)\equiv p\wedge\neg q$                                   \\
    (e) $\neg(p\leftrightarrow q)\equiv(p\wedge\neg q)\vee(q\wedge\neg p)$ &
    also $p\leftrightarrow q\equiv(p\wedge q)\vee(\neg p\wedge\neg q)$
  \end{tabular}
\end{rulebox}
\erratum{The slide lists (b) as $p\to q\equiv\neg(q\to\neg p)$, which is
  false ($\neg(q\to\neg p)\equiv p\wedge q$), and omits (a) and (d). The
  versions above are Kolman's. (d) --- the negation of a conditional is
  $p\wedge\neg q$, \emph{not} $\neg p\to\neg q$ --- is a recurring MC trap.}

\Needspace{9\baselineskip}
\subsection{Theorem 2.2.3 (quantifiers)}
\index{quantifier}\index{negating a quantified statement}\index{for all ($\forall$)}\index{there exists ($\exists$)}%
\begin{rulepar}
  \begin{tabular}{@{}ll@{}}
    (a) $\neg(\forall x\,P(x))\equiv\exists x\,\neg P(x)$                        &
    (b) $\neg(\exists x\,P(x))\equiv\forall x\,\neg P(x)$                                                                          \\
    (c) $\exists x\,(P(x)\to Q(x))\equiv\forall x\,P(x)\to\exists x\,Q(x)$       &
    (d) $\exists x\,(P(x)\vee Q(x))\equiv\exists x\,P(x)\vee\exists x\,Q(x)$                                                       \\
    (e) $\forall x\,(P(x)\wedge Q(x))\equiv\forall x\,P(x)\wedge\forall x\,Q(x)$ &                                                 \\
    \multicolumn{2}{@{}l}{(f) $(\forall x\,P(x))\vee(\forall x\,Q(x))\to\forall x\,(P(x)\vee Q(x))$ is a tautology (one-way only)} \\
    \multicolumn{2}{@{}l}{(g) $\exists x\,(P(x)\wedge Q(x))\to\exists x\,P(x)\wedge\exists x\,Q(x)$ is a tautology (one-way only)}
  \end{tabular}

  Negating a quantified statement: flip $\forall\leftrightarrow\exists$ and
  negate the inside.
\end{rulepar}
\erratum{The slide omits (a), prints (d) twice, and gives\\
  ``$(\exists x\,P(x))\to(\forall x\,Q(x))\equiv\forall x\,(P(x)\to Q(x))$'',
  which is false: take domain $\{1,2\}$ with $P,Q$ true only at $1$ --- the
  right side is true, the left is $T\to F$. The list above is Kolman's.}

\Needspace{9\baselineskip}
\newpage
\subsection{Theorem 2.2.4 (tautologies behind the inference rules)}
\index{inference rules}\index{modus ponens}\index{modus tollens}\index{disjunctive syllogism}\index{hypothetical syllogism}%
\begin{rulepar}
  \begin{tabular}{@{}ll@{}}
    (a) $(p\wedge q)\to p$                                          & (b) $(p\wedge q)\to q$  \\
    (c) $p\to(p\vee q)$                                             & (d) $q\to(p\vee q)$     \\
    (e) $\neg p\to(p\to q)$                                         & (f) $\neg(p\to q)\to p$ \\
    (g) $(p\wedge(p\to q))\to q$ \quad\emph{modus ponens}           &
    (h) $(\neg p\wedge(p\vee q))\to q$ \quad\emph{disjunctive syllogism}                      \\
    (i) $(\neg q\wedge(p\to q))\to\neg p$ \quad\emph{modus tollens} &
    (j) $((p\to q)\wedge(q\to r))\to(p\to r)$ \quad\emph{hyp.\ syllogism}
  \end{tabular}
\end{rulepar}

(i) and (j) are from Kolman; the slides list (a)--(h), with
$(p\to q)\to(\neg q\to\neg p)$ as their (f) (also a tautology).

\begin{example}[verifying (g)]\leavevmode\\[2pt]
  \begin{tabular}[t]{@{}cc|cc|c@{}}
    $p$ & $q$ & $p\to q$ & $p\wedge(p\to q)$ & $(p\wedge(p\to q))\to q$ \\\hline
    T   & T   & T        & T                 & T                        \\ T&F&F&F&T\\ F&T&T&F&T\\ F&F&T&F&T
  \end{tabular}
  \quad Last column all T $\Rightarrow$ tautology.
\end{example}

\begin{example}[MC practice: which tautology is used in ``Mum Kathy is a stone; a stone cannot fly; $\therefore$ Mum Kathy cannot fly''?]
  \index{inference rules!identifying from an English argument}%
  \textbf{Method:} turn the sentences into symbols, then match the
  \emph{shape} (premises above the line, conclusion below) to (a)--(j).
  Let $p$ = ``Mum Kathy is a stone'' and $q$ = ``Mum Kathy cannot fly''.
  The second premise, ``a stone cannot fly'', is a general rule applied to
  her: $p\to q$. So the argument is
  \[
    \begin{array}{c} p\\ p\to q\\\hline \therefore q\end{array}
  \]
  which is \textbf{modus ponens}, tautology (g):
  $(p\wedge(p\to q))\to q$. (Look for the option where the premises are
  joined by $\wedge$ and the conclusion after $\to$ is a bare letter that
  appeared as the \emph{consequent} of the conditional.) Wrong shapes to
  rule out: (h) has a negated premise and a $\vee$; (i) concludes a
  \emph{negation}; (j) chains two conditionals.
\end{example}

% =====================================================================
\kolman{3}
\section{Methods of proof}\src{Kolman 2.3}
\index{proof}\index{proof!direct}\index{proof!by contrapositive}\index{proof!by contradiction}\index{valid argument form}%

A proof shows that from premises $p_1,\dots,p_n$ the conclusion $q$ must
follow, i.e.\ $(p_1\wedge\dots\wedge p_n)\to q$ is a tautology. Valid
argument forms:
\begin{rulepar}
  \[
    \begin{array}{c} p\\ p\to q\\\hline \therefore q\end{array}
    \ \text{(modus ponens)}\qquad
    \begin{array}{c} \neg q\\ p\to q\\\hline \therefore \neg p\end{array}
    \ \text{(modus tollens)}\qquad
    \begin{array}{c} p\to q\\ q\to r\\\hline \therefore p\to r\end{array}
    \ \text{(hypothetical syllogism)}
  \]
\end{rulepar}

\begin{keybox}[title=Three proof methods]
  \begin{tabular}{@{}lp{0.7\linewidth}@{}}
    \textbf{Direct}         & To show $p\to q$: assume $p$, derive $q$.      \\
    \textbf{Contrapositive} & Prove $\neg q\to\neg p$ instead (equivalent).  \\
    \textbf{Contradiction}  & To show $p$: assume $\neg p$, derive something
    of the form $r\wedge\neg r$; so $p$ must hold.
  \end{tabular}
\end{keybox}

\begin{example}[$n^2$ even $\iff$ $n$ even]
  ($\Rightarrow$ direction by contrapositive) $n$ odd $\Rightarrow$ $n=2k+1$
  $\Rightarrow$ $n^2=4k^2+4k+1=2(2k^2+2k)+1$ odd.
  ($\Leftarrow$ direct) $n=2k$ $\Rightarrow$ $n^2=4k^2=2(2k^2)$ even.
\end{example}
\begin{example}[direct: odd $+$ odd $=$ even]
  $a=2m+1$, $b=2n+1$ $\Rightarrow$ $a+b=2(m+n+1)$, even.
\end{example}
\begin{example}[contradiction: $\sqrt2$ is irrational]
  Suppose $\sqrt2=p/q$ with $\gcd(p,q)=1$. Then $2q^2=p^2$, so $p^2$ is
  even, so $p$ is even: $p=2k$. Then $2q^2=4k^2$, $q^2=2k^2$, so $q$ is
  even too --- contradicting $\gcd(p,q)=1$.
\end{example}

% =====================================================================
\kolman{4}
\section{Mathematical induction}\src{Kolman 2.4}
\index{mathematical induction}\index{induction!base case}\index{induction!hypothesis}\index{induction!strong}%

\begin{keybox}[title={Template (same shape as a recursive function: base case + recursive step)}]
  \begin{enumerate}[label=\arabic*.,leftmargin=*]
    \item \textbf{Base case}: show $P(n_0)$ (usually $n_0=1$ or $0$).
    \item \textbf{Induction hypothesis}: assume $P(k)$ for some $k\ge n_0$.
    \item \textbf{Induction step}: using the hypothesis, prove $P(k+1)$.
    \item \textbf{Conclude}: by induction, $P(n)$ holds for all $n\ge n_0$.
  \end{enumerate}
  \textbf{Strong induction}: in step 2 assume \emph{all} of
  $P(n_0),\dots,P(k)$. Needed when $P(k+1)$ depends on several earlier
  cases --- and then you need as many base cases as the recurrence looks
  back.
\end{keybox}

\begin{example}[2.4 Ex.~1: $1+2+\dots+n=\frac{n(n+1)}{2}$]
  Base: $1=\frac{1\cdot2}{2}$. Step:
  $1+\dots+k+(k+1)=\frac{k(k+1)}{2}+(k+1)=\frac{k(k+1)+2(k+1)}{2}
    =\frac{(k+1)(k+2)}{2}$.
\end{example}

\begin{example}[2.4 Ex.~2: De Morgan for $n$ sets,
    $\overline{\bigcup_{i=1}^{n}A_i}=\bigcap_{i=1}^{n}\overline{A_i}$]
  Base $n=1$ trivial. Step:
  $\overline{\bigcup_{i=1}^{k+1}A_i}
    =\overline{\bigl(\bigcup_{i=1}^{k}A_i\bigr)\cup A_{k+1}}
    =\overline{\bigcup_{i=1}^{k}A_i}\cap\overline{A_{k+1}}
    =\bigl(\bigcap_{i=1}^{k}\overline{A_i}\bigr)\cap\overline{A_{k+1}}
    =\bigcap_{i=1}^{k+1}\overline{A_i}$
  (two-set De Morgan, then the hypothesis).
\end{example}

\begin{example}[2.4 Ex.~4: $n!\ge2^{n-1}$]
  Base: $1!=1=2^0$. Step: $(k+1)!=(k+1)k!\ge(k+1)2^{k-1}\ge2\cdot2^{k-1}=2^k$
  (since $k+1\ge2$).
\end{example}

\begin{example}[strong induction]
  $a_1=2$, $a_2=4$, $a_3=8$, $a_{n+3}=a_{n+2}+a_{n+1}+a_n$. Claim
  $a_n\le2^n$. Base: $P(1),P(2),P(3)$ (three, because the recurrence looks
  back three). Step ($k\ge3$):
  $a_{k+1}=a_k+a_{k-1}+a_{k-2}\le2^k+2^{k-1}+2^{k-2}=7\cdot2^{k-2}<8\cdot2^{k-2}=2^{k+1}$.
\end{example}
\erratum{The conclusion slide says $a_n<2^n$; the claim is
  $a_n\le2^n$ (equality holds at $n=1,2,3$).}

\kolman{5}
\section{From pattern to conjecture to proof}\src{Kolman 2.5}
\index{conjecture}%
A \textbf{conjecture} is a statement believed true from a pattern but not
yet proved: examples $\to$ pattern $\to$ conjecture $\to$ proof.

\begin{example}[2.5 Ex.~2: $\sum_{m=0}^{n}2^m=2^{n+1}-1$]
  Pattern: $1,3,7,15,\dots$ Base $n=0$: $1=2^1-1$. Step:
  $S_{k+1}=S_k+2^{k+1}=(2^{k+1}-1)+2^{k+1}=2^{k+2}-1$.
\end{example}
\begin{example}[divisibility: $16\mid(19\cdot9^n-8n-3)$, $n\ge0$]
  Base: $19-3=16$. Step:
  $19\cdot9^{k+1}-8(k+1)-3=9(19\cdot9^k-8k-3)+64k+16
    =9\cdot(\text{multiple of }16)+16(4k+1)$.
\end{example}
\begin{example}[practice: $1+3+\dots+(2n-1)=n^2$]
  Base: $1=1^2$. Step: $k^2+(2k+1)=(k+1)^2$.
\end{example}

% =====================================================================
\newpage
\part{Counting}

\kolman{1}
\section{Permutations and combinations}\src{Kolman 3.1--3.2; own notes}
\subsection{Multiplication principle}
\index{multiplication principle}\index{addition principle}If a task is done in steps, with $n_1$ ways to do the first step,
$n_2$ ways for the second, and so on, independently of each other,
the number of ways to do the whole task is
$n_1 \cdot n_2 \cdots n_k$. It is nested loops: the inner body runs
$n_1 \cdot n_2 \cdots n_k$ times. ``This \emph{and} that'' $\to$
multiply; ``this \emph{or} that'' (one or the other, not both)
$\to$ add.

\begin{example}[MC practice: license plates, 2 letters from $A$--$J$ (repeats allowed) then 3 distinct digits]
  \index{multiplication principle!license plates}%
  Fill the five slots left to right and count the choices per slot:
  letters $10\cdot10$ (repeats allowed, so the pool never shrinks);
  digits $10\cdot9\cdot8$ (no repeats, so each pick removes one digit and
  order matters, i.e.\ $P(10,3)$). Total
  $10^2\cdot720=\mathbf{72\,000}$. The letter part and the digit part are
  independent, hence multiplied.
\end{example}

\subsection{Factorial}
\index{factorial}$n! = n \cdot (n-1) \cdots 2 \cdot 1$, with $0! = 1$. E.g.\
$4! = 4 \cdot 3 \cdot 2 \cdot 1 = 24$: the number of ways to put $n$
things in a row.

\subsection{Working with factorials without a calculator (exam method)}
\index{factorial!unrolling}\index{factorial!matching MC options}%
\textbf{Key idea:} $n! = n \cdot (n-1)!$ --- the recursive definition,
\texttt{fact(n) = n * fact(n-1)}. You may stop unrolling anywhere and
leave the rest as a factorial; all of these are the same number:
\[
  11! = 11 \cdot 10! = 11 \cdot 10 \cdot 9! = 11 \cdot 10 \cdot 9 \cdot 8 \cdot 7!
\]

\begin{enumerate}[label=\arabic*.]
  \item \textbf{Unroll to the right place:} down to the factorial the
        MC options end with (e.g.\ $7!$), or to the biggest factorial
        in the denominator.
  \item \textbf{Cancel} the denominator against the unrolled numbers.
  \item \textbf{No direct match?} Split the denominator into small
        factors (primes) and cancel them piece by piece.
  \item \textbf{Compare options} by rewriting them all with the same
        tail factorial --- then only the numbers in front matter; you
        never need the value of $7!$.
\end{enumerate}

\begin{example}[Direct match: $\frac{11!}{8}$]
  $\dfrac{11 \cdot 10 \cdot 9 \cdot 8 \cdot 7!}{8} = 11 \cdot 10 \cdot 9 \cdot 7!$
\end{example}

\begin{example}[Split into factors: $\frac{10!}{3! \cdot 2!} = \frac{10!}{12}$, options ending in $7!$]
  $12$ is not among $10, 9, 8$, but $12 = 3 \cdot 4$, with
  $9 / 3 = 3$ and $8 / 4 = 2$:
  \[
    \frac{10 \cdot 9 \cdot 8 \cdot 7!}{3 \cdot 4}
    = 10 \cdot 3 \cdot 2 \cdot 7! = 60 \cdot 7!.
  \]
\end{example}

\begin{example}[Comparing options for $\frac{11!}{8}$ --- all in terms of $7!$]
  Target: $\frac{11 \cdot 10 \cdot 9 \cdot 8}{8} \cdot 7! = 990 \cdot 7!$.
  \quad $8! = 8 \cdot 7!$;\quad $3 \cdot 8! = 24 \cdot 7!$;\quad
  $3! \cdot 8! = 48 \cdot 7!$;\quad $11! = 7\,920 \cdot 7!$;\quad
  $11 \cdot 10 \cdot 9 \cdot 7! = 990 \cdot 7!$ \checkmark
\end{example}

\noindent \textbf{The same trick inside $C(n, k)$:} unroll $n!$ down
to the bigger factorial in the denominator and it cancels completely
--- what remains is the ``count down $k$ numbers, divide by $k!$''
shortcut:
\[
  C(12, 4) = \frac{12!}{4! \cdot 8!}
  = \frac{12 \cdot 11 \cdot 10 \cdot 9 \cdot 8!}{4! \cdot 8!}
  = \frac{12 \cdot 11 \cdot 10 \cdot 9}{24} = 495.
\]

\subsection{Order matters vs.\ order does not matter}
\index{permutation $P(n,k)$}\index{combination $C(n,k)$}\index{n choose k@$n$ choose $k$|see{combination}}\index{order matters / does not matter}Choosing $k$ things out of $n$ (each thing at most once):

\begin{rulepar}
  \begin{center}
    \renewcommand{\arraystretch}{1.6}
    \begin{tabular}{@{}p{3.6cm} l p{4.4cm}@{}}
      \toprule
                    & Formula
                    & Quick by hand                                      \\
      \midrule
      \textbf{Order matters}
      (permutation) & $P(n, k) = \dfrac{n!}{(n-k)!}$
                    & multiply $k$ numbers counting down from $n$        \\
      \textbf{Order does not matter}
      (combination) & $C(n, k) = \dbinom{n}{k} = \dfrac{n!}{k!\,(n-k)!}$
                    & the same, then divide by $k!$                      \\
      \bottomrule
    \end{tabular}
  \end{center}
\end{rulepar}

\noindent $C(n, k)$ is read ``$n$ choose $k$''. \textbf{Why divide by
  $k!$:} $P(n, k)$ counts every ordering of the same $k$ things
separately, and $k$ things can be ordered in $k!$ ways --- so each
unordered group was counted $k!$ times.

\begin{rulepar}
  \noindent \textbf{Shortcut:} $C(n, k) = C(n, n-k)$ --- choosing which
  $k$ to take is the same as choosing which $n - k$ to leave. So
  $C(6, 5) = C(6, 1) = 6$. Also $C(n, 0) = C(n, n) = 1$ and
  $C(n, 1) = n$.
\end{rulepar}

\begin{example}[Bag with 7 red, 6 blue, 8 black; draw 3 red, 5 blue and 2 black, order does not matter]
  Three independent choices, one per colour --- ``and'', so
  multiply. Order does not matter, so each choice is a combination:
  \begin{align*}
    \text{red: } C(7, 3)   & = \frac{7 \cdot 6 \cdot 5}{3 \cdot 2 \cdot 1}
    = \frac{210}{6} = 35                                                       \\
    \text{blue: } C(6, 5)  & = C(6, 1) = 6                                     \\
    \text{black: } C(8, 2) & = \frac{8 \cdot 7}{2 \cdot 1} = \frac{56}{2} = 28
  \end{align*}
  \[
    35 \cdot 6 \cdot 28 = 210 \cdot 28 = 5\,880 \text{ ways.}
  \]
\end{example}

\subsection{Pascal's triangle (look-up table)}
\index{Pascal's triangle}\index{binomial coefficient}\index{combination $C(n,k)$!table}%
Row $n$, column $k$ is $C(n,k)$. Only $k\le6$ is shown: for larger $k$ use
the symmetry $C(n,k)=C(n,n-k)$, so $C(12,8)=C(12,4)=495$.
\begin{center}
  \small\renewcommand{\arraystretch}{1.15}
  \begin{tabular}{@{}r|rrrrrrr@{}}
    $n\backslash k$ & 0 & 1 & 2 & 3 & 4 & 5 & 6 \\\hline
      0 & 1 &  &  &  &  &  &  \\
      1 & 1 & 1 &  &  &  &  &  \\
      2 & 1 & 2 & 1 &  &  &  &  \\
      3 & 1 & 3 & 3 & 1 &  &  &  \\
      4 & 1 & 4 & 6 & 4 & 1 &  &  \\
      5 & 1 & 5 & 10 & 10 & 5 & 1 &  \\
      6 & 1 & 6 & 15 & 20 & 15 & 6 & 1 \\
      7 & 1 & 7 & 21 & 35 & 35 & 21 & 7 \\
      8 & 1 & 8 & 28 & 56 & 70 & 56 & 28 \\
      9 & 1 & 9 & 36 & 84 & 126 & 126 & 84 \\
      10 & 1 & 10 & 45 & 120 & 210 & 252 & 210 \\
      11 & 1 & 11 & 55 & 165 & 330 & 462 & 462 \\
      12 & 1 & 12 & 66 & 220 & 495 & 792 & 924 \\
  \end{tabular}
\end{center}
\begin{rulepar}
  \textbf{Pascal's rule:} $C(n,k)=C(n-1,k-1)+C(n-1,k)$ --- each entry is
  the sum of the two above it (e.g.\ $C(7,3)=15+20=35$); this is how a row
  is built from the one before. \textbf{Row sums:} $\sum_kC(n,k)=2^n$
  (every subset of an $n$-set has some size $k$, so this is the $2^n$
  subsets count). \textbf{Symmetry:} each row reads the same backwards.
  \textbf{Beyond the table:} cancel factorials, $C(n,k)=\frac{n(n-1)\cdots(n-k+1)}{k!}$.
\end{rulepar}

\subsection{Arranging letters with repeats}
\index{permutation $P(n,k)$!with repeated letters}\index{distinguishable permutations}%
\begin{rulepar}
  A word of $n$ letters, where the repeated letters occur
  $n_1, n_2, \ldots, n_k$ times, has
  \[
    \frac{n!}{n_1! \, n_2! \cdots n_k!}
  \]
  distinguishable arrangements (Kolman 3.1). \textbf{Why:} $n!$ treats
  the two M's as different (M$_1$, M$_2$); swapping them gives the same
  visible word, so each word was counted $2!$ times --- divide it away,
  for each repeated letter. Letters that occur once give $1! = 1$ and
  can be ignored.
\end{rulepar}

\begin{example}[MATHEMATICS]
  $11$ letters; M, A and T occur twice each (H, E, I, C, S once):
  \[
    \frac{11!}{2! \cdot 2! \cdot 2!} = \frac{11!}{8}
    = \frac{11 \cdot 10 \cdot 9 \cdot 8 \cdot 7!}{8}
    = 11 \cdot 10 \cdot 9 \cdot 7! = 4\,989\,600.
  \]
\end{example}

\noindent \textbf{Matching expanded-factorial options:} see
``Working with factorials without a calculator'' (3.1.3).

\subsection{Repetition allowed: stars and bars}
\index{stars and bars}\index{combination $C(n,k)$!with repetition}%
\begin{rulepar}
  Choosing $k$ things from $n$ \emph{kinds}, order does not matter, and
  the same kind may be chosen again:
  \[
    C(n + k - 1,\ k).
  \]
\end{rulepar}
\textbf{Why:} draw the $k$ chosen things as stars $\star$ and separate
the $n$ kinds with $n - 1$ bars $|$. Every choice is one row of $k$
stars and $n - 1$ bars --- $n + k - 1$ symbols --- and a row is fixed
by which $k$ positions hold the stars.

\noindent \textbf{Typical wording:} \emph{identical} objects
distributed over kinds/colours/boxes, or ``how many of each''.
``Identical'' is the key word: if the objects were different
(chair 1, chair 2, \ldots), each would pick a kind independently,
giving $n^k$ instead.

\begin{example}[8 identical chairs, 5 colours; how many ways to choose the colours?]
  Only how many chairs get each colour matters: $k = 8$ chairs,
  $n = 5$ colours, repetition allowed, order irrelevant. One choice,
  e.g.\ 3 red, 0 green, 2 blue, 3 black, 0 orange:
  \[
    \bigstar\bigstar\bigstar \;\big|\; \;\big|\; \bigstar\bigstar \;\big|\; \bigstar\bigstar\bigstar \;\big|\;
  \]
  8 stars and 4 bars, 12 positions:
  \[
    C(5 + 8 - 1,\ 8) = C(12, 8) = C(12, 4)
    = \frac{12 \cdot 11 \cdot 10 \cdot 9}{4 \cdot 3 \cdot 2 \cdot 1}
    = \frac{11\,880}{24} = 495.
  \]
  Wrong options: $5^8$ (chairs treated as different), ${}_8C_5$ (no
  repetition, numbers swapped).
\end{example}

\noindent \textbf{Notation:} Moodle and some books write ${}_{n}C_{k}$
and ${}_{n}P_{k}$\index{nCk@${}_nC_k$, ${}_nP_k$ notation} for
$C(n, k)$ and $P(n, k)$ --- the same thing.

\subsection{Exactly, at least, at most}
\index{exactly / at least / at most}\index{at least one (complement trick)}%
When the question puts a condition on a \emph{special} group (aces,
red balls, women on a committee), split everything into the special
group and the rest, then count each part:

\begin{rulepar}
  \begin{center}
    \renewcommand{\arraystretch}{1.4}
    \begin{tabular}{@{}ll@{}}
      \toprule
      Wording                       & Method                                                \\
      \midrule
      \textbf{exactly} $j$ special  & $C(\text{special}, j) \cdot C(\text{rest}, k - j)$    \\
      \textbf{at least one} special & all $-$ none $= C(\text{all}, k) - C(\text{rest}, k)$ \\
      \textbf{at most} $j$ special  & exactly $0$ $+$ exactly $1$ $+ \cdots +$ exactly $j$  \\
      \bottomrule
    \end{tabular}
  \end{center}
\end{rulepar}

\noindent Taking the other cards from ``the rest'' is what makes it
\emph{exactly}: no further special one can slip in.

\begin{example}[5-card hand from 52 cards with exactly one ace]
  Special: the $4$ aces; rest: $52 - 4 = 48$ non-aces. Choose $1$ ace
  \emph{and} $4$ non-aces --- multiply:
  \[
    C(4, 1) \cdot C(48, 4) = 4 \cdot {}_{48}C_4
    = 4 \cdot 194\,580 = 778\,320.
  \]
  Traps: $48 \cdot {}_{48}C_4$ (only $4$ aces to choose from);
  $52 \cdot {}_{52}C_4$ (lets more aces in, and counts hands several
  times). For comparison, \emph{at least} one ace would be
  $C(52, 5) - C(48, 5)$.
\end{example}

\subsection{When the MC options are products of primes}
\index{MC options as products of primes}\index{exponents!multiplying powers}Options like $7^2 \cdot 5 \cdot 3 \cdot 2^3$ are the answer in prime
factorized form. Do \emph{not} multiply out --- factor each piece into
primes and collect equal primes. Multiplying powers of the same base
\textbf{adds the exponents}: $p^a \cdot p^b = p^{a+b}$.

\begin{example}[the bag example above]
  \begin{align*}
    C(7, 3) & = 35 = 5 \cdot 7   \\
    C(6, 5) & = 6 = 2 \cdot 3    \\
    C(8, 2) & = 28 = 2^2 \cdot 7
  \end{align*}
  \[
    (5 \cdot 7)(2 \cdot 3)(2^2 \cdot 7)
    = 7^{1+1} \cdot 5 \cdot 3 \cdot 2^{1+2}
    = 7^2 \cdot 5 \cdot 3 \cdot 2^3.
  \]
  Check: $49 \cdot 5 \cdot 3 \cdot 8 = 5\,880$ \checkmark.
  Compare the options prime by prime (exponent of $7$, of $5$, of
  $3$, of $2$); exactly one matches.
\end{example}

\noindent This form is easier without a calculator: the numbers stay
small. Tip: cancel in prime form from the start, e.g.\
$C(8, 2) = \frac{8 \cdot 7}{2} = \frac{2^3 \cdot 7}{2} = 2^2 \cdot 7$.

\subsection{Which formula? The four cases}
\index{counting!the four cases}%
Choosing $k$ from $n$ always lands in one of four boxes. Ask:
\emph{does order matter?} and \emph{can the same thing be picked
  again?}

\begin{rulepar}
  \begin{center}
    \renewcommand{\arraystretch}{1.5}
    \begin{tabular}{@{}lcc@{}}
      \toprule
                                     & \textbf{no} repetition & repetition \textbf{allowed} \\
      \midrule
      \textbf{order matters}         & $P(n, k)$              & $n^k$                       \\
      \textbf{order does not matter} & $C(n, k)$              & $C(n + k - 1,\ k)$          \\
      \bottomrule
    \end{tabular}
  \end{center}
\end{rulepar}

\noindent \textbf{Recognising the problem type:}
\begin{itemize}[nosep]
  \item ``order does not matter'', ``choose'', ``a group/committee/hand
        of'' $\to$ $C(n, k)$
  \item ``order matters'', ``arrange'', ``first/second/third place'',
        ``line up'' $\to$ $P(n, k)$
  \item \emph{different} objects each getting one of $n$ options
        (each chair picks a colour, each digit of a PIN) $\to$ $n^k$
  \item \emph{identical} objects split over $n$ kinds, ``how many of
        each'' $\to$ $C(n + k - 1, k)$
  \item several independent choices made together (``and'')
        $\to$ multiply the counts
\end{itemize}

% =====================================================================

\subsection{Worked slide examples}
\index{multiplication principle!password}\index{subsets!number of ($2^n$)}\index{distinguishable permutations!MISSISSIPPI}%
\begin{rulepar}
  \textbf{3.1 Theorem 2.} Tasks $T_1,\dots,T_k$ in sequence, $T_i$ done in
  $n_i$ ways: $n_1n_2\cdots n_k$ ways in all.
  \textbf{Theorem 3.} Sequences of length $r$ from $n$ objects, repeats
  allowed: $n^r$. Special case: each element of an $n$-set is in or out, so
  an $n$-set has $2\cdot2\cdots2=\mathbf{2^n}$ subsets.
\end{rulepar}
\begin{itemize}[nosep]
  \item \textbf{Password}, 8 characters, first a letter, the other seven
        letter or digit, repeats allowed:
        $26\cdot36^7=2\,037\,468\,266\,496$.
  \item \textbf{Three-digit codes}, repeats allowed: $10^3=1000$.
  \item \textbf{GILBREATH}, 4-letter arrangements, no repeats (9 distinct
        letters): ${}_9P_4=9\cdot8\cdot7\cdot6=3024$.
  \item \textbf{Prizes}: 10 students, 1st/2nd/3rd prize:
        ${}_{10}P_3=10\cdot9\cdot8=720$. \textbf{Arranging all} $n$:
        ${}_nP_n=n!$ (5 books: $120$).
  \item \textbf{Distinguishable permutations} (Theorem 5,
        $n!/(k_1!\cdots k_t!)$): BANANA $=6!/(3!\,2!)=60$;
        MISSISSIPPI (M$1$, I$4$, S$4$, P$2$) $=11!/(1!\,4!\,4!\,2!)=34\,650$.
  \item \textbf{Hands}: $\binom{52}{5}=2\,598\,960$. \textbf{Groups}:
        $\binom{12}{4}=495$. Symmetry $\binom nr=\binom n{n-r}$, so
        $\binom{10}3=\binom{10}7=120$ (choosing who is in = choosing who is
        left out).
  \item \textbf{Committee}, 3 women from 20 \emph{and} 4 men from 30
        (multiply): $\binom{20}3\binom{30}4=1140\cdot27\,405=31\,241\,700$.
  \item \textbf{Ice cream}, 3 scoops from 5 flavours, repeats allowed,
        order irrelevant: $\binom{5+3-1}{3}=\binom73=35$.
\end{itemize}
\erratum{The committee slide prints $31{,}242{,}600$; the product is
  $1140\cdot27\,405=31\,241\,700$.}

\kolman{3}
\section{Pigeonhole principle}\src{Kolman 3.3}
\index{pigeonhole principle}\index{generalized pigeonhole principle}\index{ceiling function}%
\begin{rulebox}[title={3.3 Theorem 1: the pigeonhole principle}]
  If $n+1$ or more objects (pigeons) are put into $n$ boxes
  (pigeonholes), some box holds at least $2$ objects.\\
  \textbf{Generalized:} $N$ objects in $k$ boxes $\Rightarrow$ some box
  holds at least $\lceil N/k\rceil$.
  To force $m$ objects into one box you need $(m-1)k+1$ objects.
\end{rulebox}
\textbf{Recognising it:} (1) what are the objects, (2) what are the
categories (the boxes), (3) are there more objects than boxes? If so, a
crowded box is guaranteed. It is a \emph{worst-case} argument: the answer
to ``how many must I pick to be sure'' is ``everything that can go wrong,
plus one''.

\begin{itemize}[nosep]
  \item \textbf{Socks}, black and white: $2$ boxes, so $3$ socks
        guarantee a matching pair.
  \item \textbf{Birthdays}: $366$ students, $365$ days $\Rightarrow$ two
        share a birthday.
  \item \textbf{Tutorial groups}: $30$ students in $7$ groups $\Rightarrow$
        some group has $\ge\lceil30/7\rceil=5$.
  \item \textbf{Bicycles}: $50$ bicycles in $7$ colours $\Rightarrow$
        some colour has $\ge\lceil50/7\rceil=8$.
  \item \textbf{Three in one month}: worst case $2$ per month
        $=2\cdot12=24$ people without three, so $\mathbf{25}$ people
        guarantee it.
\end{itemize}

\kolman{4}
\section{Probability}\src{Kolman 3.4}
\index{probability}\index{sample space}\index{event}\index{mutually exclusive events}\index{equally likely outcomes}\index{expected value}\index{complement!of an event}%

\textbf{Terms.} Random experiment (e.g.\ rolling a die); outcome (a 4);
\textbf{sample space} $S$ = all outcomes ($\{1,\dots,6\}$);
\textbf{event} $E\subseteq S$ (even: $\{2,4,6\}$); \textbf{disjoint}
events share no outcomes.

\begin{rulepar}
  Each outcome $i$ has a probability $p_i$ with $0\le p_i\le1$ and $\sum p_i=1$ (something always happens), and $P(E)$ is the sum of the $p_i$ over the outcomes in $E$. With equally likely outcomes this becomes $|E|/|S|$ (``favourable over possible''); with an \emph{unfair} die add the given probabilities instead.
\end{rulepar}

\begin{rulebox}[title=Probability rules]
  \index{probability rules}%
  $0\le P(E)\le1$, \quad $P(S)=1$, \quad $P(\emptyset)=0$;\quad
  mutually exclusive $E_i$: $P(E_1\cup\dots\cup E_k)=\sum P(E_i)$.\\
  Equally likely outcomes: $\displaystyle P(E)=\frac{|E|}{|S|}$.\qquad
  Complement: $P(E)=1-P(\overline{E})$.\\
  Overlap: $P(E\cup F)=P(E)+P(F)-P(E\cap F)$.\qquad
  Expected value: $E(X)=\sum_i x_i\,P(X=x_i)$.
\end{rulebox}

\begin{example}[3.4 Ex.~8, unequal weights]
  $P(\{1\}),\dots,P(\{6\})=\frac1{12},\frac1{12},\frac13,\frac16,\frac14,\frac1{12}$.
  $P(\text{even})=\frac1{12}+\frac16+\frac1{12}=\frac13$.
\end{example}
\begin{example}[two dice, sum $\le3$]
  $|S|=36$, $E=\{(1,1),(1,2),(2,1)\}$, $P=\frac3{36}=\frac1{12}$.
\end{example}
\begin{example}[3.4 Ex.~9, four kings in a 4-card hand]
  $\binom{52}{4}=270725$ hands, one of them all kings:
  $P=\frac1{270725}$.
\end{example}
\begin{example}[3.4 Ex.~10, 6 red + 4 green, draw 4]
  Exactly 2 red: $\frac{\binom62\binom42}{\binom{10}4}=\frac{90}{210}=\frac37$.
  At most 2 red: $\frac{\binom60\binom44+\binom61\binom43+\binom62\binom42}{210}
    =\frac{1+24+90}{210}=\frac{23}{42}$.
  At most 3 red: $1-\frac{\binom64\binom40}{210}=1-\frac{15}{210}=\frac{13}{14}$.
\end{example}
\begin{example}[3.4 Ex.~11, die thrown 3 times: all equal OR no 4]
  $|F|=6$ (all equal), $|G|=5^3=125$ (no 4), $|F\cap G|=5$.
  $|F\cup G|=6+125-5=126$, $P=\frac{126}{216}=\frac7{12}$.
\end{example}
\begin{example}[sum is 8 OR doubles, two dice]
  $|E|=5$, $|F|=6$, $E\cap F=\{(4,4)\}$: $P=\frac{5+6-1}{36}=\frac5{18}$.
\end{example}
\begin{example}[fair die]
  $E(X)=\frac{1+2+3+4+5+6}{6}=3.5$.
\end{example}
\subsection{Adding fractions (exam method)}
\index{fractions, adding}\index{common denominator}Only fractions with the same denominator can be added: rewrite each
over a \textbf{common denominator} (the $\operatorname{lcm}$ of the
denominators), multiplying top and bottom by the same number, then add
the numerators. Finally simplify by dividing top and bottom by their
$\gcd$.

\begin{example}[Unfair die: $p_1 = \frac16$, $p_2 = \frac38$, $p_3 = \frac1{12}$,
    $p_4 = \frac1{24}$, $p_5 = \frac1{12}$, $p_6 = \frac14$; probability of an even number]
  Event $E = \{2, 4, 6\}$, so
  $P(E) = p_2 + p_4 + p_6 = \frac38 + \frac1{24} + \frac14$.
  Common denominator $\operatorname{lcm}(8, 24, 4) = 24$:
  \[
    \frac{3 \cdot 3}{8 \cdot 3} + \frac{1}{24} + \frac{1 \cdot 6}{4 \cdot 6}
    = \frac{9}{24} + \frac{1}{24} + \frac{6}{24}
    = \frac{16}{24}
    = \frac{16 \div 8}{24 \div 8}
    = \frac{2}{3}.
  \]
  \textbf{Sanity check} --- all six must sum to $1$:
  $\frac{4 + 9 + 2 + 1 + 2 + 6}{24} = \frac{24}{24} = 1$ \checkmark.
  So also $P(\text{odd}) = 1 - \frac23 = \frac13$.
\end{example}


% =====================================================================
\kolman{5}
\section{Recurrence relations}\src{Kolman 3.5}
\index{recurrence relation}\index{backtracking}\index{initial conditions}\index{linear homogeneous recurrence}%

A recurrence defines $a_n$ from earlier terms; with \textbf{initial
  conditions} it fixes the sequence.

\textbf{Backtracking} (repeated substitution --- unrolling the recursion):
\begin{itemize}
  \item $a_n=a_{n-1}+3$, $a_1=4$: $a_n=a_{n-2}+2\cdot3=\dots=a_1+3(n-1)=4+3(n-1)$.
  \item $b_n=2b_{n-1}+1$, $b_1=7$: $b_n=2^{n-1}b_1+2^{n-2}+\dots+2+1
          =7\cdot2^{n-1}+2^{n-1}-1=4\cdot2^n-1$.
\end{itemize}

\textbf{Linear homogeneous} of degree $k$:
$a_n=r_1a_{n-1}+r_2a_{n-2}+\dots+r_ka_{n-k}$, constants $r_i$. The degree
is how far back it reaches: $c_n=-2c_{n-1}$ (1),
$a_n=3a_{n-1}-2a_{n-2}$ (2), $b_n=2b_{n-1}+b_{n-3}$ (3).

\begin{rulebox}[title=Theorem 3.5.1 (degree 2)]
  \index{characteristic equation}\index{recurrence relation!degree 2 solution}\index{repeated root}%
  $a_n=r_1a_{n-1}+r_2a_{n-2}$ has characteristic equation
  $x^2-r_1x-r_2=0$.
  \begin{enumerate}[label=(\alph*)]
    \item Two distinct roots $s_1\neq s_2$: \quad $a_n=u\,s_1^{\,n}+v\,s_2^{\,n}$.
    \item One (repeated) root $s$: \quad $a_n=u\,s^n+v\,n\,s^n$.
          \hfill\textbf{Not on the slides} (Kolman 3.5 Thm 1b).
  \end{enumerate}
  Then solve for $u,v$ from the initial conditions, e.g.\
  $a_1=A$, $a_2=B$ $\Rightarrow$ $us_1+vs_2=A$, $us_1^2+vs_2^2=B$.\\
  \emph{Roots give the form; initial conditions pick the specific sequence.}
\end{rulebox}

\textbf{Why it works.} If $s^2=r_1s+r_2$, multiply by $s^{n-2}$:
$s^n=r_1s^{n-1}+r_2s^{n-2}$, so $a_n=s^n$ satisfies the recurrence.
Because the recurrence is linear, any combination $us_1^n+vs_2^n$ does too.

\notslides{The roots can be complex (negative discriminant); the same
  formula (a) still applies with complex $s_1,s_2$ and $u,v$. Quadratic
  formula: $x=\dfrac{r_1\pm\sqrt{r_1^2+4r_2}}{2}$.}

\begin{example}[3.5 Ex.~7: $c_n=3c_{n-1}-2c_{n-2}$, $c_1=5$, $c_2=3$]
  $x^2-3x+2=(x-1)(x-2)=0$, so $c_n=u+v2^n$. Then $u+2v=5$, $u+4v=3$
  $\Rightarrow v=-1$, $u=7$: \ $c_n=7-2^n$.
\end{example}
\begin{example}[3.5 Ex.~9: Fibonacci $f_n=f_{n-1}+f_{n-2}$, $f_1=f_2=1$]
  $x^2-x-1=0$, $s_{1,2}=\frac{1\pm\sqrt5}{2}$, $u=\frac1{\sqrt5}$,
  $v=-\frac1{\sqrt5}$:\quad
  $f_n=\frac1{\sqrt5}\left(\frac{1+\sqrt5}{2}\right)^n-\frac1{\sqrt5}\left(\frac{1-\sqrt5}{2}\right)^n$.
\end{example}
\begin{example}[$a_n=5a_{n-1}-6a_{n-2}$, $a_1=5$, $a_2=13$]
  $(x-2)(x-3)=0$; $2u+3v=5$, $4u+9v=13$ $\Rightarrow u=v=1$:\ $a_n=2^n+3^n$.
\end{example}
\textbf{MC shortcut:} plug $n=1,2$ into each candidate formula and compare
with the given initial terms --- quicker than re-deriving, and catches
$n$ vs.\ $n-1$ offset distractors.
\subsection{Exponent rule}
\index{exponent rule}\index{exponents!dividing powers}\begin{rulepar}
  Dividing powers of the same base: \textbf{subtract the exponents}.
  \[
    \frac{r^{n}}{r^{n-2}} = r^{\,n-(n-2)} = r^{2},
    \qquad
    \frac{r^{n-1}}{r^{n-2}} = r^{\,(n-1)-(n-2)} = r^{1} = r,
    \qquad
    \frac{r^{n-2}}{r^{n-2}} = r^{0} = 1.
  \]
\end{rulepar}

\noindent \textbf{Why this matters:} guess a solution of the form
$a_n = r^n$ and substitute it:
$r^{n} = c_1 r^{n-1} + c_2 r^{n-2}$. Dividing every term by $r^{n-2}$
(the rule above) gives $r^{2} = c_1 r + c_2$ --- the
\emph{characteristic equation}\index{characteristic equation} of step 1.

\subsection{The four-step recipe for distinct roots (exam method)}
\index{recurrence relation!four-step recipe}\index{general solution}\begin{rulepar}
  \begin{enumerate}[label=\arabic*.]
    \item \textbf{Characteristic equation.} Replace $a_n$ by $r^2$,
          $a_{n-1}$ by $r$ and $a_{n-2}$ by $1$:
          $r^2 = c_1 r + c_2$, i.e.\ $r^2 - c_1 r - c_2 = 0$.
    \item \textbf{Solve the quadratic} to get the roots $r_1$ and $r_2$.
    \item \textbf{General solution:}
          $a_n = A\,r_1^{\,n} + B\,r_2^{\,n}$ (when $r_1 \neq r_2$).
    \item \textbf{Fit the starting values.} Put $n = 1$ and $n = 2$
          into step 3 to get two equations in $A$ and $B$, and solve
          them.
  \end{enumerate}
\end{rulepar}

\begin{example}[$a_n = 7a_{n-1} - 12a_{n-2}$, $a_1 = 1$, $a_2 = 7$]
  \textbf{Step 1.} $r^2 = 7r - 12$, so $r^2 - 7r + 12 = 0$.

  \textbf{Step 2.} Two numbers that multiply to $12$ and add to $-7$:
  $-3$ and $-4$. So $(r - 3)(r - 4) = 0$, giving $r_1 = 3$,
  $r_2 = 4$.

  \textbf{Step 3.} $a_n = A \cdot 3^n + B \cdot 4^n$.

  \textbf{Step 4.}
  \begin{align*}
    n = 1: \quad & 3A + 4B = 1  \\
    n = 2: \quad & 9A + 16B = 7
  \end{align*}
  Multiply the first equation by $3$: $9A + 12B = 3$. Subtract it
  from the second: $4B = 4$, so $B = 1$. Then $3A + 4 = 1$, so
  $3A = -3$ and $A = -1$.
  \[
    a_n = 4^n - 3^n
  \]
  Check: $a_1 = 4 - 3 = 1$ \checkmark, $a_2 = 16 - 9 = 7$
  \checkmark, $a_3 = 7 \cdot 7 - 12 \cdot 1 = 37 = 64 - 27$
  \checkmark
\end{example}

\begin{example}[Repeated root: $a_n = 6a_{n-1} - 9a_{n-2}$, $a_1 = -3$, $a_2 = 9$]
  \index{repeated root}\index{recurrence relation!repeated root example}%
  \textbf{Steps 1--2.} $r^2 - 6r + 9 = (r-3)^2 = 0$: one root, $r = 3$
  (a perfect square --- spot it when the discriminant $c_1^2 + 4c_2$ is
  $0$).

  \textbf{Step 3 changes.} With one root the form is
  $a_n = (A + Bn)\,3^n$ --- the extra $n$ replaces the missing second root
  (Theorem 3.5.1b).

  \textbf{Step 4.}
  \begin{align*}
    n = 1: \quad & 3(A + B) = -3 \Rightarrow A + B = -1 \\
    n = 2: \quad & 9(A + 2B) = 9 \Rightarrow A + 2B = 1
  \end{align*}
  Subtract: $B = 2$, then $A = -3$. So
  $a_n = (-3 + 2n)\,3^n = -3\cdot3^n + 2n\cdot3^n$.

  Check: $a_3 = 6\cdot9 - 9\cdot(-3) = 81$ and
  $(-3+6)\cdot27 = 81$ \checkmark. MC trap: the options with
  $(-3)^n$ mix up the root $3$ with the coefficient $-3$.
\end{example}

\subsection{Linear and homogeneous --- recognising the type}
\index{linear homogeneous recurrence!recognising}
\begin{itemize}[nosep]
  \item \textbf{Linear}\index{linear (recurrence)}: each $a$ term appears on
        its own, only multiplied by a constant --- no $a_{n-1}^2$, no
        $a_{n-1} \cdot a_{n-2}$.
  \item \textbf{Homogeneous}\index{homogeneous (recurrence)}: every term on the
        right contains an $a$ --- no lone constant or $n$, as in
        $a_n = 3a_{n-1} - a_{n-2} + 1$.
  \item \textbf{Constant coefficients}\index{constant coefficients}: the
        numbers in front of the $a$'s are fixed numbers --- \emph{no $n$
          in them}. $a_n = (3 + n)\,a_{n-1} - a_{n-2}$,
        $a_n = 3a_{n-1} - n\,a_{n-2}$ and $a_n = 3n!\,a_{n-1}$ all fail
        this: the swap to $x$ would leave an $n$ in the equation, so
        there is no characteristic equation.
  \item \textbf{Degree}\index{degree (recurrence)}: how far back it reaches;
        $a_{n-2}$ as the furthest back means degree 2 (a quadratic
        characteristic equation), $a_{n-3}$ means degree 3.
\end{itemize}

\subsection{Working backwards: from characteristic equation to recurrence}
\index{characteristic equation!working backwards}
Undo step 1: get $x^2$ (or $r^2$) alone on the left --- moving terms
across $=$ \textbf{flips their sign} --- then swap $x^2 \to a_n$,
$x \to a_{n-1}$ and the plain number $\to a_{n-2}$. The letter
($x$ or $r$) does not matter.

\begin{example}[Which recurrence has characteristic equation $x^2 - 3x + 1 = 0$?]
  \[
    x^2 - 3x + 1 = 0
    \quad\Longrightarrow\quad
    x^2 = 3x - 1
    \quad\Longrightarrow\quad
    a_n = 3a_{n-1} - a_{n-2}.
  \]
  Check forwards: $a_n \to x^2$, $a_{n-1} \to x$, $a_{n-2} \to 1$
  gives $x^2 = 3x - 1$, i.e.\ $x^2 - 3x + 1 = 0$ \checkmark
\end{example}

\noindent \textbf{Method for ``which recurrence has this characteristic
  equation?'':} (1) rule out the wrong types first --- $n$ in a
coefficient, a lone constant, a square or product; (2) swap the
remaining options to $x$ and compare.

\noindent \textbf{Traps in the MC options:}
\begin{itemize}[nosep]
  \item \textbf{Starting values are irrelevant} --- the characteristic
        equation comes only from the $a_n = \ldots$ part.
  \item \textbf{Shifted indices are the same relation}\index{recurrence relation!shifted indices}:
        $a_n = 3a_{n-1} - a_{n-2}$, $a_{n+1} = 3a_n - a_{n-1}$ and
        $a_{n+2} = 3a_{n+1} - a_n$ all give $x^2 - 3x + 1 = 0$.
        Read them as highest index $\to x^2$, next $\to x$,
        lowest $\to 1$.
  \item \textbf{Sign not flipped}: $a_n = 3a_{n-1} + a_{n-2}$ gives
        $x^2 - 3x - 1 = 0$ --- wrong.
  \item \textbf{Coefficients swapped}: $a_n = a_{n-1} - 3a_{n-2}$ gives
        $x^2 = x - 3$, i.e.\ $x^2 - x + 3 = 0$ --- wrong.
  \item \textbf{$n$ in a coefficient}, e.g.\ $(3 + n)\,a_{n-1}$ --- looks
        like the right numbers, but the coefficients are not constant.
  \item A lone constant ($+1$) makes it not homogeneous; a square or
        product of $a$'s makes it not linear; $a_{n-3}$ makes it
        degree 3.
\end{itemize}

\noindent \textbf{Checking a formula by computer:}\index{recurrence.py (checker script)}
\texttt{discrete-math/code/recurrence.py} computes the terms straight
from the recurrence and compares them with a formula, e.g.
\begin{center}
  \texttt{python3 recurrence.py 7 -12 1 7 --formula "4**n - 3**n"}
\end{center}

\subsection{Sigma notation and the geometric series}
\index{sigma notation ($\Sigma$)}\index{sum|see{sigma notation}}$\Sigma$ (capital sigma) means ``add up''. The numbers below and above
it give where the counter starts and stops; the expression to the
right is what gets added for each value:
\[
  \sum_{i=1}^{4} 3^{i-1} = 3^{0} + 3^{1} + 3^{2} + 3^{3}
  = 1 + 3 + 9 + 27 = 40.
\]
It is a for-loop:
\texttt{for i in range(1, 5): total += 3**(i-1)}.
Remember $r^{0} = 1$, so the first term is $1$.

\paragraph{Geometric series.}
\index{geometric series}\begin{rulepar}
  A sum where each term is the previous one times a fixed ratio $r$ is
  a \textbf{geometric series}. For $r \neq 1$:
  \[
    1 + r + r^{2} + \cdots + r^{n-1}
    = \sum_{i=1}^{n} r^{i-1}
    = \frac{r^{n} - 1}{r - 1}.
  \]
  Note the count: $n$ terms, and the last exponent is $n - 1$.
\end{rulepar}

\paragraph{Why --- multiply by $r$ and subtract.} Multiplying shifts
every term one place up; subtracting cancels everything in the middle:
\begin{align*}
  S      & = 1 + r + r^{2} + \cdots + r^{n-1}                      \\
  rS     & = \phantom{1 + {}} r + r^{2} + \cdots + r^{n-1} + r^{n} \\
  rS - S & = r^{n} - 1
  \quad\Longrightarrow\quad
  S = \frac{r^{n} - 1}{r - 1}.
\end{align*}

\begin{example}[$\sum_{i=1}^{20} 3^{i-1}$]
  Here $r = 3$, $n = 20$ terms ($i = 1, \ldots, 20$):
  \[
    S = 1 + 3 + 9 + \cdots + 3^{19}
    = \frac{3^{20} - 1}{3 - 1}
    = \frac{3^{20} - 1}{2}
    = \frac{3\,486\,784\,401 - 1}{2}
    = 1\,743\,392\,200.
  \]
  Small check: $1 + 3 + 9 = 13$ and $\frac{3^{3} - 1}{2} =
    \frac{26}{2} = 13$ \checkmark\\
  On the MC exam (no calculator) expect the answer in the form
  $\frac{3^{20} - 1}{2}$.
\end{example}

\noindent \textbf{Link to recurrences:} the terms $1, 3, 9, 27, \ldots$
satisfy $a_n = 3\,a_{n-1}$ --- each term is the previous one times a
fixed number, which is where the $r^{n}$ in the recurrence recipe (3.5.2)
comes from.


% =====================================================================
\newpage
\part{Relations and Digraphs}

\kolman{1}
\section{Relations and digraphs}\src{Kolman 4.1--4.3; own notes}
\index{relation}\index{digraph}%
A \textbf{relation} $R$ on a set $A$ is a set of ordered pairs
$(x, y)$; $x \mathrel{R} y$ means $(x, y) \in R$, ``$x$ is related to
$y$''. In the \textbf{digraph} (directed graph) each pair is an arrow
$x \to y$. \textbf{First step for any digraph question: write out the
  list of pairs} --- then every option is a lookup. (Think of $R$ as a
table of \texttt{(from, to)} rows.)

\begin{center}
  \renewcommand{\arraystretch}{1.3}
  \begin{tabular}{@{}lll@{}}
    \toprule
    Notation                & Meaning                      & Read from the digraph                           \\
    \midrule
    $x \mathrel{R} y$       & $x$ is related to $y$        & arrow $x \to y$                                 \\
    $R(x)$                  & everything $x$ is related to & where $x$'s arrows point                        \\
    $\operatorname{Dom}(R)$ & domain                       & elements with an arrow \textbf{out}             \\
    $\operatorname{Ran}(R)$ & range                        & elements with an arrow \textbf{in}              \\
    $x \mathrel{R^n} y$     & path of length $n$           & walk $x \to \cdots \to y$ in exactly $n$ arrows \\
    \bottomrule
  \end{tabular}
\end{center}
\index{domain (relation)}\index{range (relation)}\index{R(x)@$R(x)$, relative set}\index{path!of length $n$ ($R^n$)}%
\noindent $\operatorname{Dom}(R)$ is \texttt{SELECT DISTINCT from},
$\operatorname{Ran}(R)$ is \texttt{SELECT DISTINCT to}. For $R^n$ one
path of exactly $n$ steps is enough; a shorter path does not count.

\begin{example}[Which statement is false for this digraph?]
  \leavevmode\par\noindent
  \begin{minipage}[c]{0.36\textwidth}
    \centering
    \begin{tikzpicture}[>=stealth, every node/.style={circle, draw, inner sep=1.5pt, minimum size=5mm}]
      \node (a) at (0, 2.6)   {$a$};
      \node (b) at (2.6, 2.6) {$b$};
      \node (c) at (1.3, 1.4) {$c$};
      \node (d) at (3.8, 1.4) {$d$};
      \node (e) at (2.6, 0)   {$e$};
      \draw[->] (a) to[bend left=10] (b);
      \draw[->] (a) -- (c);
      \draw[->] (a) to[bend right=40] (e);
      \draw[->] (b) -- (c);
      \draw[->] (b) to[bend left=25] (d);
      \draw[->] (c) -- (e);
      \draw[->] (e) to[bend right=8] (b);
    \end{tikzpicture}
  \end{minipage}\hfill
  \begin{minipage}[c]{0.6\textwidth}
    $R = \{(a,b), (a,c), (a,e), (b,c), (b,d), (c,e), (e,b)\}$.
    \begin{itemize}[nosep]
      \item $a \mathrel{R} b$: arrow $a \to b$ --- true.
      \item $R(a) = \{b, c, e\}$ --- true.
      \item $\operatorname{Dom}(R) = \{a, b, c, e\}$ ($d$ has no
            arrow out) --- true.
      \item $\operatorname{Ran}(R) = \{a, b\}$ --- \textbf{false}:
            the range is $\{b, c, d, e\}$ (nothing points to $a$).
      \item $b \mathrel{R^3} b$: $b \to c \to e \to b$ --- true.
      \item $a \mathrel{R^3} b$: $a \to c \to e \to b$ --- true.
    \end{itemize}
  \end{minipage}
\end{example}

\subsection{Computing \texorpdfstring{$R^2$}{R2} from a list of pairs (exam method)}
\index{path!computing $R^2$}\index{composition of relations}%
$(x, z) \in R^2$ when there is a middle element $y$ with
$(x, y) \in R$ \emph{and} $(y, z) \in R$ --- two steps, $x \to y \to z$.
(SQL: \texttt{SELECT r1.from, r2.to FROM R r1 JOIN R r2 ON r1.to = r2.from}.)
\begin{enumerate}[label=\arabic*.]
  \item \textbf{One-step table:} for each element, list where it can go
        in one step. Check the total equals the number of pairs.
  \item \textbf{Second step:} for each starting element, look up every
        landing place in the same table and collect the results.
\end{enumerate}
For $R^3$, take one more step from the $R^2$ results.

\begin{example}[$R = \{(1,2), (2,2), (3,4), (4,3), (2,1), (4,1), (3,2)\}$ on $\{1, 2, 3, 4\}$; find $R^2$]
  One-step table: $1 \to 2$;\quad $2 \to 1, 2$;\quad $3 \to 2, 4$;\quad
  $4 \to 1, 3$ \quad ($1 + 2 + 2 + 2 = 7$ pairs \checkmark).
  \begin{itemize}[nosep]
    \item from $1$: $\to 2 \to \{1, 2\}$
    \item from $2$: $\to 1 \to \{2\}$ and $\to 2 \to \{1, 2\}$, together $\{1, 2\}$
    \item from $3$: $\to 2 \to \{1, 2\}$ and $\to 4 \to \{1, 3\}$, together $\{1, 2, 3\}$
    \item from $4$: $\to 1 \to \{2\}$ and $\to 3 \to \{2, 4\}$, together $\{2, 4\}$
  \end{itemize}
  \[
    R^2 = \{(1,1), (1,2), (2,1), (2,2), (3,1), (3,2), (3,3), (4,2), (4,4)\}.
  \]
\end{example}

\noindent \textbf{Exam layout (pen and paper, about 2 minutes):} one line
per element, then replace each line by the merged lines of its
targets --- duplicates collapse. (Each line is a row of bits; the new
row is the OR of the rows it points to --- Kolman's Boolean matrix
product, written as lists.)
\begin{center}
  \begin{tabular}{@{}l@{\qquad}l@{}}
    \toprule
    $R$ (one step) & $R^2$ (merge the targets' lines)                           \\
    \midrule
    $1: 2$         & line 2 $= \{1, 2\}$                                        \\
    $2: 1, 2$      & line 1 $+$ line 2 $= \{2\} \cup \{1, 2\} = \{1, 2\}$       \\
    $3: 2, 4$      & line 2 $+$ line 4 $= \{1, 2\} \cup \{1, 3\} = \{1, 2, 3\}$ \\
    $4: 1, 3$      & line 1 $+$ line 3 $= \{2\} \cup \{2, 4\} = \{2, 4\}$       \\
    \midrule
    $7$ pairs      & $2 + 2 + 3 + 2 = 9$ pairs                                  \\
    \bottomrule
  \end{tabular}
\end{center}
\index{path!computing $R^2$!exam layout}%
\noindent \textbf{Step 2 as ``friends of friends'':} read each line as
\emph{who that element knows}. For line $3$: (1) who does $3$ know?
$2$ and $4$; (2) ask each of them who \emph{they} know --- line $2$
gives $\{1, 2\}$, line $4$ gives $\{1, 3\}$; (3) put the answers
together, each number once: $\{1, 2, 3\}$. Watch out:
\begin{itemize}[nosep]
  \item always look up in the \textbf{original} one-step table, never
        in the $R^2$ lines you are building;
  \item your own friends list is \emph{not} part of the answer --- only
        what comes out of step (2) ($3 \to 4$ directly does not give
        $(3,4) \in R^2$);
  \item a loop counts: $2 \in$ line $2$, so asking $2$ also returns $2$.
\end{itemize}

Then match line by line: find the options whose $3$-pairs are exactly
$(3,1), (3,2), (3,3)$, and so on --- usually one or two lines leave a
single option.

\noindent \textbf{Eliminating MC options fast:} (1) count the pairs ---
here $9$, which rules out options with another count; (2) test one or
two easy pairs, e.g.\ $(1,1)$ via $1 \to 2 \to 1$ and $(2,1)$ via
$2 \to 2 \to 1$. \textbf{Trap:} an option that is just $R$ itself.

\noindent \textbf{$R^2$ is built from scratch, not ``$R$ plus extra
  pairs''.} Think of $R$ as one-way roads and $R^2$ as ``reachable by
driving \emph{exactly} two roads''. A pair can be in $R$ but not in
$R^2$: above, $4 \to 1$ is one road, but every two-road trip from $4$
ends at $2$ or $4$, so $(4, 1) \notin R^2$. A loop like $(2, 2)$ counts
as a road too ($1 \to 2 \to 2$ gives $(1, 2) \in R^2$).

\subsection{\texorpdfstring{$R^\infty$}{R-infinity}: the connectivity relation (exam method)}
\index{connectivity relation $R^\infty$}\index{path!$R^\infty$ reachability}\index{transitive closure}%
\textbf{Picture $R$ as a map of one-way roads} between towns: each pair
$(a,b)\in R$ is a road from $a$ to $b$, and $(a,a)$ is a loop road back
to the same town. Then $R^2$ is ``towns reachable by driving
\emph{exactly} two roads'' (see above), and
\[
  R^\infty = R \cup R^2 \cup R^3 \cup \cdots
\]
is \textbf{``towns reachable by driving one or more roads, any number''}:
$(x,y)\in R^\infty$ means \emph{you can get from $x$ to $y$ at all}. It
is reachability in a directed graph (a graph search from each node), and
the smallest transitive relation containing $R$ (transitive closure).

\begin{rulepar}
  \textbf{Method (one town at a time):}
  \begin{enumerate}[label=\arabic*.,nosep]
    \item Write the one-step table: for each town, where its roads lead.
    \item Pick a starting town. Collect its neighbours; then the
          neighbours' neighbours; keep going and \emph{keep a town only
            once}, until nothing new turns up.
    \item Every town you collected gives a pair (start, town) in
          $R^\infty$. Repeat for each starting town.
  \end{enumerate}
  On $n$ towns, paths longer than $n$ never add anything, so the search
  always ends.
\end{rulepar}

\begin{example}[MC practice: $R = \{(1,2),(2,1),(2,3),(3,3)\}$ on $\{1,2,3\}$; find $R^\infty$]
  \textbf{Roads:} $1\to2$, \ $2\to1$, \ $2\to3$, \ $3\to3$ (a loop).
  One-step table: $1 \to 2$;\quad $2 \to 1, 3$;\quad $3 \to 3$.
  \begin{itemize}[nosep]
    \item \textbf{Start at $1$:} the only road goes to $2$. From $2$
          you can go to $1$ and $3$. So $2$, $1$ and $3$ are all
          reachable ($1$ again via $1\to2\to1$). Reaches $\{1,2,3\}$ :
          pairs $(1,1),(1,2),(1,3)$.
    \item \textbf{Start at $2$:} roads to $1$ and $3$; from $1$ you
          can drive back to $2$. Reaches $\{1,2,3\}$ :
          pairs $(2,1),(2,2),(2,3)$.
    \item \textbf{Start at $3$:} the only road is the loop back to
          $3$, so you can never leave. Reaches $\{3\}$ : pair $(3,3)$.
  \end{itemize}
  \[
    R^\infty = \{(1,1),(1,2),(1,3),(2,1),(2,2),(2,3),(3,3)\}\quad(7\text{ pairs}).
  \]
  \textbf{Two things that trip people up.}
  (1)~$(1,1)\in R^\infty$ although $(1,1)\notin R$: you can leave $1$ and
  come back ($1\to2\to1$), and a path is allowed to return to its start.
  (2)~$(3,1)\notin R^\infty$: $3$ is a dead end that only loops on
  itself, so nothing outside $\{3\}$ is reachable from it. This
  eliminates every option containing $(3,1)$ or $(3,2)$.
\end{example}

\subsection{Slide results for 4.1--4.3}
\index{product set}\index{Cartesian product}\index{partition}\index{block of a partition}\index{image of a set under a relation}\index{relation!matrix of a relation}%
\begin{itemize}[nosep]
  \item \textbf{Product set}
        $A\times B=\{(a,b)\mid a\in A,\,b\in B\}$; a relation from $A$ to
        $B$ is any subset $R\subseteq A\times B$ (on $A$ when $B=A$).
        $|A_1\times\dots\times A_m|=|A_1|\cdots|A_m|$ (the multiplication
        principle again).
  \item \textbf{Partition} of $A$: nonempty blocks $A_1,\dots,A_k$ with
        $A_1\cup\dots\cup A_k=A$ (nothing left out) and $A_i\cap A_j=\varnothing$
        for $i\ne j$ (nothing in two blocks).
  \item \textbf{Matrix of $R$} from $A=\{a_1..a_m\}$ to $B=\{b_1..b_n\}$:
        $m\times n$, rows $=A$, columns $=B$, $m_{ij}=1$ iff
        $(a_i,b_j)\in R$. Digraph (on one set): vertex per element, arrow
        $a_i\to a_j$ for $a_iRa_j$; \textbf{in-degree} = arrows entering,
        \textbf{out-degree} = arrows leaving.
\end{itemize}
\begin{rulebox}[title={4.2 Theorems 1 and 2}]
  For $R$ from $A$ to $B$ and $A_1,A_2\subseteq A$, with $R(A_i)$ the set of
  partners of the elements of $A_i$:
  (1) $A_1\subseteq A_2\Rightarrow R(A_1)\subseteq R(A_2)$;
  (2) $R(A_1\cup A_2)=R(A_1)\cup R(A_2)$;
  (3) $R(A_1\cap A_2)\subseteq R(A_1)\cap R(A_2)$ --- \textbf{only
    $\subseteq$}: with $R=\{(1,a),(2,a)\}$, $A_1=\{1\}$, $A_2=\{2\}$ the
  left side is $\varnothing$, the right side $\{a\}$.\\
  If $R(a)=S(a)$ for every $a\in A$, then $R=S$.
\end{rulebox}
\begin{rulepar}
  \textbf{4.3 Theorems 1 and 2.} $M_{R^2}=M_R\odot M_R$ and, for
  $n\ge2$, $M_{R^n}=M_R\odot\cdots\odot M_R$ ($n$ factors): $(a,b)\in R^n$
  iff there is a path of length $n$ from $a$ to $b$. Reachability:
  $R^\infty=R\cup R^2\cup R^3\cup\cdots$.
\end{rulepar}

\kolman{4}
\section{Properties of relations}\src{Kolman 4.4; own notes}
\index{relation!properties}\index{reflexive}\index{irreflexive}\index{symmetric}\index{antisymmetric}\index{asymmetric}\index{transitive}%
\textbf{Pen-and-paper tool:} draw $R$ as a $0/1$ grid (the matrix
$M_R$) --- rows ``from'', columns ``to'', a $1$ for each pair. Check
the count of $1$'s against the number of pairs. Each property is then
a visual check, and \textbf{one counterexample breaks any of them}.

\begin{rulepar}
  \begin{center}
    \renewcommand{\arraystretch}{1.3}
    \begin{tabular}{@{}lll@{}}
      \toprule
      Property      & Definition                                                          & Check in the grid                                      \\
      \midrule
      reflexive     & $x \mathrel{R} x$ for all $x$                                       & diagonal all $1$                                       \\
      irreflexive   & $x \mathrel{R} x$ for no $x$                                        & diagonal all $0$                                       \\
      symmetric     & $x \mathrel{R} y \Rightarrow y \mathrel{R} x$                       & mirror image across the diagonal                       \\
      antisymmetric & $x \mathrel{R} y$ and $y \mathrel{R} x \Rightarrow x = y$           & never $1$ in both mirror positions off the diagonal    \\
      asymmetric    & $x \mathrel{R} y \Rightarrow$ not $y \mathrel{R} x$                 & antisymmetric \emph{and} diagonal all $0$              \\
      transitive    & $x \mathrel{R} y$ and $y \mathrel{R} z \Rightarrow x \mathrel{R} z$ & every chain $x \to y \to z$ has the shortcut $x \to z$ \\
      \bottomrule
    \end{tabular}
  \end{center}
\end{rulepar}

\noindent \textbf{Traps:} a relation can be \emph{neither} reflexive nor
irreflexive (diagonal mixed). Antisymmetric ignores the diagonal;
asymmetric is the stricter version that also forbids it --- any
$(x, x)$ makes a relation not asymmetric.

\begin{rulepar}
  \textbf{Special relations on $A$} \emph{Empty} $R=\varnothing$;
  \emph{full} $R=A\times A$; \emph{identity} (diagonal)
  $\Delta=I_A=\{(a,a)\mid a\in A\}$; \emph{inverse}
  $R^{-1}=\{(b,a)\mid(a,b)\in R\}$. Asymmetric $=$ $aRb\Rightarrow\neg\,bRa$,
  which also forbids every loop $aRa$; antisymmetric only restricts
  \emph{distinct} pairs. In the matrix: antisymmetric means no $m_{ij}=m_{ji}=1$
  with $i\ne j$.
\end{rulepar}


\begin{example}[$R = \{(1,1), (1,4), (2,1), (2,3), (3,1), (3,3), (3,4), (4,2)\}$ on $\{1, 2, 3, 4\}$: which is true?]
  \leavevmode\par\noindent
  \begin{minipage}[c]{0.28\textwidth}
    \centering
    $\begin{array}{c|cccc}
          & 1 & 2 & 3 & 4 \\ \hline
        1 & 1 & 0 & 0 & 1 \\
        2 & 1 & 0 & 1 & 0 \\
        3 & 1 & 0 & 1 & 1 \\
        4 & 0 & 1 & 0 & 0
      \end{array}$
  \end{minipage}\hfill
  \begin{minipage}[c]{0.7\textwidth}
    \begin{itemize}[nosep]
      \item reflexive: diagonal $1, 0, 1, 0$ --- no, $(2,2)$ missing
      \item irreflexive: $(1,1)$ present --- no
      \item symmetric: $(1,4)$ but no $(4,1)$ --- no
      \item \textbf{antisymmetric}: the six mirror pairs
            $(1,2)/(2,1)$, $(1,3)/(3,1)$, \ldots, $(3,4)/(4,3)$ never
            both $1$ --- \textbf{yes}
      \item asymmetric: $(1,1)$ on the diagonal --- no
      \item transitive: $1 \to 4 \to 2$ but no $(1,2)$ --- no
    \end{itemize}
  \end{minipage}
\end{example}

\begin{example}[MC practice: $x \mathrel{R} y \iff x \ne y$ on $\Z^+$, how many properties hold?]
  \index{relation!given by a formula}\index{relation!properties}\index{counterexample}%
  Test each property with the smallest possible witnesses, $1$ and $2$:
  \begin{itemize}[nosep]
    \item symmetric: $x \ne y \Rightarrow y \ne x$ --- \textbf{yes}
    \item irreflexive: $x \ne x$ never holds, so no $(x,x)$ --- \textbf{yes}
    \item reflexive: $(1,1) \notin R$ --- no
    \item asymmetric: $(1,2)$ and $(2,1)$ are both in $R$ --- no
          (symmetric and nonempty always kills asymmetric)
    \item transitive: $1 \to 2 \to 1$ but $(1,1) \notin R$ --- no
  \end{itemize}
  Answer: $\mathbf{2}$ properties (symmetric and irreflexive).
\end{example}

\kolman{5}
\section{Equivalence relations}\src{Kolman 4.5}
\index{equivalence relation}\index{equivalence class}\index{reflexive}\index{symmetric}\index{transitive}\index{relation!matrix of a relation}%

The two exercises below (and in 4.7) are the ones worked at the end of
the lecture. Row $i$, column $j$ of $M_R$ is
$1$ iff $(i,j)\in R$.
\begin{itemize}
  \item \textbf{Reflexive}: all diagonal entries $1$.
        \textbf{Symmetric}: $M_R=M_R^T$.
        \textbf{Transitive}: $(a,b),(b,c)\in R\Rightarrow(a,c)\in R$.
  \item \textbf{Equivalence relation} = reflexive + symmetric + transitive;
        it splits the set into equivalence classes.
\end{itemize}

\begin{rulebox}[title={4.5 Theorems 1, 2 and Lemma 1}]
  \textbf{Thm 1.} If $P$ is a partition of $A$ and $aRb$ means ``$a$ and $b$
  lie in the same block'', then $R$ is an equivalence relation.
  (Reflexive: $a$ shares a block with itself. Symmetric: ``same block'' is
  symmetric. Transitive: $a,b$ and $b,c$ in the same block put $a,b,c$ in one
  block.)\\
  \textbf{Lemma 1.} For an equivalence relation, $aRb\iff R(a)=R(b)$, where
  the class of $a$ is $[a]=R(a)=\{x\mid aRx\}$.\\
  \textbf{Thm 2.} The distinct classes $\{[a]\mid a\in A\}$ form a
  \textbf{partition} of $A$. So equivalence relations and partitions are the
  same thing seen two ways.
\end{rulebox}
\begin{itemize}[nosep]
  \item Mnemonic: \textbf{equivalence} $=$ reflexive $+$ symmetric $+$
        transitive; \textbf{partial order} $=$ reflexive $+$ antisymmetric $+$
        transitive.
  \item $a\le b$ on $\Z$ is reflexive and transitive but \textbf{not
        symmetric} ($1\le2$, $2\not\le1$), so not an equivalence.
  \item \textbf{Congruence mod $n$}: $aRb\iff n\mid(a-b)$ (same remainder
        on division by $n$) is an equivalence relation on $\Z$. Mod $3$ it
        has three classes
        $[0]=\{\dots,-3,0,3,\dots\}$, $[1]=\{\dots,-2,1,4,\dots\}$,
        $[2]=\{\dots,-1,2,5,\dots\}$, and $\{[0],[1],[2]\}$ is a partition
        of $\Z$.
  \item A set such as $A_3\cup A_4\cup A_5$ with blocks
        $\{a,c,e,g\},\{b,d\},\{f,h\}$ is a partition of
        $\{a,\dots,h\}$ (union is everything, blocks disjoint); a family
        with overlapping blocks, like $\{a,b,c,d\}$ and $\{a,c,e,f,g\}$, is
        not.
\end{itemize}

\begin{example}[4.5.1]
  (1) $M_R=\begin{bmatrix}1&0&0\\0&1&1\\0&1&1\end{bmatrix}$ on $\{a,b,c\}$:
  reflexive, symmetric, transitive $\Rightarrow$ equivalence, classes
  $\{a\}$, $\{b,c\}$.
  (2) $M_R=\begin{bmatrix}1&0&1\\0&1&0\\0&0&1\end{bmatrix}$: $(a,c)\in R$
  but $(c,a)\notin R$, not symmetric $\Rightarrow$ not an equivalence
  (answer not on slides).
\end{example}
\kolman{7}
\section{Operations on relations}\src{Kolman 4.7}
\index{relation!complement}\index{relation!inverse}\index{relation!union and intersection}%
\begin{rulepar}
  Complement $\overline R$: flip every bit. $R\cap S$: AND. $R\cup S$: OR.
  Inverse $R^{-1}$: transpose, $M_{R^{-1}}=M_R^T$.
\end{rulepar}

\begin{example}[4.7.9]
  $M_R=\begin{bmatrix}1&1&0&1\\0&0&0&1\\1&1&1&0\end{bmatrix}$,
  $M_S=\begin{bmatrix}0&1&1&0\\1&0&0&1\\1&1&0&0\end{bmatrix}$:\\[3pt]
  $M_{\overline S}=\begin{bmatrix}1&0&0&1\\0&1&1&0\\0&0&1&1\end{bmatrix}$,\
  $M_{R\cap S}=\begin{bmatrix}0&1&0&0\\0&0&0&1\\1&1&0&0\end{bmatrix}$,\
  $M_{R\cup S}=\begin{bmatrix}1&1&1&1\\1&0&0&1\\1&1&1&0\end{bmatrix}$,\
  $M_{R^{-1}}=\begin{bmatrix}1&0&1\\1&0&1\\0&0&1\\1&1&0\end{bmatrix}$.
\end{example}

\begin{rulebox}[title={4.7 Theorems 1--4}]
  \textbf{Thm 1.} $R\subseteq S\Rightarrow R^{-1}\subseteq S^{-1}$ and
  $\overline S\subseteq\overline R$;
  $(R\cap S)^{-1}=R^{-1}\cap S^{-1}$; $(R\cup S)^{-1}=R^{-1}\cup S^{-1}$;
  $\overline{R\cap S}=\overline R\cup\overline S$ (De Morgan).\\
  \textbf{Thm 2.} $R$ reflexive $\Rightarrow$ $R^{-1}$ reflexive; $R,S$
  reflexive $\Rightarrow$ $R\cap S$ and $R\cup S$ reflexive;
  $R$ reflexive $\iff$ $\overline R$ irreflexive.\\
  \textbf{Thm 3.} $R$ symmetric $\iff R=R^{-1}$;
  antisymmetric $\iff R\cap R^{-1}\subseteq\Delta$;
  asymmetric $\iff R\cap R^{-1}=\varnothing$.\\
  \textbf{Thm 4.} $R$ symmetric $\Rightarrow$ $R^{-1}$ and $\overline R$
  symmetric; $R,S$ symmetric $\Rightarrow$ $R\cap S$, $R\cup S$ symmetric.
\end{rulebox}

\begin{example}[4.7 Ex.~4 and Ex.~6 (slides)]
  $M_R=\begin{bmatrix}1&1&0\\0&0&1\\0&0&1\end{bmatrix}$,
  $M_S=\begin{bmatrix}0&1&1\\1&0&0\\0&0&0\end{bmatrix}$:\
  $M_{\overline R}=\begin{bmatrix}0&0&1\\1&1&0\\1&1&0\end{bmatrix}$,\
  $M_{R^{-1}}=M_R^T=\begin{bmatrix}1&0&0\\1&0&0\\0&1&1\end{bmatrix}$,\
  $M_{R\cap S}=\begin{bmatrix}0&1&0\\0&0&0\\0&0&0\end{bmatrix}$,\
  $M_{R\cup S}=\begin{bmatrix}1&1&1\\1&0&1\\0&0&1\end{bmatrix}$.\\[3pt]
  On $\{1,2,3\}$ let $R=\{(1,1),(1,2),(2,1),(1,3),(3,1)\}$ and
  $S=\{(1,1),(1,2),(2,1),(2,2),(3,3)\}$ (both symmetric). Then
  $R^{-1}=R$, $\overline R=\{(2,2),(2,3),(3,2),(3,3)\}$,
  $R\cap S=\{(1,1),(1,2),(2,1)\}$,
  $R\cup S=\{(1,1),(1,2),(1,3),(2,1),(2,2),(3,1),(3,3)\}$ --- all symmetric,
  as Theorem 4 predicts.
\end{example}

\begin{rulepar}
  \textbf{Closure} of $R$ for a property (reflexive, symmetric, transitive):
  the \emph{smallest} relation containing $R$ with that property --- add
  only the pairs that are needed. \textbf{Composition}: for $R$ from $A$ to
  $B$ and $S$ from $B$ to $C$, $a\,(S\circ R)\,c\iff\exists b:\ aRb\wedge
  bSc$; note the matrix order, $M_{S\circ R}=M_R\odot M_S$ (R first). With
  the $M_R,M_S$ above, $M_{S\circ R}=\begin{bmatrix}1&1&1\\0&0&0\\0&0&0\end{bmatrix}$.
\end{rulepar}

% =====================================================================
\newpage
\part{Functions}

\kolman{1}
\section{Functions}\src{Kolman 5.1}
\index{function}\index{identity function}\index{composition of functions}\index{injective (one-to-one)}\index{surjective (onto)}\index{bijective}\index{everywhere defined function}%

\begin{defn}
  A \textbf{function} $f:A\to B$ is a relation where each $a\in A$ is
  related to \emph{exactly one} $b\in B$; write $f(a)=b$. As a set of
  pairs: $f=\{(a,f(a))\mid a\in A\}$. $f(x)=x^2$ on $\R$ is a function
  ($f(2)=f(-2)=4$ is fine --- two inputs may share an output, but one input
  may not have two outputs).
\end{defn}

\textbf{Identity}: $1_A(a)=a$; e.g.\ $1_{\{1,2,3\}}=\{(1,1),(2,2),(3,3)\}$.
\textbf{Composition}: $f:A\to B$, $g:B\to C$, $(g\circ f)(a)=g(f(a))$ ---
apply $f$ first.

\begin{keybox}[title={Special types, for $f:A\to B$}]
  \textbf{Injective} (one-to-one): $f(a)=f(a')\Rightarrow a=a'$.\quad
  \textbf{Surjective} (onto): every $b\in B$ is hit.\quad
  \textbf{Bijective}: both.\\
  \textbf{Everywhere defined}: every $a\in A$ has an image.
\end{keybox}

\begin{example}[5.1 Ex.~12]
  $f_1=\{(a_1,b_2),(a_2,b_3),(a_3,b_1)\}$: everywhere defined, one-to-one,
  onto.
  $f_2=\{(a_1,d_2),(a_2,d_1),(a_3,d_4)\}$ into $D=\{d_1,\dots,d_4\}$:
  one-to-one, not onto.
  $f_3=\{(b_1,c_2),(b_2,c_2),(b_3,c_1)\}$: onto, not one-to-one.
  $f_4=\{(d_1,b_1),(d_2,b_2),(d_3,b_1)\}$: not everywhere defined ($d_4$),
  not one-to-one, not onto.
\end{example}

\subsection{Checking a function from a table (exam method)}
\index{function}\index{everywhere defined}\index{one-to-one|see{injective}}\index{onto|see{surjective}}\index{injective}\index{surjective}\index{bijection}\index{inverse function}\index{composition of functions}%
In Kolman a \textbf{function} $f$ from $A$ to $B$ is a relation where each
$a$ has \emph{at most one} partner $f(a)$ --- a lookup table with no
duplicate keys. Given as a list of pairs, write it as a two-row table
first; every property below is then a glance at one of the rows.
Kolman also says \emph{one-to-one} for injective and
\emph{one-to-one correspondence} for bijection.

\begin{rulepar}
  \begin{center}
    \renewcommand{\arraystretch}{1.3}
    \begin{tabular}{@{}lll@{}}
      \toprule
      Property           & Meaning                               & Check in the table                 \\
      \midrule
      function at all    & no $a$ has two outputs                & no repeat in the \emph{top} row    \\
      everywhere defined & every $a \in A$ has an output         & top row is all of $A$              \\
      injective          & $a \ne a' \Rightarrow f(a) \ne f(a')$ & no repeat in the \emph{bottom} row \\
      surjective (onto)  & every $b \in B$ is hit                & bottom row is all of $B$           \\
      bijection          & all of the above                      & both rows full, no repeats         \\
      \bottomrule
    \end{tabular}
  \end{center}
\end{rulepar}

\textbf{Inverse:} $f^{-1}$ swaps every pair. It is always a relation;
it is a \emph{function} exactly when $f$ is injective (swapping the
rows must not create a repeat in the new top row). To evaluate
$f^{-1}(b)$, find $b$ in the bottom row and read the $a$ above it.

\textbf{Preimage} $f^{-1}(\{b\})$ \index{preimage}\index{inverse image}\index{f inverse of a set@$f^{-1}(\{b\})$}(a \emph{set} inside the braces): the
set of inputs that $f$ sends to $b$. Find $b$ in the bottom row and collect
\emph{every} $a$ above it. It always exists, even when $f$ is not
injective (then the set can have several elements) or not onto (then it
can be $\emptyset$). So $f^{-1}(\{4\})=\{3\}$ for $f(3)=4$, not the
output side $\{f(4)\}$: read \emph{backwards}, bottom row to top row.

\textbf{Composition:} $(g \circ f)(a) = g(f(a))$ --- \textbf{right to
  left}, innermost first, like nested calls \texttt{g(f(a))}. So
$(f \circ f)(3)$ means look up $3$, then look up the answer.

\textbf{Finite shortcut:} for $f\colon A \to A$ everywhere defined on a
\emph{finite} set, injective $\iff$ surjective (pigeonhole: $n$ inputs
with $n$ distinct outputs fill all $n$ slots). One check gives you
bijection, and then $f^{-1}$ is a function too.

\begin{example}[$f = \{(1,2), (2,1), (3,4), (4,5), (5,6), (6,3)\}$ on $\{1,\dots,6\}$: which is false?]
  \leavevmode\par\noindent
  \begin{minipage}[c]{0.36\textwidth}
    \centering
    $\begin{array}{c|cccccc}
        x    & 1 & 2 & 3 & 4 & 5 & 6 \\ \hline
        f(x) & 2 & 1 & 4 & 5 & 6 & 3
      \end{array}$
  \end{minipage}\hfill
  \begin{minipage}[c]{0.6\textwidth}
    \begin{itemize}[nosep]
      \item bottom row $2,1,4,5,6,3$: no repeats, all of $1..6$ ---
            bijection, so injective, onto, everywhere defined and
            ``$f^{-1}$ is a function'' are all \textbf{true}
      \item $f^{-1}(3)$: $3$ is under $6$, so $f^{-1}(3) = 6$ --- true
      \item $(f \circ f)(3) = f(f(3)) = f(4) = 5 \ne 3$ --- \textbf{false}
    \end{itemize}
  \end{minipage}

  \smallskip\noindent
  Only the computed claims needed work; the structural ones fell
  together from one look at the bottom row. Cycle view:
  $f = (1\,2)(3\,4\,5\,6)$, and $f \circ f$ moves $3$ two steps round
  its 4-cycle. Only elements in cycles of length $1$ or $2$ satisfy
  $(f \circ f)(x) = x$.
\end{example}

\begin{example}[MC practice: $f=\{(1,6),(2,3),(3,4),(4,1),(5,5),(6,2)\}$ on $\{1,\dots,6\}$; which is false?]
  Table: $x=1,\dots,6$ gives $f(x)=6,3,4,1,5,2$. The bottom row has no
  repeat and is all of $\{1,\dots,6\}$, so $f$ is everywhere defined,
  injective and onto (a bijection). Then $f(5)=5$ is true;
  $(f\circ f)(1)=f(f(1))=f(6)=2$ is true; but $f^{-1}(\{4\})$ is the input
  that gives $4$, namely $3$ ($f(3)=4$), so $f^{-1}(\{4\})=\{3\}$ and the
  claim ``$=\{1\}$'' is \textbf{false} (it confuses $f(4)=1$ with $f^{-1}(4)$).
\end{example}

\subsection{Injective or surjective from a formula (exam method)}
\index{injective}\index{surjective}\index{function!formula on $\mathbb{Z}$}\index{floor function!injective or surjective}%
For $f\colon\mathbb{Z}\to\mathbb{Z}$ given by a formula, test the two
properties separately and hunt for a \emph{counterexample} first.

\begin{rulebox}[title={Two quick tests}]
  \begin{itemize}[nosep]
    \item \textbf{Not injective}: find $a\neq a'$ with $f(a)=f(a')$.
          Typical culprits: even powers and anything symmetric
          ($f(1)=f(-1)$), and a floor or ceiling, which squeezes several
          inputs into one output.
    \item \textbf{Not surjective}: find an integer that is never an output.
          Typical culprits: a bounded range ($x^2+1\ge1$), a coefficient
          that skips values ($2x$ misses every odd number), or a gap
          between outputs ($(x-1)^3$ gives only perfect cubes).
  \end{itemize}
  If the counterexample hunt fails, prove it: injective by solving
  $f(a)=f(a')$ and getting $a=a'$; surjective by solving $f(x)=y$ for $x$
  and checking that $x$ is an integer for every integer $y$.
\end{rulebox}

\begin{center}
  \renewcommand{\arraystretch}{1.3}
  \begin{tabular}{@{}lll@{}}
    \toprule
    $f(x)$ on $\mathbb{Z}$ & injective? & surjective? \\
    \midrule
    $x+c$                  & yes & yes ($x=y-c$) \\
    $2x$, $2x+1$, $kx$ ($|k|\ge2$) & yes & no (skips values) \\
    $x^2+1$, $x^2$         & no ($f(1)=f(-1)$) & no (nothing negative) \\
    $x^3$, $(x-1)^3$       & yes & no (only cubes) \\
    $\lfloor x/2\rfloor$   & no ($f(2n)=f(2n+1)$) & yes (hits every $n$) \\
    \bottomrule
  \end{tabular}
\end{center}
\textbf{Floor check:} $\lfloor x/2+\tfrac12\rfloor$ is $n$ for both
$x=2n-1$ and $x=2n$ (not injective), and every $n$ is reached
(surjective). \textbf{Increasing functions} (strictly rising on $\mathbb{Z}$,
such as $x\cdot2^{x^2}$) are injective, but the outputs jump, so they are
rarely onto.

\subsection{Inverses}
\index{inverse function}\index{inverse function!of a composition}%
$f^{-1}$ reverses the pairs: $f(a)=b\iff f^{-1}(b)=a$.
\begin{rulebox}
  \textbf{Thm 5.1.1:} $f^{-1}$ is a function $\iff$ $f$ is bijective.\\
  \textbf{Thm 5.1.2:} $1_B\circ f=f$, \ $f\circ1_A=f$; for bijective $f$:
  $f^{-1}\circ f=1_A$, \ $f\circ f^{-1}=1_B$.\\
  \textbf{Thm 5.1.3:} (a) if $g\circ f=1_A$ and $f\circ g=1_B$ then $f,g$
  are bijections and $g=f^{-1}$.\
  (b) $(g\circ f)^{-1}=f^{-1}\circ g^{-1}$ (reverse order, like
  $(AB)^T=B^TA^T$).
\end{rulebox}

\begin{example}[5.1 Ex.~17]
  $f(x)=2x^3-1$, $g(y)=\sqrt[3]{(y+1)/2}$ on $\R$:
  $g(f(x))=\sqrt[3]{x^3}=x$ and $f(g(y))=2\cdot\frac{y+1}2-1=y$, so both are
  bijections and $g=f^{-1}$.
\end{example}
\begin{example}[5.1.24--25, is $f^{-1}$ a function?]
  $f=\{(1,a),(2,a),(3,c),(4,d)\}$: $f^{-1}$ sends $a$ to both $1$ and $2$
  --- \textbf{not} a function.
  $f=\{(1,a),(2,c),(3,b),(4,d)\}$: bijective, $f^{-1}$ \textbf{is} a
  function.
\end{example}
\begin{example}[5.1 Ex.~9 and 5.1.10]
  $f(a)=a+1$, $g(b)=2b$: $(g\circ f)(a)=2a+2$.\\
  $f(a)=a+1$, $g(b)=b^2+2$:
  $(g\circ f)(-2)=g(-1)=3$;\ $(f\circ g)(-2)=f(6)=7$;\
  $(g\circ f)(x)=x^2+2x+3$;\ $(f\circ g)(x)=x^2+3$;\
  $(f\circ f)(y)=y+2$;\ $(g\circ g)(y)=y^4+4y^2+6$.
\end{example}
\begin{example}[MC practice test 5: composition read off a digraph]\leavevmode\par\noindent
  \begin{minipage}[c]{0.2\textwidth}
    \centering
    \begin{tikzpicture}[every node/.style={circle,draw,minimum size=6mm,inner sep=0pt,font=\small},>=stealth]
      \node (1) at (0,1.3) {1};
      \node (2) at (1.3,1.3) {2};
      \node (3) at (0,0) {3};
      \node (4) at (1.3,0) {4};
      \draw[->] (1) to[loop above,looseness=8] (1);
      \draw[->] (3)--(1);
      \draw[<->] (4)--(2);
    \end{tikzpicture}
  \end{minipage}\hfill
  \begin{minipage}[c]{0.77\textwidth}
    $F$ on $A=\{1,2,3,4\}$ is given by the digraph: $F(1)=1$, $F(3)=1$,
    $F(2)=4$, $F(4)=2$ (the arrow between $2$ and $4$ points both ways).
    $G=\{(1,c),(2,d),(3,e),(4,a)\}$. Then
    $(G\circ F)(2)=G(F(2))=G(4)=a$. Apply $F$ \emph{first}, then $G$.
  \end{minipage}
\end{example}

% =====================================================================
\kolman{2}
\section{Functions for computer science}\src{Kolman 5.2}
\index{floor and ceiling}\index{hash function}\index{fuzzy set}%

Common examples: characteristic function $\chi_A$, $m\bmod n$, $n!$,
floor $\lfloor x\rfloor$, ceiling $\lceil x\rceil$, string length,
power set, polynomials, $2^x$, $\log_bx$, transpose, $\gcd$, $\lcm$.
Boolean-valued functions have codomain $\{\text{true},\text{false}\}$.

\textbf{Hash function}: $h(k)=k\bmod m$ maps key $k$ to one of $m$ table
positions. A \textbf{collision} is $k_1\neq k_2$ with $h(k_1)=h(k_2)$:
with $m=10$, $h(27)=h(37)=7$.

\textbf{Fuzzy set}: membership $\mu_A(x)\in[0,1]$ ($0$ not in, $1$ fully
in, between = partly). E.g.\ ``tall'': 160\,cm $0.1$, 170 $0.3$, 180 $0.7$,
190 $1.0$.
\begin{rulebox}[title=Fuzzy operations (Theorem 5.2.1)]
  \index{fuzzy set!min and max}%
  $(f\cap g)(x)=\min\{f(x),g(x)\}$ \qquad
  $(f\cup g)(x)=\max\{f(x),g(x)\}$ \qquad
  $\overline f(x)=1-f(x)$
\end{rulebox}
\begin{example}[5.2 Ex.~13, answers not on slides]
  On $[0,10]$:
  $f(x)=\frac{x^2}{16}$ on $[0,4)$, $1$ on $[4,6)$, $-\frac x2+4$ on $[6,8)$,
  $0$ on $[8,10]$;\
  $g(x)=\frac x3$ on $[0,3)$, $-\frac x4+\frac74$ on $[3,5)$, $\frac12$ on
  $[5,7)$, $-\frac x6+\frac{10}6$ on $[7,10]$.\\
  $f(3.5)=\frac{49}{64}$, $g(3.5)=\frac78$; $f(7)=g(7)=\frac12$.
  So (a) $3.5$ in $f\cap g$: $\frac{49}{64}$; (b) $7$ in $f\cap g$:
  $\frac12$; (c) $3.5$ in $f\cup g$: $\frac78$; (d) $7$ in $f\cup g$: $\frac12$.
\end{example}

\begin{example}[5.2.22--23, propositional functions]
  $P(x,y)=(x\vee y)\wedge\neg y$: $P(T,T)=F$, $P(F,T)=F$, $P(T,F)=T$.\\
  $Q(x):\exists y\in\Z^+,\ xy=60$: $Q(3)=T$ ($y=20$), $Q(7)=F$,
  $Q(-6)=F$ ($y=-10\notin\Z^+$), $Q(15)=T$ ($y=4$).
\end{example}

% =====================================================================
\kolman{3}
\section{Growth of functions}\src{Kolman 5.3}
\index{growth of functions}\index{order of growth}\index{Big-O}\index{Big-Theta@Big-$\Theta$}\index{Theta-class@$\Theta$-class}\index{same order of growth}\index{polynomial versus exponential}\index{doubling the input}%

\begin{rulebox}[title={Growth order (slowest to fastest, for large $n$)}]
  \index{order of growth!ranking of common functions}%
  \centering
  $1<\log n<n<n\log n<n^2<n^3<2^n$\\
  constant $\to$ logarithmic $\to$ linear $\to$ linearithmic $\to$ polynomial
  $\to$ exponential
\end{rulebox}

\begin{defn}[Big-O]
  $f=O(g)$ if there are constants $c>0$, $k\in\Z^+$ with
  $|f(n)|\le c\,|g(n)|$ for all $n\ge k$. An \emph{upper bound}: $f$ grows
  no faster than $g$. To prove it, exhibit a $c$ and $k$.
\end{defn}
\begin{defn}[Big-$\Theta$]
  $f=\Theta(g)$ if $f=O(g)$ and $g=O(f)$; equivalently
  $c_1|g(n)|\le|f(n)|\le c_2|g(n)|$ for $n\ge k$. \emph{Same order of
    growth.} Theorem 5.3.1: $\Theta$ is an equivalence relation (reflexive,
  symmetric, transitive), so functions fall into $\Theta$-classes.
\end{defn}

$100n^2=O(n^3)$ but $100n^2=\Theta(n^2)$, not $\Theta(n^3)$.

\begin{rulepar}
  \textbf{Finding the $\Theta$-class}: drop constant factors and lower-order
  terms; the highest power dominates a polynomial; floor/ceiling don't
  change the class; any exponential $a^n$ ($a>1$) beats any polynomial.
\end{rulepar}

\begin{example}[Big-O proofs]
  $100n^2\le100n^3$ for $n\ge1$: $c=100$, $k=1$.\\
  $2n^3+n^2+7\le2n^3+n^3+n^3=4n^3$ for $n\ge2$: $c=4$, $k=2$.\\
  $n!=1\cdot2\cdots n\le n\cdot n\cdots n=n^n$: $c=1$, $k=1$.\\
  $1+\dots+n=\frac{n(n+1)}2\le\frac{n\cdot2n}2=n^2$: $c=1$, $k=1$.
\end{example}
\begin{example}[5.3.9: $5n^2+4n+3$ and $n^2+100n$ have the same order]
  $f\le5n^2+4n^2+3n^2=12n^2\le12g$, and
  $g\le n^2+100n^2=101n^2\le\frac{101}5f$. So $f=\Theta(g)=\Theta(n^2)$.
\end{example}
\begin{example}[5.3 Ex.~7--8]
  $4n^4-6n^7+25n^3\in\Theta(n^7)$;\ $\lfloor\lg n\rfloor-3n\in\Theta(n)$;\
  $1.1^n+n^{15}\in\Theta(1.1^n)$.\\
  Ordering:
  $\Theta(1{,}000{,}000)=\Theta(1)<\Theta(n^{0.2})<\Theta(n+10^7)=\Theta(n)
    <\Theta(n\lg n)<\Theta(1000n^2-n)=\Theta(n^2)<\Theta(1.3^n)$.
\end{example}
\begin{example}[MC practice test 5: highest order of growth]
  Options: $3^n+n^7$, $2^n+n^7$, $n^7$, $n^8+\log_2n$, $\log_2n$,
  $n^2\cdot3^n+\log_2n$. Keep each dominant term: $3^n$, $2^n$, $n^7$,
  $n^8$, $\log_2n$, $n^23^n$. Since $\log n<n^k<a^n$ ($a>1$) and
  $2^n<3^n<n^2\,3^n$, the answer is $n^2\cdot3^n+\log_2n$. A polynomial
  factor in front of an exponential still counts.
\end{example}
\begin{example}[MC practice: which option has the same order as $n^2$?]
  \index{same order of growth!picking from options}\index{dominant term}%
  Reduce each option to its dominant term, then compare with $n^2$; only
  $\Theta(n^2)$ matches. $n^2+2^n\to2^n$ and $n^2+1.1^n\to1.1^n$ (any
  exponential wins, even base $1.1$); $n^3-n^2\to n^3$;
  $8n^2\log_2n\to n^2\log n$ (extra log factor, strictly faster);
  $4n\log_2n\to n\log n$ (strictly slower);
  $4n^2+n\log_2n\to4n^2$, since $\frac{n\log n}{n^2}=\frac{\log n}n\to0$
  --- so it is $\Theta(n^2)$. Constant factors (the $4$, the $8$) never
  change the class; a $\log$ factor \emph{does}.
\end{example}

\begin{example}[MC practice: which has the same order as $n^3+2^n\log_2n$?]
  \index{same order of growth!log base}\index{logarithm!change of base}\index{exponential times logarithm}%
  Dominant term: $2^n\log_2n$ ($n^3$ is much smaller than any exponential).
  Two facts decide the options:
  \begin{itemize}[nosep]
    \item \textbf{The log base never matters:} $\log_{10}n=\dfrac{\log_2n}{\log_210}$
          is a constant multiple of $\log_2n$, so $2^n\log_{10}n=\Theta(2^n\log_2n)$.
          (Constant factors, even $\frac14$ or $\frac1{\log_210}$, never change the class.)
    \item \textbf{Exponential $\times$ log is its own class:} $2^n\log_2n$ is
          strictly above $2^n$ (ratio $\log_2n\to\infty$) and strictly
          below $2^nn^3$ and $4^n$. So $2^n$ and $n^3+2^n$ (both
          $\Theta(2^n)$) do \emph{not} match.
  \end{itemize}
  Answer: $2^n\log_{10}n$. The distractors: $2^n$ and $n^3+2^n$ (log
  factor dropped), $\frac14 4^n\log_2n$ (a bigger base), $2^nn^3$ (a
  polynomial factor, not a log), $n^3$ (too small).
\end{example}

\begin{example}[5.3.1: what doubling the input does]\leavevmode\\
  \begin{tabular}[t]{@{}ll@{\qquad}ll@{}}
    $1001$   & no change & $1.4\lg n$ & $f(2n)=f(n)+1.4$               \\
    $3n$     & $\times2$ & $n\lg n$   & $f(2n)=2f(n)+2n$               \\
    $5n^2$   & $\times4$ & $100n^4$   & $\times16$                     \\
    $2.5n^3$ & $\times8$ & $2^n$      & $f(2n)=f(n)^2$ (not on slides)
  \end{tabular}
\end{example}

\subsection{Limits and L'Hôpital}
\index{Lhopital's rule@L'H\^opital's rule}\index{limit}\index{derivative!table of derivatives}%
\begin{rulebox}
  $\displaystyle\lim_{n\to\infty}\frac{f(n)}{g(n)}=L$ with $0<L<\infty$
  $\Rightarrow f=\Theta(g)$;\qquad
  limit $=0$ $\Rightarrow$ $f$ grows strictly slower than $g$.\\[4pt]
  \textbf{L'Hôpital}: if $\frac{f}{g}$ has the form $\frac00$ or
  $\frac\infty\infty$, then
  $\displaystyle\lim\frac{f(x)}{g(x)}=\lim\frac{f'(x)}{g'(x)}$ (if the
  right side exists).
\end{rulebox}

\begin{rulepar}
  \begin{center}
    \begin{tabular}{@{}c|cccccc@{}}
      $f(x)$  & $c$ & $x$ & $x^n$      & $\ln x$ & $e^x$ & $\sin x$ \quad $\cos x$  \\\hline
      $f'(x)$ & $0$ & $1$ & $nx^{n-1}$ & $1/x$   & $e^x$ & $\cos x$ \quad $-\sin x$
    \end{tabular}
    \qquad e.g.\ $(5x^4)'=20x^3$
  \end{center}
\end{rulepar}

\begin{example}[limits]
  $\lim\frac{100n^2}{n^3}=\lim\frac{100}n=0$: $100n^2$ grows slower than $n^3$.\\
  $\lim\frac{2n^3+n^2+7}{n^3}=\lim\left(2+\frac1n+\frac7{n^3}\right)=2$:
  $\Theta(n^3)$.\\
  $\lim_{x\to0}\frac{\sin x}x$ ($\frac00$) $=\lim\frac{\cos x}1=1$.\\
  $\lim_{x\to\infty}\frac{\ln x}{x^k}$ ($\frac\infty\infty$, $k>0$)
  $=\lim\frac{1/x}{kx^{k-1}}=\lim\frac1{kx^k}=0$: $\ln x$ grows slower
  than every positive power of $x$.
\end{example}

% =====================================================================
\newpage
\part{Order Relations and Structures}

\kolman{1}
\section{Partial orders and posets}\src{Kolman 6.1}
\index{partial order}\index{poset}\index{antisymmetric}\index{reflexive}\index{transitive}%

\begin{defn}
  A relation $R$ on a set $A$ is a \textbf{partial order} if it is
  reflexive, antisymmetric and transitive. The pair $(A,R)$ (written
  $(A,\le)$ when $R$ is denoted $\le$) is called a \textbf{poset}
  (partially ordered set).
\end{defn}

\begin{example}[6.1 Ex.~1--3]
  $(\mathcal P(S),\subseteq)$ is a poset for any set $S$ (Example 1);
  $(\Z^{+},\le)$ is a poset (Example 2); divisibility, $aRb\iff a\mid b$,
  is a poset on $\Z^{+}$ (Example 3) --- each of reflexivity,
  antisymmetry and transitivity for $\le$ and $\mid$ is the familiar
  number fact.
\end{example}

\begin{keybox}[title={Strict order: why $<$ is not a partial order}]
  \index{strict order}\index{asymmetric}\index{irreflexive}%
  $<$ fails reflexivity ($a<a$ is never true), so it is not a partial
  order on $\Z^+$/$\R$ --- but it is a useful relation in its own right,
  a \textbf{strict partial order}. Two more relation properties: \textbf{Irreflexive}: $\forall x\in A,\ \neg(xRx)$ (no
  element related to itself). \textbf{Asymmetric}: $xRy\Rightarrow
    \neg(yRx)$ (never both ways).
  \emph{Asymmetric and antisymmetric are not the same}: antisymmetric
  allows $xRy\wedge yRx$ when $x=y$ (as $\le$ does); asymmetric forbids
  $xRy\wedge yRx$ outright, which is exactly why $<$ (irreflexive,
  transitive) is asymmetric rather than merely antisymmetric.
\end{keybox}

\subsection{Comparability, chains, duals}
\index{comparable elements}\index{incomparable elements}\index{chain (linear order)}\index{dual poset}%
Two elements $a,b\in A$ are \textbf{comparable} if $a\le b$ or $b\le a$;
otherwise \textbf{incomparable}. Under divisibility, $2$ and $6$ are
comparable ($2\mid6$) but $2$ and $3$ are incomparable ($2\nmid3$,
$3\nmid2$). If \emph{every} pair is comparable, $(A,\le)$ is
\textbf{linearly ordered} (a \textbf{chain}); e.g.\ $(\Z,\le)$, but not
divisibility in general.

\begin{rulebox}[title={Is it a linear order? Counterexample $\mathcal P(S)$}]
  \index{linear order!$\mathcal P(S)$ is not}\index{power set!not linearly ordered}\index{incomparable elements!counterexample}%
  To show a poset is \textbf{not} linear, exhibit \emph{one} incomparable
  pair. For $(\mathcal P(S),\subseteq)$ with $|S|\ge2$, pick distinct
  $a,b\in S$: $\{a\}\not\subseteq\{b\}$ and $\{b\}\not\subseteq\{a\}$, so
  $\{a\},\{b\}$ are incomparable and $\subseteq$ is \textbf{not} a linear
  order. (It is still a partial order --- Example~1.) If $|S|\le1$ the
  power set \emph{is} a chain ($\emptyset\subseteq\{a\}$), so state the
  case $|S|\ge2$. To show it \emph{is} linear you must prove \emph{every}
  pair comparable, which needs a general argument, not an example.
\end{rulebox}

The \textbf{dual poset} of $(A,\le)$ is $(A,\ge)$ --- same set, order
reversed; e.g.\ the dual of $(\mathcal P(S),\subseteq)$ is
$(\mathcal P(S),\supseteq)$.

\subsection{Product and lexicographic order}
\index{product order}\index{lexicographic order}%
\begin{rulebox}[title=Theorem 1 (product partial order)]
  For posets $(A,\le)$, $(B,\le)$, define an order on $A\times B$ by
  \[(a,b)\le(a',b') \iff a\le a' \text{ and } b\le b'.\]
\end{rulebox}
$(2,3)\le(4,5)$ since $2\le4$ and $3\le5$. But $(2,5)$ and $(3,4)$ are
\textbf{incomparable} under the product order: $(2,5)\le(3,4)$ needs
$5\le4$ (false), and $(3,4)\le(2,5)$ needs $3\le2$ (false).

\begin{keybox}[title=Lexicographic (dictionary) order]
  $(a,b)\prec(a',b')$ iff $a<a'$, or $a=a'$ and $b<b'$ --- compare the
  first component; only if it ties, compare the second. Unlike the
  product order, every pair of pairs is comparable this way.
\end{keybox}
$(2,100)$ vs.\ $(3,1)$: \textbf{product order} --- $(2,100)\le(3,1)$
needs $100\le1$ (false), so incomparable. \textbf{Lexicographic order}
--- $2<3$ settles it immediately: $(2,100)\prec(3,1)$.
\index{lexicographic order}\index{order relation!lexicographic}%
Like sorting by surname, then first name: \texttt{ORDER BY a, b}.

\textbf{Consequence:} a larger first coordinate wins outright --- the
second coordinate is never even looked at. $(5, \emptyset)$ is above
$(2, \{a, b, c\})$, however big that set is.

\textbf{Trap when $B$ is only a partial order} (e.g.\ $\subseteq$ on a
power set): two pairs with the same first coordinate can be
\emph{incomparable}. $\{a\}$ and $\{b\}$: neither $\subseteq$ the other,
so $(3, \{a\})$ and $(3, \{b\})$ are incomparable.


\begin{rulebox}[title=Theorem 2 (no cycles)]
  A partial order's digraph has no cycle of length greater than $1$.
  \emph{Proof idea}: a cycle $2\le5\le8\le2$ gives, by transitivity,
  $2\le8$ and $8\le2$; antisymmetry then forces $2=8$, contradiction
  (distinct elements).
\end{rulebox}

\subsection{Exam method: is a relation a partial order?}
\index{partial order!matrix test}\index{partial order!given by a formula}\index{relation!given by a formula}\index{partial order!which relation on $\mathbb{Z}^+$?}\index{empty relation}\index{same parity}\index{transitive!row trick}\index{equivalence relation}%
Given as a matrix, or as a formula on $\mathbb{Z}^+$:

\begin{rulepar}
  \textbf{Matrix test, cheapest check first} --- stop at the first failure:
  \begin{enumerate}[nosep]
    \item \textbf{Reflexive:} the whole diagonal is $1$. A single $0$ on
          the diagonal and it is \emph{not} an order. (Strict orders like
          $<$ have diagonal all $0$ --- not partial orders in Kolman's
          sense.)
    \item \textbf{Antisymmetric:} for every off-diagonal $1$ at
          $(i,j)$, the mirror cell $(j,i)$ is $0$. Fold the matrix along
          the diagonal: no two $1$'s may land on top of each other.
          Shortcut: if everything below the diagonal is $0$ (upper
          triangular), it is automatically antisymmetric.
    \item \textbf{Transitive --- row trick:} for every $1$ at $(i,j)$
          with $i \ne j$, every $1$ in row $j$ must also be a $1$ in row
          $i$ (``if $i$ reaches $j$, then $i$ reaches everything $j$
          reaches''). Off-diagonal $1$'s only; the diagonal ones are
          trivially fine.
  \end{enumerate}
  \index{transitive!row trick}%
  Do the transitivity check last: it is the slowest, and steps 1--2
  usually eliminate most MC options.
\end{rulepar}

\begin{example}[Is ``divides'' on $\{1, 2, 3, 6\}$ a partial order?]
  \leavevmode\par\noindent
  \begin{minipage}[c]{0.28\textwidth}
    \centering
    $\begin{array}{c|cccc}
          & 1 & 2 & 3 & 6 \\ \hline
        1 & 1 & 1 & 1 & 1 \\
        2 & 0 & 1 & 0 & 1 \\
        3 & 0 & 0 & 1 & 1 \\
        6 & 0 & 0 & 0 & 1
      \end{array}$
  \end{minipage}\hfill
  \begin{minipage}[c]{0.7\textwidth}
    \begin{enumerate}[nosep]
      \item reflexive: diagonal $1,1,1,1$ --- yes
      \item antisymmetric: everything below the diagonal is $0$ ---
            yes, no mirror pair can be $1/1$
      \item transitive, row trick over the off-diagonal $1$'s:
            \begin{itemize}[nosep]
              \item $(1,2)$, $(1,3)$, $(1,6)$: row 1 is all $1$'s,
                    so it contains any row --- fine
              \item $(2,6)$: row 6 $= 0001 \subseteq$ row 2 $= 0101$ --- fine
              \item $(3,6)$: row 6 $= 0001 \subseteq$ row 3 $= 0011$ --- fine
            \end{itemize}
    \end{enumerate}
    All three hold: \textbf{a partial order}. Had row 2 been $0100$
    instead, $(1,2)$ would still pass but $(2,6)$ would vanish ---
    always recheck after a change.
  \end{minipage}
\end{example}

\noindent\textbf{Failing transitivity looks like this:} rows $1{:}\ 1100$,
$2{:}\ 0110$, $3{:}\ 0010$, $4{:}\ 0001$. Reflexive and antisymmetric
pass, but $(1,2)$ is a $1$ and row 2 has a $1$ in column 3 that row 1
lacks: $1 \to 2 \to 3$ with no shortcut $(1,3)$.

\textbf{Relations given by a formula} (``which relation on $\mathbb{Z}^+$
\index{partial order!which relation on $\mathbb{Z}^+$?}\index{empty relation}\index{same parity}%
is a partial order?''): first rewrite each condition into something you
recognise, then run reflexive $\to$ antisymmetric $\to$ transitive.
\begin{itemize}[nosep]
  \item $x - y \le 0$ is just $x \le y$ --- the usual order, so a
        partial order. $x - y < 0$ is $x < y$: strict, so
        \emph{not} reflexive.
  \item \textbf{Empty or everything:} on positive numbers $x + y \le 0$
        and $x + y < 0$ are never true (empty relation, not reflexive),
        and $xy \ge 0$ is always true (every pair related, so
        $1\,R\,2$ and $2\,R\,1$ breaks antisymmetry).
  \item $2 \mid (x + y)$ means ``same parity'': reflexive and symmetric
        (an equivalence relation), but $1\,R\,3$ and $3\,R\,1$ with
        $1 \ne 3$ kills antisymmetry.
\end{itemize}
A symmetric relation that relates any two \emph{different} elements is
never antisymmetric.

\subsection{Hasse diagrams}
\index{Hasse diagram}\index{cover relation}%
Drawing every self-loop (reflexivity) and every transitive edge makes a
poset's digraph unreadable. A \textbf{Hasse diagram} simplifies it:
remove self-loops, remove edges implied by transitivity (keep only the
\textbf{cover relations}), place larger elements higher, and draw no
arrowheads --- direction is ``upward.'' Read it \textbf{bottom to top}:
an upward path from $a$ to $b$ means $a\le b$. The \textbf{dual} Hasse
diagram is the same picture upside down (order reversed).

\begin{example}[6.1 Ex.~12: divisibility on $\{1,2,3,4,12\}$]
  Divisibility pairs: $1\mid2,\,1\mid3,\,1\mid4,\,1\mid12,\,2\mid4,\,
    2\mid12,\,3\mid12,\,4\mid12$. Dropping what transitivity already gives
  (e.g.\ $1\mid2\mid4\Rightarrow1\mid4$ is redundant, likewise $1\mid12$,
  $2\mid12$) leaves the cover relations
  $1\prec2,\ 1\prec3,\ 2\prec4,\ 3\prec12,\ 4\prec12$:
  \begin{center}
    \begin{tikzpicture}[every node/.style={circle,draw,minimum size=6mm,inner sep=0pt,font=\small}]
      \node (n1) at (0,0) {1};
      \node (n2) at (-1,1) {2};
      \node (n3) at (1,1) {3};
      \node (n4) at (-1,2) {4};
      \node (n12) at (0,3) {12};
      \draw (n1)--(n2) (n1)--(n3) (n2)--(n4) (n4)--(n12) (n3)--(n12);
    \end{tikzpicture}
  \end{center}
\end{example}

\begin{example}[power set of $\{a,b,c\}$ under $\subseteq$]
  The familiar cube: singletons cover $\emptyset$, pairs cover the
  singletons they contain, and $\{a,b,c\}$ covers every pair.
  \begin{center}
    \begin{tikzpicture}[every node/.style={rectangle,draw,rounded corners,inner sep=2pt,font=\small}]
      \node (e)   at (0,0)    {$\emptyset$};
      \node (a)   at (-2,1.2) {$\{a\}$};
      \node (b)   at (0,1.2)  {$\{b\}$};
      \node (c)   at (2,1.2)  {$\{c\}$};
      \node (ab)  at (-2,2.4) {$\{a,b\}$};
      \node (ac)  at (0,2.4)  {$\{a,c\}$};
      \node (bc)  at (2,2.4)  {$\{b,c\}$};
      \node (abc) at (0,3.6)  {$\{a,b,c\}$};
      \draw (e)--(a) (e)--(b) (e)--(c);
      \draw (a)--(ab) (a)--(ac) (b)--(ab) (b)--(bc) (c)--(ac) (c)--(bc);
      \draw (ab)--(abc) (ac)--(abc) (bc)--(abc);
    \end{tikzpicture}
  \end{center}
\end{example}

\subsection{Exam method: from a matrix to a Hasse diagram}
\index{Hasse diagram!from a matrix}\index{partial order!matrix to Hasse diagram}\index{cover relation!from a matrix}\index{matrix of a relation!to Hasse diagram}%
Given the matrix $M_R$ of a partial order on $A=\{1,\dots,n\}$ (row $i$,
column $j$ is $1$ iff $i\,R\,j$, i.e.\ $i\le j$):

\begin{rulepar}
  \begin{enumerate}[nosep]
    \item \textbf{Read the pairs.} Ignore the diagonal (loops). Every
          other $1$ at $(i,j)$ is a relation $i\le j$.
    \item \textbf{Cross out transitive edges.} Keep $i\to j$ only if
          there is \emph{no} $k$ ($k\ne i,j$) with $i\le k$ and $k\le j$.
          Matrix trick: $i\to j$ is redundant iff some $k\ne i,j$ has a
          $1$ in row $i$ (so $i\le k$) \emph{and} a $1$ in column $j$ (so
          $k\le j$) --- a $k$ whose entry in row $i$ and entry in column
          $j$ are both $1$.
    \item \textbf{Find the bottom and top.} \emph{Minimal} elements: the
          column has only its diagonal $1$. \emph{Maximal} elements: the
          row has only its diagonal $1$.
    \item \textbf{Assign heights.} Minimal elements at the bottom; every
          other element goes one level above the highest element it
          covers. Draw the covers as plain upward lines, no arrowheads.
  \end{enumerate}
  Sanity check: row sum of $i$ minus $1$ $=$ number of elements above $i$;
  column sum of $j$ minus $1$ $=$ number below $j$. The picture must agree.
\end{rulepar}

\begin{example}[a branch]
  \[
    M=\begin{bmatrix}
      1 & 1 & 1 & 1 & 1 \\ 0&1&0&0&1\\ 0&0&1&1&1\\ 0&0&0&1&1\\ 0&0&0&0&1
    \end{bmatrix}
  \]
  Pairs: $1\le2,3,4,5$; $2\le5$; $3\le4,5$; $4\le5$. Redundant:
  $1\to4$ (via $3$), $1\to5$ (via $2$), $3\to5$ (via $4$). Covers:
  $1\prec2,\ 1\prec3,\ 3\prec4,\ 2\prec5,\ 4\prec5$. Minimal: $1$ only
  (column~1 is $1,0,0,0,0$); maximal: $5$ only. Heights: $1$ bottom, $2$
  and $3$ above it, $4$ above $3$, $5$ on top.
  \begin{center}
    \begin{tikzpicture}[every node/.style={circle,draw,minimum size=6mm,inner sep=0pt,font=\small}]
      \node (n1) at (0,0) {1};
      \node (n2) at (-1,1) {2};
      \node (n3) at (1,1) {3};
      \node (n4) at (1,2) {4};
      \node (n5) at (0,3) {5};
      \draw (n1)--(n2) (n1)--(n3) (n3)--(n4) (n2)--(n5) (n4)--(n5);
    \end{tikzpicture}
  \end{center}
\end{example}

\begin{example}[two minimal elements]
  \[
    M=\begin{bmatrix}
      1 & 0 & 1 & 1 & 1 \\ 0&1&1&1&1\\ 0&0&1&1&1\\ 0&0&0&1&0\\ 0&0&0&0&1
    \end{bmatrix}
  \]
  Pairs: $1\le3,4,5$; $2\le3,4,5$; $3\le4,5$. Redundant: $1\to4,\
    1\to5,\ 2\to4,\ 2\to5$ (each goes through $3$). Covers: $1\prec3,\
    2\prec3,\ 3\prec4,\ 3\prec5$. Minimal: $1$ and $2$ (columns $1,2$ have
  only a diagonal $1$); maximal: $4$ and $5$ (rows $4,5$ likewise), so
  there is neither a least nor a greatest element. $4,5$ are
  incomparable (both zeros at $(4,5)$ and $(5,4)$), so they sit side by
  side.
  \begin{center}
    \begin{tikzpicture}[every node/.style={circle,draw,minimum size=6mm,inner sep=0pt,font=\small}]
      \node (n1) at (-1,0) {1};
      \node (n2) at (1,0) {2};
      \node (n3) at (0,1) {3};
      \node (n4) at (-1,2) {4};
      \node (n5) at (1,2) {5};
      \draw (n1)--(n3) (n2)--(n3) (n3)--(n4) (n3)--(n5);
    \end{tikzpicture}
  \end{center}
\end{example}

\subsection{Topological sorting}
\index{topological sorting}%
A poset only orders \emph{some} pairs; a \textbf{topological sort}
extends it to a full linear ordering that respects every $a\preceq b$
already fixed by the partial order. If $A\prec C$, $B\prec C$,
$C\prec D$ with $A,B$ incomparable, both $A,B,C,D$ and $B,A,C,D$ are
valid topological sorts --- \textbf{a topological sort need not be
  unique}.

\subsection{Isomorphism of posets}
\index{isomorphism!of posets}%
\begin{defn}
  A bijection $f:A\to A'$ between posets is an \textbf{isomorphism} if
  $a\le b\iff f(a)\le f(b)$ for all $a,b\in A$ --- the element names
  differ, but the order structure is identical.
\end{defn}
\begin{example}[6.1 Ex.~17]
  $A=\{1,2,3,6\}$ under $\mid$ and $A'=\{\emptyset,\{a\},\{b\},\{a,b\}\}$
  under $\subseteq$ are isomorphic via
  $1\mapsto\emptyset,\ 2\mapsto\{a\},\ 3\mapsto\{b\},\ 6\mapsto\{a,b\}$.
\end{example}

\begin{keybox}[title=Key points]
  Partial order $=$ reflexive $+$ antisymmetric $+$ transitive; a set
  with one is a \textbf{poset}. Pairs are \textbf{comparable} or not; if
  all pairs are, the order is \textbf{linear} (a chain). A \textbf{Hasse
    diagram} drops loops and transitive edges. A \textbf{topological
    sort} extends a partial order to a linear one, not necessarily
  uniquely. \textbf{Isomorphic} posets share order structure regardless
  of element names.
\end{keybox}

\begin{example}[Ex.~21--24: divisibility posets, linearly ordered?]\leavevmode\\[2pt]
  \begin{tabular}[t]{@{}ll@{}}
    $A_1=\{1,2,3,5,6,10,15,30\}$       & not linear ($2\nmid3$, $3\nmid2$)                  \\
    $A_2=\{2,4,8,16,32\}$              & \textbf{linear} --- a chain, each divides the next \\
    $A_3=\{3,6,12,36,72\}$             & \textbf{linear} --- a chain, each divides the next \\
    $A_4=\{1,\dots,6,10,12,15,30,60\}$ & not linear ($2\nmid3$, $3\nmid2$)
  \end{tabular}\\[2pt]
  A set of divisors forms a chain exactly when it can be listed so each
  divides the next; branching (two elements dividing a common one but
  not each other) breaks linearity.
\end{example}

% =====================================================================
\kolman{2}
\section{Extremal elements of posets}\src{Kolman 6.2}
\index{greatest element}\index{least element}\index{maximal element}\index{minimal element}%

\begin{defn}
  Let $(A,\le)$ be a poset and $a\in A$. Going \emph{up}:
  \begin{itemize}
    \item $a$ is the \textbf{greatest} element if $x\le a$ for all
          $x\in A$ --- it is above everything;
    \item $a$ is \textbf{maximal} if there is no $x\in A$ with $a<x$ ---
          nothing is strictly above it.
  \end{itemize}
  Going \emph{down}: $a$ is the \textbf{least} element if $a\le x$ for
  all $x\in A$, and \textbf{minimal} if there is no $x\in A$ with $x<a$.
\end{defn}

\begin{keybox}[title={Greatest vs.\ maximal}]
  \index{greatest element!versus maximal}%
  \begin{minipage}{0.62\textwidth}
    A greatest element is always maximal, but a maximal element need not
    be greatest. A greatest element, if it exists, is \textbf{unique}
    (Theorem 2); a poset may have \textbf{several} maximal elements, and
    then it has no greatest one. Left: $d$ is greatest (and maximal).
    Right: $b$ and $c$ are both maximal, so there is no greatest element.

    \smallskip
    \textbf{Reading a Hasse diagram:} maximal = no edge leaves it upward;
    minimal = no edge leaves it downward.

    \smallskip
    \textbf{Not on the slides.} In a \emph{finite} poset, a greatest
    element exists exactly when there is only \emph{one} maximal element
    (follow edges upward from anywhere and you must land on it). This
    fails for infinite posets.
  \end{minipage}\hfill
  \begin{minipage}{0.34\textwidth}
    \centering
    \begin{tikzpicture}[hasse]
      \node[pt,label=below:$a$] (a) at (0,0) {};
      \node[pt,label=left:$b$]  (b) at (-0.8,1) {};
      \node[pt,label=right:$c$] (c) at (0.8,1) {};
      \node[pt,label=above:$d$] (d) at (0,2) {};
      \draw (a)--(b)--(d)--(c)--(a);
    \end{tikzpicture}\qquad
    \begin{tikzpicture}[hasse]
      \node[pt,label=below:$a$] (a) at (0,0) {};
      \node[pt,label=above:$b$] (b) at (-0.8,1.5) {};
      \node[pt,label=above:$c$] (c) at (0.8,1.5) {};
      \draw (b)--(a)--(c);
    \end{tikzpicture}
  \end{minipage}
\end{keybox}

\begin{example}[Ex.~1: a zigzag]\leavevmode\par\noindent
  \begin{minipage}[c]{0.45\textwidth}
    \centering
    \begin{tikzpicture}[hasse]
      \node[pt,label=above:$a_1$] (a1) at (0,1.4) {};
      \node[pt,label=above:$a_2$] (a2) at (2,1.4) {};
      \node[pt,label=above:$a_3$] (a3) at (3.2,1.4) {};
      \node[pt,label=below:$b_1$] (b1) at (0.8,0) {};
      \node[pt,label=below:$b_2$] (b2) at (2.6,0) {};
      \node[pt,label=below:$b_3$] (b3) at (4.6,0) {};
      \draw (a1)--(b1)--(a2)--(b2)--(a3)--(b3);
    \end{tikzpicture}
  \end{minipage}\hfill
  \begin{minipage}[c]{0.52\textwidth}
    Maximal: $a_1,a_2,a_3$. Minimal: $b_1,b_2,b_3$. So there is no
    greatest and no least element \emph{(answer not on slides)}.
  \end{minipage}
\end{example}

\begin{example}[Ex.~2: $D_{12}=\{1,2,3,4,6,12\}$ under divisibility]\leavevmode\par\noindent
  \begin{minipage}[c]{0.3\textwidth}
    \centering
    \begin{tikzpicture}[hasse,y=0.9cm]
      \node[dv] (1) at (0,0) {1};
      \node[dv] (2) at (-0.8,1) {2};
      \node[dv] (3) at (0.8,1) {3};
      \node[dv] (4) at (-0.8,2) {4};
      \node[dv] (6) at (0.8,2) {6};
      \node[dv] (12) at (0,3) {12};
      \draw (1)--(2)--(4)--(12)--(6)--(3)--(1) (2)--(6);
    \end{tikzpicture}
  \end{minipage}\hfill
  \begin{minipage}[c]{0.66\textwidth}
    Greatest element $12$ (every element divides $12$); least element $1$
    ($1$ divides everything).
  \end{minipage}
\end{example}

\subsection{Upper and lower bounds, LUB and GLB}
\index{upper bound}\index{lower bound}\index{least upper bound (LUB)}\index{greatest lower bound (GLB)}\index{lower bound!of a one-element set}%
\begin{defn}
  Let $B\subseteq A$. An element $u\in A$ is an \textbf{upper bound} of
  $B$ if $b\le u$ for all $b\in B$; $\ell\in A$ is a \textbf{lower bound}
  of $B$ if $\ell\le b$ for all $b\in B$.
  The \textbf{least upper bound} $\mathrm{LUB}(B)$ (the
  \textbf{supremum}, $\sup B$) is the least of the upper bounds; the
  \textbf{greatest lower bound} $\mathrm{GLB}(B)$ (the
  \textbf{infimum}, $\inf B$) is the greatest of the lower bounds.
\end{defn}
Bounds are elements of $A$, not necessarily of $B$. In a Hasse
diagram: the upper bounds of $B$ are the elements reachable by an
upward path from \emph{every} element of $B$.

\begin{example}[Ex.~3: $B=\{1,2,3\}$ in $D_{12}$]
  Upper bounds: $6$ and $12$, so $\mathrm{LUB}(B)=6$. Lower bounds:
  just $1$, so $\mathrm{GLB}(B)=1$.
\end{example}

\begin{keybox}[title={Extremal elements need not exist}]
  A poset need not have a greatest or least element, and a subset need
  not have a LUB or GLB. $(\Z^+,\le)$ has least element $1$ but no
  greatest. A LUB can fail to exist in two ways: there are \textbf{no}
  upper bounds at all (Ex.~24 below), or there are upper bounds but
  \textbf{no least one} among them (Ex.~4 below: $d$ and $e$ are
  incomparable).
\end{keybox}

\begin{example}[Ex.~4]\leavevmode\par\noindent
  \begin{minipage}[t]{0.3\textwidth}
    \vspace{0pt}\centering
    \begin{tikzpicture}[hasse]
      \node[pt,label=below:$a$] (a) at (0,0) {};
      \node[pt,label=left:$b$]  (b) at (-0.8,1) {};
      \node[pt,label=right:$c$] (c) at (0.8,1) {};
      \node[pt,label=left:$d$]  (d) at (-1.4,2) {};
      \node[pt,label=right:$e$] (e) at (0,2) {};
      \node[pt,label=right:$f$] (f) at (1.6,1.6) {};
      \node[pt,label=above:$g$] (g) at (-1.4,3) {};
      \node[pt,label=above:$h$] (h) at (0,3) {};
      \draw (a)--(b)--(d)--(g) (a)--(c)--(e)--(h);
      \draw (b)--(e) (c)--(d) (c)--(f);
    \end{tikzpicture}
  \end{minipage}\hfill
  \begin{minipage}[t]{0.66\textwidth}
    \vspace{0pt}
    Greatest: none. Least: $a$. Maximal: $f,g,h$. Minimal: $a$.\\[2pt]
    For $B=\{b,c\}$: upper bounds $\{d,e,g,h\}$; \textbf{no LUB}, since
    the two lowest upper bounds $d$ and $e$ are incomparable
    ($d\not\le e$, $e\not\le d$). Lower bounds $\{a\}$, so
    $\mathrm{GLB}(B)=a$.
  \end{minipage}
\end{example}

\paragraph{Dual poset.} Reversing the order, $(A,\le)\to(A,\ge)$, swaps
every ``up'' concept with its ``down'' twin: greatest $\leftrightarrow$
least, maximal $\leftrightarrow$ minimal, upper bound $\leftrightarrow$
lower bound, LUB $\leftrightarrow$ GLB. Prove one, get the other free.

\subsection{Exam methods: greatest in a lexicographic order, bounds from a diagram, product order}
\index{greatest element!versus maximal}\index{lexicographic order}\index{product order}%
\textbf{Method for ``which is greatest'' in a lexicographic order:}
\begin{enumerate}[nosep]
  \item Find the largest first coordinate among the candidates. Drop
        every candidate with a smaller one --- they are all below.
  \item Among the survivors (all with the same first coordinate),
        compare the second coordinates in $B$'s order.
  \item The greatest is the one whose second coordinate is above
        \emph{all} the other survivors' second coordinates. If two
        survivors are incomparable and neither is beaten by a third,
        there is no greatest.
\end{enumerate}

\begin{example}[Greatest element, lexicographic order]
  $\mathbb{Z}$ with $\le$, $P(\{a,b\})$ with $\subseteq$, lexicographic
  order on the pairs. Which is greatest of $(1, \{a,b\})$, $(3, \{a\})$,
  $(3, \{b\})$, $(3, \{a,b\})$?
  \begin{enumerate}[nosep]
    \item Largest first coordinate is $3$: drop $(1, \{a,b\})$, despite
          its big set.
    \item Survivors' second coordinates: $\{a\}$, $\{b\}$, $\{a,b\}$.
    \item $\{a\} \subseteq \{a,b\}$ and $\{b\} \subseteq \{a,b\}$:
          $\{a,b\}$ is above both.
  \end{enumerate}
  Greatest: $(3, \{a,b\})$. \emph{Without} that candidate, $(3,\{a\})$ and
  $(3,\{b\})$ would both be maximal but incomparable --- no greatest.
\end{example}

\textbf{Method for ``lower bounds of $S$'' from a Hasse diagram:}
\index{lower bound!method}\index{lower bound!of a one-element set}\index{upper bound!method}\index{Hasse diagram!reading bounds}%
\begin{enumerate}[nosep]
  \item Start at each element of $S$ and walk \emph{downwards} along
        the lines, collecting everything you reach.
  \item \textbf{Include the element itself.} $x \le x$, so every element
        of $S$ is a bound of $S$ whenever it is below all the others.
        For a one-element set $S = \{b\}$, the lower bounds are $b$ and
        everything below $b$.
  \item For a bigger $S$, keep only what is below \emph{every} element
        of $S$. Upper bounds: the same, walking upwards.
\end{enumerate}
\textbf{Trap:} answering ``everything strictly below''. That drops $b$
itself --- the usual wrong MC option.

\begin{example}[Lower bounds of $\{b\}$]
  \leavevmode\par\noindent
  \begin{minipage}[c]{0.28\textwidth}
    \centering
    \begin{tikzpicture}[nhasse, scale=0.8]
      \node (a) at (0,2.4) {$a$};
      \node (c) at (-1.2,1.2) {$c$}; \node (b) at (1.2,1.2) {$b$};
      \node (d) at (-1.2,0) {$d$}; \node (f) at (0,0) {$f$}; \node (e) at (1.2,0) {$e$};
      \draw (a) -- (c) (a) -- (b) (c) -- (d) (b) -- (f) (b) -- (e);
    \end{tikzpicture}
  \end{minipage}\hfill
  \begin{minipage}[c]{0.7\textwidth}
    Below $b$: the lines go down to $e$ and $f$. Add $b$ itself:
    lower bounds of $\{b\}$ are \textbf{$b, e, f$}. Not $d$ ($d$ hangs
    under $c$, not $b$), and not $a$ or $c$ (above or beside $b$).
    \begin{itemize}[nosep]
      \item \textbf{GLB} of $\{b\}$ $= b$ (the greatest lower bound).
      \item Upper bounds of $\{b\}$: $a, b$; \textbf{LUB} $= b$.
    \end{itemize}
    $e, f$ alone is the wrong answer --- it forgets $b \le b$.
  \end{minipage}
\end{example}

\textbf{Product order} on pairs: $(x,y) \le (x',y') \iff x \le x'$
\emph{and} $y \le y'$. Picture points on a grid: ``below'' means down
and to the left. $(1,4)$ and $(3,2)$ are incomparable --- one coordinate
is larger, the other smaller. LUB of a set = (largest $x$, largest $y$),
GLB = (smallest $x$, smallest $y$), each taken separately.
\index{product order}%

\begin{example}[Product order on $\mathbb{Z}^+ \times \mathbb{Z}^+$]
  $S = \{(1,2), (2,1), (2,3), (4,1), (1,5)\}$.
  \begin{itemize}[nosep]
    \item \textbf{Minimal:} $(1,2)$ and $(2,1)$ --- nothing in $S$ is
          below either. Incomparable, so \textbf{no least element}.
    \item \textbf{Maximal:} $(2,3)$, $(4,1)$, $(1,5)$ --- nothing in $S$
          is above any of them. $(1,2) \le (2,3)$ and $(2,1) \le (2,3)$,
          so those two are not maximal. Three incomparable maximal
          elements, so \textbf{no greatest element}.
    \item \textbf{LUB} $= (\max x, \max y) = (4, 5)$, not in $S$ --- fine.
    \item \textbf{GLB} $= (\min x, \min y) = (1, 1)$.
  \end{itemize}
  \textbf{Trap:} $(0,0)$ is below everything too, but $0 \notin
    \mathbb{Z}^+$ --- it isn't in the poset, so it can't be a bound. Always
  check the ambient set. And if one element of $S$ is $\le$ all the
  others, it \emph{is} the least element and the GLB at once; a least
  element is always unique.
\end{example}

\subsection{Topological sorting: algorithm SORT}
\index{SORT algorithm}\index{topological sorting}%
In a finite poset you can always pick a minimal element (Theorem 1),
so this loop terminates with every element listed:
\begin{enumerate}
  \item Select a minimal element and append it to the ordering.
  \item Remove it from the poset.
  \item Repeat on what is left until nothing remains.
\end{enumerate}
It is the textbook ``Kahn's algorithm'' for build-order/dependency
resolution: repeatedly emit a node with no unmet prerequisites.

\begin{example}\leavevmode\par\noindent
  \begin{minipage}[c]{0.3\textwidth}
    \centering
    \begin{tikzpicture}[hasse]
      \node[pt,label=below:$d$] (d) at (-1,0) {};
      \node[pt,label=below:$e$] (e) at (1,0) {};
      \node[pt,label=right:$c$] (c) at (0,1) {};
      \node[pt,label=above:$a$] (a) at (-1,2) {};
      \node[pt,label=above:$b$] (b) at (1,2) {};
      \draw (d)--(c)--(e) (a)--(c)--(b);
    \end{tikzpicture}
  \end{minipage}\hfill
  \begin{minipage}[c]{0.66\textwidth}
    Minimal: $d,e$ --- take $d$. Then $e$ is minimal; take it. Then $c$.
    Then $a$ and $b$ are both minimal; take $b$, then $a$. One ordering:
    $d,e,c,b,a$. Each tie is a free choice, so the answer is not unique.
  \end{minipage}
\end{example}

\subsection{Exam method: which list is \emph{not} a topological sorting?}
\index{topological sorting}\index{Hasse diagram!topological sorting}%
\textbf{Method for ``which is \emph{not} a topological sorting'':}
\begin{enumerate}[nosep]
  \item From the Hasse diagram, write down the edges as ``lower $<$
        upper'' (each line, bottom end first).
  \item For each element with more than one line going \emph{up} from
        it, note everything above it --- these are the constraints
        most easily broken.
  \item In each option, check that every element appears after
        everything directly below it. The first violation rules the
        option out.
\end{enumerate}
Check which way the options run: if they all end with the top element,
the list goes bottom-up.

\begin{example}[Which is \emph{not} a topological sorting?]
  \leavevmode\par\noindent
  \begin{minipage}[c]{0.28\textwidth}
    \centering
    \begin{tikzpicture}[nhasse, scale=0.8]
      \node (d) at (-1.2,0) {$d$}; \node (e) at (0,0) {$e$}; \node (f) at (1.2,0) {$f$};
      \node (b) at (-1.2,1.2) {$b$}; \node (c) at (1.2,1.2) {$c$};
      \node (a) at (0,2.4) {$a$};
      \draw (d) -- (b) (e) -- (b) (e) -- (c) (f) -- (c) (b) -- (a) (c) -- (a);
    \end{tikzpicture}
  \end{minipage}\hfill
  \begin{minipage}[c]{0.7\textwidth}
    Constraints: $b$ after $d, e$; $c$ after $e, f$; $a$ last.
    $e$ sits under \emph{both} $b$ and $c$, so look for $e$ first.
    \begin{itemize}[nosep]
      \item $e, f, c, d, b, a$ --- fine
      \item $d, e, b, f, c, a$ --- fine ($b$ may come before $f$:
            they are incomparable)
      \item $d, f, b, c, e, a$ --- \textbf{not}: $e$ comes after
            $b$ and $c$
    \end{itemize}
  \end{minipage}
\end{example}

\subsection{Theorems}
\begin{rulebox}[title={6.2 Theorems 1--4}]
  \textbf{Th.~1.} A finite nonempty poset has at least one maximal and
  at least one minimal element.\\
  \textbf{Th.~2.} A poset has at most one greatest and at most one least
  element. \emph{Proof}: if $a,b$ are both greatest, then $a\le b$ and
  $b\le a$, so $a=b$ by antisymmetry.\\
  \textbf{Th.~3.} A subset $B$ has at most one LUB and at most one GLB.
  It says \emph{if} they exist they are unique --- not that they exist.\\
  \textbf{Th.~4.} An isomorphism $f:(A,\le)\to(A',\le')$
  ($x\le y\iff f(x)\le' f(y)$) carries maximal/minimal,
  greatest/least, upper/lower bounds and LUB/GLB to the same thing on
  the other side.
\end{rulebox}

\subsection{Exercises covered}
Figure numbers are Kolman's; the slides redraw the same diagrams.

\noindent\textbf{Ex.~1--4 (maximal/minimal)} and \textbf{Ex.~9--12
  (greatest/least)} \emph{(answers not on slides)}:\\[2pt]
\begin{tabular}{@{}lll@{\qquad}lll@{}}
  \toprule
  Ex.         & Maximal & Minimal & Ex.          & Greatest & Least                   \\
  \midrule
  1 (Fig.~27) & $3,5$   & $1,6$   & 9 (Fig.~31)  & $f$      & $a$                     \\
  2 (Fig.~28) & $f,g$   & $a,b,c$ & 10 (Fig.~32) & $e$      & none ($a$, $b$ minimal) \\
  3 (Fig.~29) & $e,f$   & $a$     & 11 (Fig.~33) & none     & none                    \\
  4 (Fig.~30) & $4,7$   & $1,8,9$ & 12 (Fig.~34) & $5$      & none ($1$, $3$ minimal) \\
  \bottomrule
\end{tabular}

\medskip
\noindent\textbf{Ex.~23--26 (bounds of $B$):}\\[2pt]
\begin{tabular}{@{}llllll@{}}
  \toprule
  Ex.          & $B$             & Upper bounds & Lower bounds & LUB  & GLB  \\
  \midrule
  23 (Fig.~35) & $\{c,d,e\}$     & $\{f,g,h\}$  & $\{a,b,c\}$  & $f$  & $c$  \\
  24 (Fig.~36) & $\{1,2,3,4,5\}$ & none         & none         & none & none \\
  25 (Fig.~37) & $\{b,c,d\}$     & $\{d,e,f\}$  & $\{a,b\}$    & $d$  & $b$  \\
  26 (Fig.~38) & $\{3,4,6\}$     & $\{5\}$      & $\{1,2,3\}$  & $5$  & $3$  \\
  \bottomrule
\end{tabular}

\medskip
\noindent\textbf{Ex.~29.} $A=\R$ with $\le$, $B=(1,2)$: upper bounds
$[2,\infty)$, lower bounds $(-\infty,1]$, $\mathrm{LUB}(B)=2$,
$\mathrm{GLB}(B)=1$ --- even though $1\notin B$ and $2\notin B$.

\noindent\textbf{Ex.~33.} SORT on Figure 35 (the poset of Ex.~23)
gives, for example, $a,b,c,d,e,f,g,h$. Another valid ordering is
$b,a,c,e,d,f,h,g$.

% =====================================================================
\kolman{3}
\section{Lattices}\src{Kolman 6.3}
\index{lattice}\index{join}\index{meet}\index{lattice!test from a Hasse diagram}%

\begin{defn}
  A \textbf{lattice} is a poset $(L,\le)$ in which every pair $a,b$ has
  both a LUB and a GLB, written
  \[
    a\vee b=\mathrm{LUB}(a,b)\ \text{(the \textbf{join})},\qquad
    a\wedge b=\mathrm{GLB}(a,b)\ \text{(the \textbf{meet})}.
  \]
\end{defn}
\textbf{In a Hasse diagram:} $a\vee b$ = go \emph{up} to the first
element both reach; $a\wedge b$ = go \emph{down} to the first element
both reach.
\notslides{The symbols are chosen on purpose: in $(\mathcal P(S),\subseteq)$
  join is $\cup$ and meet is $\cap$, and in logic they are OR and AND.
  That correspondence is where 6.4 (Boolean algebras) is heading.}

\begin{example}[Ex.~1: $(\mathcal P(S),\subseteq)$]
  For $A,B\subseteq S$: $A\cup B$ contains both and is the smallest set
  that does, so $A\vee B=A\cup B$; dually $A\wedge B=A\cap B$. Every
  pair has both, so $(\mathcal P(S),\subseteq)$ is a lattice. (The slide
  stops mid-sentence; this is the conclusion.)
\end{example}

\begin{example}[Ex.~2: $(\Z^+,\mid)$]
  An upper bound $u$ of $\{a,b\}$ needs $a\mid u$ and $b\mid u$ --- a
  common multiple --- so $a\vee b=\lcm(a,b)$. A lower bound $\ell$ needs
  $\ell\mid a$, $\ell\mid b$, so $a\wedge b=\gcd(a,b)$. Every pair has
  both, so it is a lattice. For $\{4,6\}$: upper bounds
  $12,24,36,\dots$, LUB $12$; lower bounds $1,2$, GLB $2$ \emph{(answer
    not on slides)}.
\end{example}

\begin{example}[Ex.~3--4: divisor lattices $D_{20}$ and $D_{30}$]\leavevmode\par\noindent
  \begin{minipage}[c]{0.24\textwidth}
    \centering
    \begin{tikzpicture}[hasse,y=0.9cm]
      \node[dv] (1) at (0,0) {1};
      \node[dv] (2) at (-0.8,1) {2};
      \node[dv] (5) at (0.8,1) {5};
      \node[dv] (4) at (-0.8,2) {4};
      \node[dv] (10) at (0.8,2) {10};
      \node[dv] (20) at (0,3) {20};
      \draw (1)--(2)--(4)--(20)--(10)--(5)--(1) (2)--(10);
    \end{tikzpicture}
  \end{minipage}
  \begin{minipage}[c]{0.24\textwidth}
    \small $2\vee5=10$\\ $4\wedge10=2$\\ $4\vee5=20$\\ $4\wedge5=1$
  \end{minipage}
  \begin{minipage}[c]{0.28\textwidth}
    \centering
    \begin{tikzpicture}[hasse,y=0.9cm]
      \node[dv] (1) at (0,0) {1};
      \node[dv] (2) at (-1.1,1) {2};
      \node[dv] (3) at (0,1) {3};
      \node[dv] (5) at (1.1,1) {5};
      \node[dv] (6) at (-1.1,2) {6};
      \node[dv] (10) at (0,2) {10};
      \node[dv] (15) at (1.1,2) {15};
      \node[dv] (30) at (0,3) {30};
      \draw (1)--(2) (1)--(3) (1)--(5);
      \draw (2)--(6) (2)--(10) (3)--(6) (3)--(15) (5)--(10) (5)--(15);
      \draw (6)--(30) (10)--(30) (15)--(30);
    \end{tikzpicture}
  \end{minipage}
  \begin{minipage}[c]{0.2\textwidth}
    \small $2\vee3=6$\\ $6\wedge10=2$\\ $6\vee15=30$\\ $10\wedge15=5$\\
    Lattice? \textbf{Yes} \emph{(answer not on slides)}
  \end{minipage}

  \smallskip
  Join is $\lcm$, meet is $\gcd$ --- read the diagram or just compute.
  Every $D_n$ is a lattice, since $\gcd$ and $\lcm$ of two divisors of
  $n$ are again divisors of $n$.
\end{example}

\textbf{Common lattices at a glance} (join, meet, bottom/top):
\begin{center}
  \renewcommand{\arraystretch}{1.3}
  \begin{tabular}{@{}lllll@{}}
    \toprule
    Lattice                         & Order            & $a \vee b$  & $a \wedge b$ & $0$ / $I$               \\
    \midrule
    $P(S)$                          & $\subseteq$      & $A \cup B$  & $A \cap B$   & $\emptyset$ / $S$       \\
    $D_n$ (divisors of $n$)         & $\mid$ (divides) & $\lcm(a,b)$ & $\gcd(a,b)$  & $1$ / $n$               \\
    $\mathbb{Z}$ (any linear order) & $\le$            & $\max(a,b)$ & $\min(a,b)$  & none                    \\
    $L_1 \times L_2$ product order  & $\le$ in both    & per coord.  & per coord.   & $(0,0)$ / $(I,I)$       \\
    $B_n$ (bit strings)             & bit by bit       & bitwise OR  & bitwise AND  & $0\ldots0$ / $1\ldots1$ \\
    \bottomrule
  \end{tabular}
\end{center}
\index{product order}\index{D n@$D_n$ (divisors of $n$)}%

\begin{example}[Product order on $\mathbb{Z} \times \mathbb{Z}$: $(x,y) \le (x',y') \iff x \le x'$ and $y \le y'$. Find $(3,-2) \vee (-1,4)$ and $(3,-2) \wedge (-1,4)$.]
  Join takes the larger value in \emph{each} coordinate, meet the smaller:
  \[
    (3,-2) \vee (-1,4) = (\max(3,-1), \max(-2,4)) = (3, 4), \qquad
    (3,-2) \wedge (-1,4) = (-1, -2).
  \]
  \textbf{Trap:} the answer is usually \emph{neither} of the two given
  pairs. Don't confuse this with the lexicographic order (first
  coordinate decides), and don't add the coordinates.
\end{example}

\begin{example}[In $D_{12}$: $4 \vee 6$ and $4 \wedge 6$]
  $4 \vee 6 = \lcm(4,6) = 12$ and $4 \wedge 6 = \gcd(4,6) = 2$.
  See the gcd/lcm methods in 1.4.5.
\end{example}

\subsection{Properties of lattices}
\index{lattice!properties}\index{absorption law!lattice}\index{idempotent law!lattice}\index{lattice!consistency ($a\le b\iff a\wedge b=a$)}\index{lattice!monotonicity of $\vee$, $\wedge$}%
\notslides{Kolman 6.3 pp.~267--268. These are the tools for proving
  anything about a lattice \emph{in general} (no Hasse diagram to look
  at).}

\noindent The definition of join and meet, written out as four facts:
\begin{rulepar}
  \begin{tabular}{@{}l@{\qquad}l@{}}
    (1) $a\le a\vee b$ and $b\le a\vee b$                   & ($a\vee b$ is an upper bound)       \\
    (2) if $a\le c$ and $b\le c$, then $a\vee b\le c$       & (it is the \emph{least} one)        \\
    (1$'$) $a\wedge b\le a$ and $a\wedge b\le b$            & ($a\wedge b$ is a lower bound)      \\
    (2$'$) if $c\le a$ and $c\le b$, then $c\le a\wedge b$  & (it is the \emph{greatest} one)
  \end{tabular}
\end{rulepar}

\begin{rulebox}[title={6.3 Theorem 2 (order vs.\ join and meet)}]
  For all $a,b$ in a lattice $L$:\quad
  (a) $a\vee b=b \iff a\le b$;\quad
  (b) $a\wedge b=a \iff a\le b$;\quad
  (c) $a\wedge b=a \iff a\vee b=b$.\\
  So ``$a\le b$'' can always be swapped for an equation, and back.
\end{rulebox}

\begin{rulebox}[title={6.3 Theorem 3 (algebraic laws)}]
  \begin{tabular}{@{}l@{\qquad}l@{\qquad}l@{}}
    Idempotent  & $a\vee a=a$                     & $a\wedge a=a$                     \\
    Commutative & $a\vee b=b\vee a$               & $a\wedge b=b\wedge a$             \\
    Associative & $a\vee(b\vee c)=(a\vee b)\vee c$ & $a\wedge(b\wedge c)=(a\wedge b)\wedge c$ \\
    Absorption  & $a\vee(a\wedge b)=a$            & $a\wedge(a\vee b)=a$
  \end{tabular}\\
  These hold in \emph{every} lattice. The distributive laws do
  \emph{not} --- see below.
\end{rulebox}

\begin{rulebox}[title={6.3 Theorem 4 (monotonicity)}]
  For all $a,b,c,d$ in a lattice $L$:
  \begin{enumerate}[nosep,label=\arabic*.]
    \item If $a\le b$, then (a) $a\vee c\le b\vee c$ and (b) $a\wedge c\le b\wedge c$.
    \item $a\le c$ and $b\le c$ $\iff$ $a\vee b\le c$.
    \item $c\le a$ and $c\le b$ $\iff$ $c\le a\wedge b$.
    \item If $a\le b$ and $c\le d$, then (a) $a\vee c\le b\vee d$ and (b) $a\wedge c\le b\wedge d$.
  \end{enumerate}
\end{rulebox}

\paragraph{Method: proving an equation $u=v$ in a general lattice.}
\index{lattice!proving an identity}\index{antisymmetric!proving equality}%
Either rewrite with the laws of Theorem 3 (and Theorem 2 to turn a given
$\le$ into an equation), or prove $u\le v$ and $v\le u$ separately and
conclude $u=v$ by \textbf{antisymmetry} --- each inequality from (1)--(2$'$)
and Theorem 4. This is how Kolman proves Theorem 3 itself:

\begin{example}[Absorption $a\vee(a\wedge b)=a$ (Kolman's proof of Theorem 3, 4(a))]
  $\le$: $a\wedge b\le a$ by (1$'$) and $a\le a$, so $a$ is an upper bound
  of $a$ and $a\wedge b$; by (2), $a\vee(a\wedge b)\le a$.\\
  $\ge$: $a\le a\vee(a\wedge b)$ by (1).\\
  Both directions, so $a\vee(a\wedge b)=a$ by antisymmetry.
\end{example}

\begin{example}[A two-line proof with Theorem 2: if $a\le b$, then $a\vee(a\wedge b)=a\wedge b$]
  $a\le b$ gives $a\wedge b=a$ (Theorem 2(b)). Then
  $a\vee(a\wedge b)=a\vee a=a=a\wedge b$ (idempotent).
\end{example}

\textbf{Watch out:} a law that holds in $D_n$ or $\mathcal P(S)$ need not
hold in every lattice. To prove a statement about \emph{all} lattices,
use only the definition and Theorems 2--4; to disprove one, a single
lattice (one Hasse diagram, one choice of elements) is enough.

\subsection{Exam method: is this Hasse diagram a lattice?}
\index{lattice!is it a lattice?}\index{comparable elements}\index{incomparable elements}%
\index{lattice!Hasse diagram test}%
\begin{enumerate}[nosep]
  \item \textbf{Exactly one top and one bottom?} Every finite lattice
        has a greatest element $I$ and a least element $0$ (Kolman 6.3,
        Theorem 5). Two maximal elements at the top: \textbf{not} a
        lattice --- those two have no upper bound at all.
  \item \textbf{Every incomparable pair has a \emph{unique} lowest
          common point above, and a unique highest common point below.}
        Comparable pairs (connected by an upward path) never fail.
        The failure to look for is the \textbf{bow-tie}: two elements
        that both sit below the same two incomparable elements.
\end{enumerate}

\begin{center}
  \begin{tikzpicture}[nhasse, scale=0.8]
    \node (a) at (0,0) {$a$};
    \node (b) at (-0.8,1.2) {$b$}; \node (c) at (0.8,1.2) {$c$};
    \draw (a) -- (b) (a) -- (c);
    \node[align=center] at (0,-0.9) {\small two tops: no $b \vee c$};
  \end{tikzpicture}\hspace{1.5cm}
  \begin{tikzpicture}[nhasse, scale=0.8]
    \node (0) at (0,0) {$0$};
    \node (a) at (-0.8,1) {$a$}; \node (b) at (0.8,1) {$b$};
    \node (c) at (-0.8,2.2) {$c$}; \node (d) at (0.8,2.2) {$d$};
    \node (I) at (0,3.2) {$I$};
    \draw (0) -- (a) (0) -- (b) (a) -- (c) (a) -- (d) (b) -- (c) (b) -- (d) (c) -- (I) (d) -- (I);
    \node[align=center] at (0,-0.9) {\small bow-tie: no $a \vee b$};
  \end{tikzpicture}\hspace{1.5cm}
  \begin{tikzpicture}[nhasse, scale=0.8]
    \node (1) at (0,0) {$1$};
    \node (2) at (-0.8,1) {$2$}; \node (3) at (0.8,1) {$3$};
    \node (4) at (-1.6,2) {$4$}; \node (6) at (0,2) {$6$};
    \node (12) at (-0.8,3) {$12$};
    \draw (1) -- (2) (1) -- (3) (2) -- (4) (2) -- (6) (3) -- (6) (4) -- (12) (6) -- (12);
    \node[align=center] at (-0.4,-0.9) {\small $D_{12}$: a lattice};
  \end{tikzpicture}
\end{center}
In the bow-tie, $a$ and $b$ have upper bounds $c$, $d$ and $I$; $c$ and
$d$ are both minimal among them and incomparable, so there is no
\emph{least} one. Note that it has a top and a bottom --- step 1 is
necessary, not sufficient.

\subsection{Sublattice}
\index{sublattice}%
A nonempty $S\subseteq L$ is a
\textbf{sublattice} if it is closed under the lattice's own $\vee$ and
$\wedge$: $a,b\in S\Rightarrow a\vee b\in S$ and $a\wedge b\in S$.
Example: $(D_n,\mid)$ is a sublattice of $(\Z^+,\mid)$.

\begin{rulebox}[title={6.3 Theorem 1 (product of lattices)}]
  If $(L_1,\le)$ and $(L_2,\le)$ are lattices, then $L_1\times L_2$ with
  the product order, $(a_1,a_2)\le(b_1,b_2)\iff a_1\le b_1$ and
  $a_2\le b_2$, is a lattice. (Join and meet are taken componentwise.)
\end{rulebox}

\textbf{Exam method: is $S$ a sublattice?}
\begin{itemize}[nosep]
  \item \textbf{Method:} for every incomparable pair in $S$, compute
        $a \vee b$ and $a \wedge b$ in $L$ (for $D_n$: lcm and gcd) and
        check that both are in $S$. Comparable pairs pass automatically.
  \item \textbf{Trap:} $S$ can be a lattice on its own and still not
        be a sublattice, because inside $S$ the ``join'' may be a
        higher element than the real join in $L$ (Kolman 6.3, Example 9).
\end{itemize}

\begin{example}[In $D_{36}$: are $\{1, 4, 9, 36\}$ and $\{2, 4, 6, 36\}$ sublattices?]
  \leavevmode
  \begin{itemize}[nosep]
    \item $\{1, 4, 9, 36\}$: the only incomparable pair is $4, 9$:
          $\lcm = 36 \in S$, $\gcd = 1 \in S$. \textbf{Sublattice.}
    \item $\{2, 4, 6, 36\}$: pair $4, 6$: $\gcd = 2 \in S$, but
          $\lcm = 12 \notin S$. \textbf{Not a sublattice} --- even
          though inside $S$ alone, $4$ and $6$ have the least upper
          bound $36$, so $S$ by itself is a lattice.
  \end{itemize}
\end{example}

\subsection{Bounded lattice}
\index{bounded lattice}%
$L$ is \textbf{bounded} if it has a least
element $0$ and a greatest element $I$, so $0\le a\le I$ for all $a$.
Then
\begin{rulepar}
  \[
    a\vee0=a,\qquad a\wedge0=0,\qquad a\vee I=I,\qquad a\wedge I=a.
  \]
\end{rulepar}
\notslides{Every finite lattice is bounded (Kolman 6.3 Theorem 5):
  $I=a_1\vee\dots\vee a_n$ and $0=a_1\wedge\dots\wedge a_n$.}

\subsection{Distributive lattice}
\index{distributive lattice}\index{non-distributive lattice}\index{pentagon lattice ($N_5$)}\index{diamond lattice ($M_3$)}%
$L$ is \textbf{distributive} if for
all $a,b,c$:
\[
  a\wedge(b\vee c)=(a\wedge b)\vee(a\wedge c)
  \quad\text{and}\quad
  a\vee(b\wedge c)=(a\vee b)\wedge(a\vee c).
\]
To show a lattice is \emph{not} distributive, \textbf{one
  counterexample is enough}: one triple $a,b,c$ where either identity
fails. $\mathcal P(S)$, every $D_n$ and every linear order are distributive.

\begin{example}[6.3 Ex.~17 and 18: distributive or not?]\leavevmode\par\noindent
  \begin{minipage}[t]{0.3\textwidth}
    \centering
    \begin{tikzpicture}[hasse]
      \node[pt,label=below:$0$] (0) at (0,0) {};
      \node[pt,label=left:$a$]  (a) at (-0.8,0.8) {};
      \node[pt,label=right:$c$] (c) at (0.8,0.8) {};
      \node[pt,label=left:$b$]  (b) at (-0.8,1.6) {};
      \node[pt,label=right:$d$] (d) at (0.8,1.6) {};
      \node[pt,label=above:$1$] (1) at (0,2.4) {};
      \draw (0)--(a)--(b)--(1)--(d)--(c)--(0) (a)--(d);
    \end{tikzpicture}

    \small\textbf{Ex.~17}: distributive
  \end{minipage}\hfill
  \begin{minipage}[t]{0.3\textwidth}
    \centering
    \begin{tikzpicture}[hasse]
      \node[pt,label=below:$0$] (0) at (0,0) {};
      \node[pt,label=left:$b$]  (b) at (-0.8,0.8) {};
      \node[pt,label=left:$a$]  (a) at (-0.8,1.6) {};
      \node[pt,label=right:$c$] (c) at (0.8,1.2) {};
      \node[pt,label=above:$1$] (1) at (0,2.4) {};
      \draw (0)--(b)--(a)--(1)--(c)--(0);
    \end{tikzpicture}

    \small\textbf{Ex.~18a}: pentagon $N_5$
  \end{minipage}\hfill
  \begin{minipage}[t]{0.3\textwidth}
    \centering
    \begin{tikzpicture}[hasse]
      \node[pt,label=below:$0$] (0) at (0,0) {};
      \node[pt,label=left:$a$]  (a) at (-0.9,1.2) {};
      \node[pt,label=right:$b$] (b) at (0,1.2) {};
      \node[pt,label=right:$c$] (c) at (0.9,1.2) {};
      \node[pt,label=above:$1$] (1) at (0,2.4) {};
      \draw (0)--(a)--(1)--(b)--(0)--(c)--(1);
    \end{tikzpicture}

    \small\textbf{Ex.~18b}: diamond $M_3$
  \end{minipage}

  \medskip
  \textbf{Ex.~17} satisfies both distributive laws for every triple. (Not
  on the slides: it is the $3\times2$ grid, the product of a 3-chain and a
  2-chain, and a product of chains is always distributive.)\\
  \textbf{Ex.~18a}: $a\wedge(b\vee c)=a\wedge1=a$, but
  $(a\wedge b)\vee(a\wedge c)=b\vee0=b\ne a$. Not distributive.\\
  \textbf{Ex.~18b}: $a\wedge(b\vee c)=a\wedge1=a$, but
  $(a\wedge b)\vee(a\wedge c)=0\vee0=0\ne a$. Not distributive.
\end{example}

\begin{rulebox}[title={Not on the slides: where to find the counterexample (Kolman 6.3 Theorem 6)}]
  A lattice is nondistributive \textbf{if and only if} it contains a
  sublattice shaped like one of the two in Ex.~18: the \textbf{pentagon}
  $N_5$ or the \textbf{diamond} $M_3$. So scan the Hasse diagram for an
  element with three upper covers meeting again above (diamond), or a
  square where one side has an extra element (pentagon). Then use its
  labels as the triple: the element on the long side of the pentagon, or
  any middle element of the diamond, as $a$ in $a\wedge(b\vee c)$.
\end{rulebox}

\subsection{Modular lattice}
\index{modular lattice}\index{lattice!modular}%
\notslides{Not in the lectures or the text proper: Kolman 6.3
  Exercises 26 and 39 (p.~272).}
\begin{defn}
  A lattice $(L,\le)$ is \textbf{modular} if for all $a,b,x\in L$:
  \[
    a\le b \implies a\vee(x\wedge b)=(a\vee x)\wedge b.
  \]
\end{defn}
\textbf{Reading it.} Only pairs with $a\le b$ are tested, and for them
the two ways of ``combining $a$, $x$ and $b$'' must agree: on the left,
cut $x$ down below $b$ first and then join with $a$; on the right, join
$x$ with $a$ first and then cut down below $b$. Set it side by side with
the second distributive law $a\vee(b\wedge c)=(a\vee b)\wedge(a\vee c)$:
the shapes are close, and the modular law is the weaker demand of the
two. (Kolman's Exercise 26 writes the same law with $a\le c$ and $b$ in
the middle.)

\begin{rulepar}
  \textbf{What Kolman asks (p.~272), no proofs given:}
  Ex.~26(a) every \textbf{distributive} lattice is modular;
  Ex.~26(b) the diamond in Figure 52 ($M_3$) is \textbf{nondistributive
    but modular};
  Ex.~39 prove or disprove: a sublattice of a modular lattice is modular.
\end{rulepar}

\paragraph{Method: checking the modular law.}
\index{modular lattice!checking}%
\begin{itemize}[nosep]
  \item \textbf{To disprove:} one triple $a\le b$, $x$ where the two sides
        differ is enough. Compute both sides from the Hasse diagram (go up
        for $\vee$, down for $\wedge$).
  \item \textbf{To prove for a given finite lattice:} the law must hold
        for \emph{every} triple with $a\le b$. List the pairs $a\le b$,
        and for each one the possible $x$; use Theorems 2--4 to dispose
        of whole families of cases at once rather than computing every
        triple by hand. State clearly why each family is covered.
  \item \textbf{To prove for a whole class} (e.g.\ all distributive
        lattices): work with letters only --- the given laws plus
        Theorems 2--4, never a picture.
\end{itemize}

\begin{rulebox}[title={Useful inequality (proved here, so it can be cited by this proof)}]
  In any lattice, if $a\le b$ then
  \[
    a\vee(x\wedge b)\;\le\;(a\vee x)\wedge b .
  \]
  \emph{Proof.} $a\le a\vee x$ by (1) and $a\le b$, so $a\le(a\vee x)\wedge b$
  by (2$'$). Also $x\wedge b\le x\le a\vee x$ and $x\wedge b\le b$, so
  $x\wedge b\le(a\vee x)\wedge b$ by (2$'$). Both $a$ and $x\wedge b$ are
  below $(a\vee x)\wedge b$, so their join is too, by (2). \qed\\
  Consequence: to show the modular law for a triple you only need the
  \emph{other} direction, ``$\ge$''.
\end{rulebox}

\begin{example}[The mechanics in $D_{12}$ (a distributive lattice)]
  Take $a=2\le b=12$ and $x=3$. Left: $2\vee(3\wedge12)=\lcm(2,\gcd(3,12))
  =\lcm(2,3)=6$. Right: $(2\vee3)\wedge12=\gcd(\lcm(2,3),12)=\gcd(6,12)=6$.
  Equal. With $a=2\le b=4$, $x=6$: left $2\vee\gcd(6,4)=2\vee2=2$; right
  $\gcd(\lcm(2,6),4)=\gcd(6,4)=2$. Equal again.
\end{example}

\subsection{Complement}
\index{complement!in a lattice}\index{complemented lattice}%
In a bounded lattice, $a'$ is a
\textbf{complement} of $a$ if
\[
  a\vee a'=I\quad\text{(together they reach the top)}\qquad\text{and}\qquad
  a\wedge a'=0\quad\text{(they share only the bottom)}.
\]
\notslides{$L$ is \textbf{complemented} if it is bounded and
  \emph{every} element has at least one complement (Kolman 6.3) ---
  needed for Ex.~28--29 below. In a nondistributive lattice complements
  can be non-unique: in $M_3$, $a$ has both $b$ and $c$ as complements.}

\begin{example}[Ex.~19: $L=\mathcal P(\{x,y\})$]
  Here $0=\emptyset$, $I=\{x,y\}$, $\vee=\cup$, $\wedge=\cap$. Since
  $\{x\}\cup\{y\}=I$ and $\{x\}\cap\{y\}=\emptyset=0$, the complement of
  $\{x\}$ is $\{y\}$. All of them:
  $\emptyset'=\{x,y\}$, $\{x\}'=\{y\}$, $\{y\}'=\{x\}$,
  $\{x,y\}'=\emptyset$ --- exactly the ordinary set complement.
\end{example}

\begin{rulebox}[title={6.3 Theorem 7}]
  In a bounded distributive lattice, a complement, if it exists, is
  unique --- so we can write it $a'$ unambiguously.
\end{rulebox}
\erratum{The slide says only ``a distributive lattice''. Complements are
  only defined in a \emph{bounded} lattice, and Kolman's Theorem 7 says
  ``bounded distributive''.}

\textbf{Complements, in practice.} Always $0'=I$ and $I'=0$; in a lattice with two or more elements nothing is its own complement. An element can have \textbf{no}, \textbf{one} or \textbf{several} complements; in a bounded \emph{distributive} lattice at most one.

\begin{center}
  \begin{tikzpicture}[nhasse, scale=0.8]
    \node (0) at (0,0) {$0$};
    \node (a) at (-1,1) {$a$}; \node (b) at (0,1) {$b$}; \node (c) at (1,1) {$c$};
    \node (I) at (0,2) {$I$};
    \draw (0) -- (a) (0) -- (b) (0) -- (c) (a) -- (I) (b) -- (I) (c) -- (I);
    \node[align=center] at (0,-0.9) {\small diamond $M_3$};
  \end{tikzpicture}\hspace{2cm}
  \begin{tikzpicture}[nhasse, scale=0.8]
    \node (0) at (0,0) {$0$};
    \node (a) at (-0.9,0.8) {$a$}; \node (b) at (-0.9,1.6) {$b$};
    \node (c) at (0.9,1.2) {$c$};
    \node (I) at (0,2.4) {$I$};
    \draw (0) -- (a) (a) -- (b) (b) -- (I) (0) -- (c) (c) -- (I);
    \node[align=center] at (0,-0.9) {\small pentagon $N_5$};
  \end{tikzpicture}
\end{center}

\begin{example}[Complements in the pentagon $N_5$ above]
  $c'$: try every element. $a \vee c = I$ and $a \wedge c = 0$; the same
  for $b$. So $c$ has \textbf{two} complements, $a$ and $b$; and
  $a' = c$, $b' = c$. Two complements for one element prove that $N_5$
  is not distributive (Theorem 7). In $M_3$, each of $a, b, c$ has the
  other two as complements.
\end{example}

\textbf{Complements in $D_n$:} the only candidate for $d'$ is $n/d$
(since $\gcd = 1$ and $\lcm = n$ force $d \cdot d' = n$). It works
exactly when $\gcd(d, n/d) = 1$.

\begin{example}[Complements of $5$ and $2$ in $D_{30}$ and in $D_{20}$]
  $D_{30}$: $5' = 30/5 = 6$, and $\gcd(5, 6) = 1$ --- yes.
  $D_{20}$: $5' = 4$, $\gcd(5, 4) = 1$ --- yes; but for $2$ the
  candidate is $10$ and $\gcd(2, 10) = 2 \ne 1$ --- $2$ has
  \textbf{no complement}, so $D_{20}$ is not complemented.
\end{example}

\subsection{Exercises covered}

\noindent\textbf{Ex.~1 (lattice?)} \emph{(answers not on slides)}.
Diagram 1 (Fig.~45): \textbf{yes} --- it is a $3\times3$ grid, the
product of two 3-element chains, so a lattice by Theorem 1.
Diagram 2: \textbf{no} --- $e,f$ are both maximal, so $e\vee f$ does
not exist (and $a,b$ minimal, so no $a\wedge b$). Diagram 3:
\textbf{no} --- $a,b$ are both minimal, so $a\wedge b$ does not exist.

\medskip
\noindent\textbf{Ex.~7.} $A=\{2,3,6,12,24,36,72\}$ under $\mid$:
\textbf{not a lattice}. The pair $2,3$ has no lower bound in $A$ ($1$
is missing), so $2\wedge3$ does not exist. (Every other pair is fine:
$24\vee36=72$, $24\wedge36=12$.)
\erratum{The solution slide's bullets say ``every pair has a unique
  LUB'' and ``every pair has a unique GLB'', then conclude ``not a
  lattice'' --- contradictory. The first bullet is true; the second is
  false, because of $2\wedge3$, and that is the reason it fails.}

\medskip
\noindent\textbf{Ex.~28--29 (distributive? complemented?)} Both are
\textbf{neither}. The slides work Ex.~28, with the same triple $c,f,g$
and the same complement argument for $c$; Ex.~29 is \emph{(answer not
  on slides)}.
\begin{itemize}
  \item \textbf{28} (Fig.~53; $a<b<c<e$, $a<d$, $d<c,f,g<e$): $d$ has
        three upper covers $c,f,g$ meeting at $e$ --- a diamond
        $\{d,c,f,g,e\}$. $c\wedge(f\vee g)=c\wedge e=c$ but
        $(c\wedge f)\vee(c\wedge g)=d\vee d=d$: not distributive. $c$ has no
        complement (every element is either below $c$ or above $d$, so
        $c\vee x=c$ or $c\wedge x\ge d$): not complemented.
  \item \textbf{29} (Fig.~54; $a<b<c<d<f$, $a<e<d$): $\{a,b,c,d,e\}$
        is a pentagon. $c\wedge(b\vee e)=c\wedge d=c$ but
        $(c\wedge b)\vee(c\wedge e)=b\vee a=b$: not distributive. Only $f$
        joins with anything to give $I=f$, and $b\wedge f=b\ne a$, so $b$
        has no complement: not complemented.
\end{itemize}

\begin{keybox}[title=Key points (6.2--6.3)]
  \textbf{Greatest/least} (above/below everything, unique if it exists)
  vs.\ \textbf{maximal/minimal} (nothing strictly above/below, can be
  several). \textbf{Upper/lower bounds} of $B$ live in $A$;
  \textbf{LUB}/\textbf{GLB} are unique if they exist, but may not.
  \textbf{SORT}: repeatedly remove a minimal element. A
  \textbf{lattice} has $a\vee b$ (join, LUB) and $a\wedge b$ (meet,
  GLB) for every pair; $\mathcal P(S)$: $\cup,\cap$; divisibility:
  $\lcm,\gcd$. \textbf{Sublattice}: closed under $\vee,\wedge$.
  \textbf{Bounded}: has $0$ and $I$. \textbf{Distributive}: one
  counterexample disproves it. \textbf{Complement}: $a\vee a'=I$,
  $a\wedge a'=0$; unique in a bounded distributive lattice.
\end{keybox}

% =====================================================================
\kolman{4}
\section{Finite Boolean algebras}\src{Kolman 6.4}
\index{Boolean algebra}\index{B n@$B_n$}\index{atom of a Boolean algebra}%

Three ideas from 6.3 --- \textbf{lattice}, \textbf{complement},
\textbf{distributive} --- combine into one structure, the
\textbf{Boolean algebra}.

\begin{rulebox}[title={6.4 Theorem 1}]
  If $S_1=\{x_1,\dots,x_n\}$ and $S_2=\{y_1,\dots,y_n\}$ have the same
  number of elements, then $(\mathcal P(S_1),\subseteq)$ and
  $(\mathcal P(S_2),\subseteq)$ are isomorphic lattices.
\end{rulebox}
\textbf{Why:} pair the elements up, $x_i\leftrightarrow y_i$. That
induces a map on subsets, e.g.\ $f(\{x_1,x_3\})=\{y_1,y_3\}$, and it
preserves inclusion: $A\subseteq B\iff f(A)\subseteq f(B)$. Only the
\emph{number} of elements matters, not what they are --- so call the
lattice of subsets of any $n$-element set $B_n$.

\begin{example}[6.4 Ex.~1--2: three faces of $B_3$]\leavevmode\par\noindent
  \newcommand{\cube}[8]{%
    \begin{tikzpicture}[hasse,boxed,x=0.75cm,y=0.75cm]
      \node[pt,label=below:{#1}] (0) at (0,0) {};
      \node[pt,label=left:{#2}]  (x) at (-1.2,1) {};
      \node[pt,label=right:{#3}] (y) at (0,1) {};
      \node[pt,label=right:{#4}] (z) at (1.2,1) {};
      \node[pt,label=left:{#5}]  (xy) at (-1.2,2) {};
      \node[pt,label=right:{#6}] (xz) at (0,2) {};
      \node[pt,label=right:{#7}] (yz) at (1.2,2) {};
      \node[pt,label=above:{#8}] (1) at (0,3) {};
      \draw (0)--(x) (0)--(y) (0)--(z) (x)--(xy) (x)--(xz) (y)--(xy) (y)--(yz)
      (z)--(xz) (z)--(yz) (xy)--(1) (xz)--(1) (yz)--(1);
    \end{tikzpicture}}%
  \begin{minipage}[t]{0.33\textwidth}\centering
    \cube{$\emptyset$}{$\{a\}$}{$\{b\}$}{$\{c\}$}{$\{a,b\}$}{$\{a,c\}$}{$\{b,c\}$}{$\{a,b,c\}$}\\
    \small (a) $\mathcal P(\{a,b,c\})$
  \end{minipage}%
  \begin{minipage}[t]{0.33\textwidth}\centering
    \cube{$\emptyset$}{$\{2\}$}{$\{3\}$}{$\{5\}$}{$\{2,3\}$}{$\{2,5\}$}{$\{3,5\}$}{$\{2,3,5\}$}\\
    \small (b) $\mathcal P(\{2,3,5\})$
  \end{minipage}%
  \begin{minipage}[t]{0.33\textwidth}\centering
    \cube{$000$}{$100$}{$010$}{$001$}{$110$}{$101$}{$011$}{$111$}\\
    \small (c) bit strings of length 3
  \end{minipage}

  \smallskip
  Same Hasse diagram, different labels. In (c) bit $k$ says whether
  element $k$ is in the subset, and $\le$ is ``bit by bit''.
\end{example}

\begin{defn}
  A finite lattice is a \textbf{Boolean algebra} if it is isomorphic to
  $B_n$ for some $n\ge0$. Since $|B_n|=2^n$, \textbf{a finite Boolean
    algebra always has $2^n$ elements} --- a quick test to rule one out.
\end{defn}

\begin{center}
  \begin{tikzpicture}[hasse,boxed,x=0.7cm,y=0.7cm]
    \node[pt,label=below:$\emptyset$] at (0,0) {};
    \node[font=\small] at (0,-1.1) {$n=0$};
    \begin{scope}[xshift=2.2cm]
      \node[pt,label=below:$0$] (0) at (0,0) {};
      \node[pt,label=above:$1$] (1) at (0,1.5) {};
      \draw (0)--(1);
      \node[font=\small] at (0,-1.1) {$n=1$};
    \end{scope}
    \begin{scope}[xshift=4.8cm]
      \node[pt,label=below:$00$] (0) at (0,0) {};
      \node[pt,label=left:$10$]  (a) at (-0.8,0.8) {};
      \node[pt,label=right:$01$] (b) at (0.8,0.8) {};
      \node[pt,label=above:$11$] (1) at (0,1.6) {};
      \draw (0)--(a)--(1)--(b)--(0);
      \node[font=\small] at (0,-1.1) {$n=2$};
    \end{scope}
    \begin{scope}[xshift=8.4cm]
      \node[pt,label=below:$000$] (0) at (0,0) {};
      \node[pt,label=left:$100$]  (x) at (-1.1,0.8) {};
      \node[pt,label=right:$010$] (y) at (0,0.8) {};
      \node[pt,label=right:$001$] (z) at (1.1,0.8) {};
      \node[pt,label=left:$110$]  (xy) at (-1.1,1.6) {};
      \node[pt,label=right:$101$] (xz) at (0,1.6) {};
      \node[pt,label=right:$011$] (yz) at (1.1,1.6) {};
      \node[pt,label=above:$111$] (1) at (0,2.4) {};
      \draw (0)--(x) (0)--(y) (0)--(z) (x)--(xy) (x)--(xz) (y)--(xy) (y)--(yz)
      (z)--(xz) (z)--(yz) (xy)--(1) (xz)--(1) (yz)--(1);
      \node[font=\small] at (0,-1.1) {$n=3$};
    \end{scope}
  \end{tikzpicture}
\end{center}

\begin{rulepar}
  \textbf{Properties:} a Boolean algebra is \textbf{bounded},
  \textbf{complemented} and \textbf{distributive}. So every element has a
  complement (unique, by 6.3 Theorem 7) and both distributive laws hold.
\end{rulepar}

\subsection{Exam method: is it a Boolean algebra?}
\index{Boolean algebra!not a Boolean algebra}\index{isomorphism}\index{Boolean algebra}\index{B n@$B_n$}%
$B_n$ is the set of bit strings of length $n$, ordered bit by bit
($x \le y$ when every bit of $x$ is $\le$ the same bit of $y$). Then
$x \vee y$ is \textbf{bitwise OR}, $x \wedge y$ is \textbf{bitwise AND},
and the complement $x'$ is \textbf{bitwise NOT}. This is exactly $P(S)$
with $|S| = n$, each subset stored as a bit mask (bit $k$ is $1$ when
element $k$ is in the set) --- a \texttt{[Flags]} enum: $\cup$ is
\texttt{|}, $\cap$ is \texttt{\&}, complement is \texttt{\~{}}.

\begin{center}
  \begin{tikzpicture}[nhasse, scale=0.85]
    \node (000) at (0,0) {$000$};
    \node (100) at (-1.4,1) {$100$}; \node (010) at (0,1) {$010$}; \node (001) at (1.4,1) {$001$};
    \node (110) at (-1.4,2) {$110$}; \node (101) at (0,2) {$101$}; \node (011) at (1.4,2) {$011$};
    \node (111) at (0,3) {$111$};
    \draw (000) -- (100) (000) -- (010) (000) -- (001)
    (100) -- (110) (100) -- (101) (010) -- (110) (010) -- (011) (001) -- (101) (001) -- (011)
    (110) -- (111) (101) -- (111) (011) -- (111);
    \node at (0,-0.7) {\small $B_3$};
  \end{tikzpicture}\hspace{1.5cm}
  \begin{tikzpicture}[nhasse, scale=0.85]
    \node (1) at (0,0) {$1$};
    \node (2) at (-1.4,1) {$2$}; \node (3) at (0,1) {$3$}; \node (5) at (1.4,1) {$5$};
    \node (6) at (-1.4,2) {$6$}; \node (10) at (0,2) {$10$}; \node (15) at (1.4,2) {$15$};
    \node (30) at (0,3) {$30$};
    \draw (1) -- (2) (1) -- (3) (1) -- (5)
    (2) -- (6) (2) -- (10) (3) -- (6) (3) -- (15) (5) -- (10) (5) -- (15)
    (6) -- (30) (10) -- (30) (15) -- (30);
    \node at (0,-0.7) {\small $D_{30}$, same shape};
  \end{tikzpicture}
\end{center}
The isomorphism $D_{30} \to B_3$: one bit per prime $(2, 3, 5)$ ---
$1 \mapsto 000$, $2 \mapsto 100$, $6 \mapsto 110$, $30 \mapsto 111$, and so on.

\textbf{Is $L$ a Boolean algebra? Ways to say no, quickest first:}
\begin{enumerate}[nosep]
  \item $|L|$ is not a power of $2$ ($1, 2, 4, 8, 16, \ldots$). No
        Boolean algebra has $3$, $5$, $6$ or $7$ elements.
  \item Some element has no complement, or two complements.
  \item It contains a diamond $M_3$ or pentagon $N_5$ (not distributive).
\end{enumerate}
To say \textbf{yes}: match the diagram to $B_n$ (for $8$ elements, the cube above).

\subsection{Which \texorpdfstring{$D_n$}{D n} are Boolean algebras?}
\index{D n@$D_n$ (divisors of $n$)}\index{Boolean algebra!D n test@$D_n$ as a Boolean algebra}\index{square-free}\index{distinct primes}%

\begin{example}[6.4 Ex.~3: $D_{20}$ and $D_{30}$]
  $D_{20}=\{1,2,4,5,10,20\}$ has $6$ elements, and $6\ne2^n$ for any $n$:
  \textbf{not} a Boolean algebra ($20=2^2\cdot5$).\\
  $D_{30}=\{1,2,3,5,6,10,15,30\}$ has $8=2^3$ elements --- the right
  count, so it \emph{may} be one. It is: $D_{30}\cong B_3$, and is a
  Boolean algebra.
  \notslides{The count alone doesn't prove it --- Ex.~6 below has $8$
    elements and fails. What settles $D_{30}$ is Theorem 2:
    $30=2\cdot3\cdot5$.}
\end{example}

\begin{rulebox}[title={6.4 Theorem 2}]
  If $n=p_1p_2\cdots p_k$ with $p_1,\dots,p_k$ \textbf{distinct} primes,
  then $D_n$ is a Boolean algebra, isomorphic to $B_k$.
\end{rulebox}
\notslides{The converse holds too (Kolman 6.4 Example 8): if $p^2\mid n$
  for some prime $p$, then $p$ has no complement in $D_n$, so $D_n$ is
  \emph{not} a Boolean algebra. Together: $D_n$ is a Boolean algebra
  $\iff$ no prime divides $n$ twice. E.g.\ $D_{40}$ ($40=2^3\cdot5$) and
  $D_{75}$ ($75=3\cdot5^2$) are not.}

\index{Boolean algebra!D n test@$D_n$ test}\index{square-free}%
\begin{rulepar}
  \textbf{Rule:} $D_n$ is a Boolean algebra $\iff$ $n$ is
  \textbf{square-free} --- a product of \emph{distinct} primes, no prime
  appearing twice in the factorization. (Kolman 6.4, Theorem 2 and Example 8.)
\end{rulepar}

\textbf{Method:} prime-factorize $n$ (trial division, see 1.4.2) and look for any exponent $\ge 2$. Watch for $4$, $9$, $25$,
$49$ hiding inside.

\textbf{Why:} each divisor is a choice of which primes to include, so
divisors $\leftrightarrow$ subsets of the primes, lcm $= \cup$ and
gcd $= \cap$. With $p^2 \mid n$, the divisor $p$ has no complement:
the only candidate $n/p$ still contains a factor $p$, so
$\gcd(p, n/p) \ne 1$.

\begin{example}[Which of $D_{78}$, $D_{98}$, $D_{91}$, $D_{60}$, $D_{17}$ are Boolean algebras?]
  \leavevmode
  \begin{itemize}[nosep]
    \item $78 = 2 \cdot 3 \cdot 13$ --- distinct primes, \textbf{yes} ($\cong B_3$, $8$ divisors)
    \item $98 = 2 \cdot 7^2$ --- \textbf{no}
    \item $91 = 7 \cdot 13$ --- \textbf{yes} ($\cong B_2$). $91$ looks prime but isn't.
    \item $60 = 2^2 \cdot 3 \cdot 5$ --- \textbf{no} ($12$ divisors, not a power of $2$)
    \item $17$ is prime: $D_{17} = \{1, 17\} \cong B_1$ --- \textbf{yes}.
          A prime is square-free; don't skip it.
  \end{itemize}
\end{example}

\begin{example}[6.4 Ex.~5: $D_{210}$, $D_{66}$, $D_{646}$]
  $210=2\cdot3\cdot5\cdot7$, $66=2\cdot3\cdot11$, $646=2\cdot17\cdot19$ ---
  all distinct primes, so $D_{210}\cong B_4$, $D_{66}\cong B_3$,
  $D_{646}\cong B_3$: all three are Boolean algebras.
\end{example}

\subsection{Shape of the Hasse diagram of $D_n$ (exam method)}
\index{D n@$D_n$ (divisors of $n$)}\index{Hasse diagram!of D n@of $D_n$}\index{product of chains}%
Factorize $n=p_1^{e_1}\cdots p_k^{e_k}$. Each prime is one ``direction'': its
exponent runs $0,1,\dots,e_i$, a chain of $e_i+1$ elements, and $D_n$ is the
\emph{product} of these chains (join $=\lcm$, meet $=\gcd$, as for any product
order). So
\begin{itemize}[nosep]
  \item \textbf{dots} $=$ number of divisors $=(e_1+1)\cdots(e_k+1)$;
  \item \textbf{levels} $=e_1+\dots+e_k+1$ (the height from $1$ up to $n$);
  \item one prime $\Rightarrow$ a single vertical \textbf{chain}.
\end{itemize}

\begin{center}
  \renewcommand{\arraystretch}{1.25}
  \begin{tabular}{@{}lll@{}}
    \toprule
    $n$ looks like & shape of $D_n$ & example \\
    \midrule
    $p$           & chain of 2 dots                    & $D_{7}=\{1,7\}$ \\
    $p^2$         & chain of 3 dots                    & $D_{25}=\{1,5,25\}$, $D_{9}$, $D_{49}$ \\
    $p^3$         & chain of 4 dots                    & $D_{8}$, $D_{27}$ \\
    $pq$          & diamond (square, 4 dots)           & $D_{6}$, $D_{15}$, $D_{91}$ \\
    $p^2q$        & $3\times2$ grid, 6 dots            & $D_{12}$, $D_{20}$, $D_{18}$ \\
    $p^2q^2$      & $3\times3$ grid, 9 dots            & $D_{36}$ \\
    $pqr$         & cube $B_3$, 8 dots                 & $D_{30}$, $D_{42}$, $D_{66}$ \\
    \bottomrule
  \end{tabular}
\end{center}
A prime power gives a chain, so $D_{25}$ is \emph{not} a diamond: $\{1,5,25\}$
has $1\mid5\mid25$ and nothing incomparable. \textbf{To match an MC picture:}
count the dots first (that is $d(n)$), then check chain, grid or cube. Only
the cube-shaped $8$-dot diagram is Boolean ($n$ square-free), while an
$8$-dot chain (e.g.\ $D_{128}$) or the $4\times2$ grid ($D_{24}$) is not.

\begin{example}[MC practice: $(\mathcal P(\{1,2,3\}),\subseteq)$ vs.\ $D_{30}$, $D_9$; which statement is false?]
  \index{power set!isomorphic to D n@isomorphic to $D_n$}\index{isomorphism!P(S) and D n@$\mathcal P(S)$ and $D_n$}\index{GCD and LUB in a power set}%
  $\mathcal P(S)$ with $|S|=k$ is $B_k$, and $D_n\cong B_k$ exactly when
  $n$ is a product of $k$ \emph{distinct} primes. So with $|S|=3$:
  $D_{30}$ ($2\cdot3\cdot5$), $D_{42}$, $D_{66}$ all match, with $2^3=8$
  elements. $D_9=\{1,3,9\}$ ($9=3^2$) is a 3-element \emph{chain}: not
  isomorphic, and the false statement.\\
  In any lattice the words change but the operations do not:
  \textbf{LUB $=$ join $=\cup$}, \textbf{GLB $=$ meet $=$ ``gcd'' $=\cap$} in
  a power set (gcd is only the name of the meet in $D_n$). So
  $\mathrm{LUB}(\{1\},\{2,3\})=\{1,2,3\}$ and
  $\mathrm{GLB}(\{1\},\{2,3\})=\emptyset$ are both true, and
  $|\mathcal P(S)|=2^3=8$.
\end{example}

\subsection{Fast square-free check (no factoring needed)}
\index{square-free!fast check}\index{D n@$D_n$ (divisors of $n$)}\index{Boolean algebra!D n test@$D_n$ test}%
To decide Boolean or not you do \emph{not} need the factorization, only
whether some prime square divides $n$. Test the squares of primes with
$p^2\le n$, each by a quick rule:

\begin{center}
  \renewcommand{\arraystretch}{1.25}
  \begin{tabular}{@{}cll@{}}
    \toprule
    square & divides $n$ when \ldots & only relevant when \\
    \midrule
    $4$   & last two digits are a multiple of $4$ & always \\
    $9$   & digit sum is a multiple of $9$ & $n\ge9$ \\
    $25$  & ends in $00$, $25$, $50$ or $75$ & $n\ge25$ \\
    $49$  & $n$ is a multiple of $49$ ($49,98,147,196,\dots$) & $n\ge49$ \\
    $121$ & multiples $121,242,363,\dots$ & $n\ge121$ \\
    $169$ & multiples $169,338,\dots$ & $n\ge169$ \\
    \bottomrule
  \end{tabular}
\end{center}
If none of them divides $n$, then $n$ is square-free and $D_n$ \emph{is} a
Boolean algebra. For $n<49$ only $4$, $9$ and $25$ can matter.

\begin{rulebox}[title={Not square-free, $n\le100$ (so $D_n$ is not Boolean)}]
  Every multiple of $4$: $4,8,12,\dots,96,100$.\\
  Plus the odd multiples of $9$: $9,27,45,63,81,99$,
  and the even ones not already above: $18,54,90$.\\
  Plus $25,50,75$ and $49,98$.\\
  Every other $n\le100$ is square-free, so $D_n$ \emph{is} Boolean.
\end{rulebox}

\begin{example}[$D_{37},D_{38},D_{39},D_{40},D_{41},D_{42}$]
  Only $4$, $9$, $25$ can divide these. $37$, $41$: odd, digit sums $10$
  and $5$, no. $38$: last two digits $38$, not a multiple of $4$; digit sum
  $11$. $39$: digit sum $12$, not $9$. $42$: $42$ is not a multiple of
  $4$, digit sum $6$. But $40=4\cdot10$, so $D_{40}$ is the one that is
  \textbf{not} Boolean ($40=2^3\cdot5$).
\end{example}

\subsection{Lookup: every $D_n$ seen in questions and examples}
\index{D n@$D_n$ (divisors of $n$)!lookup table}\index{Boolean algebra!D n test@$D_n$ test}%
Boolean $\iff$ $n$ square-free; dots $=$ number of divisors; shape from the
exponents (see the two subsections before).
\begin{center}
  \small\renewcommand{\arraystretch}{1.15}
  \begin{tabular}{@{}llrll@{}}
    \toprule
    $D_n$ & $n=$ & dots & Boolean? & shape \\
    \midrule
    $D_{9}$ & $3^2$ & 3 & \textbf{no} & chain of 3 \\
    $D_{12}$ & $2^2 \cdot 3$ & 6 & \textbf{no} & $3\times2$ grid \\
    $D_{17}$ & $17$ & 2 & \textbf{yes} ($B_1$) & chain of 2 \\
    $D_{20}$ & $2^2 \cdot 5$ & 6 & \textbf{no} & $3\times2$ grid \\
    $D_{25}$ & $5^2$ & 3 & \textbf{no} & chain of 3 \\
    $D_{30}$ & $2 \cdot 3 \cdot 5$ & 8 & \textbf{yes} ($B_3$) & cube $B_3$ \\
    $D_{36}$ & $2^2 \cdot 3^2$ & 9 & \textbf{no} & $3\times3$ grid \\
    $D_{37}$ & $37$ & 2 & \textbf{yes} ($B_1$) & chain of 2 \\
    $D_{38}$ & $2 \cdot 19$ & 4 & \textbf{yes} ($B_2$) & diamond \\
    $D_{39}$ & $3 \cdot 13$ & 4 & \textbf{yes} ($B_2$) & diamond \\
    $D_{40}$ & $2^3 \cdot 5$ & 8 & \textbf{no} & $4\times2$ grid \\
    $D_{41}$ & $41$ & 2 & \textbf{yes} ($B_1$) & chain of 2 \\
    $D_{42}$ & $2 \cdot 3 \cdot 7$ & 8 & \textbf{yes} ($B_3$) & cube $B_3$ \\
    $D_{60}$ & $2^2 \cdot 3 \cdot 5$ & 12 & \textbf{no} & $3\times2\times2$ product of chains \\
    $D_{63}$ & $3^2 \cdot 7$ & 6 & \textbf{no} & $3\times2$ grid \\
    $D_{66}$ & $2 \cdot 3 \cdot 11$ & 8 & \textbf{yes} ($B_3$) & cube $B_3$ \\
    $D_{75}$ & $3 \cdot 5^2$ & 6 & \textbf{no} & $3\times2$ grid \\
    $D_{78}$ & $2 \cdot 3 \cdot 13$ & 8 & \textbf{yes} ($B_3$) & cube $B_3$ \\
    $D_{84}$ & $2^2 \cdot 3 \cdot 7$ & 12 & \textbf{no} & $3\times2\times2$ product of chains \\
    $D_{90}$ & $2 \cdot 3^2 \cdot 5$ & 12 & \textbf{no} & $3\times2\times2$ product of chains \\
    $D_{91}$ & $7 \cdot 13$ & 4 & \textbf{yes} ($B_2$) & diamond \\
    $D_{92}$ & $2^2 \cdot 23$ & 6 & \textbf{no} & $3\times2$ grid \\
    $D_{98}$ & $2 \cdot 7^2$ & 6 & \textbf{no} & $3\times2$ grid \\
    $D_{105}$ & $3 \cdot 5 \cdot 7$ & 8 & \textbf{yes} ($B_3$) & cube $B_3$ \\
    $D_{210}$ & $2 \cdot 3 \cdot 5 \cdot 7$ & 16 & \textbf{yes} ($B_4$) & $B_4$ \\
    $D_{646}$ & $2 \cdot 17 \cdot 19$ & 8 & \textbf{yes} ($B_3$) & cube $B_3$ \\
    \bottomrule
  \end{tabular}
\end{center}

\subsection{Calculating in a Boolean algebra}
\index{Boolean algebra!laws}\index{De Morgan's laws!Boolean algebra}\index{absorption law}\index{involution}\index{set algebra}%

\begin{rulebox}[title={6.4 Theorem 4: Boolean algebra $\leftrightarrow$ set algebra}]
  \centering
  \begin{tabular}{@{}lccccc@{}}
    Boolean algebra & $a\vee b$ & $a\wedge b$ & $a'$          & $0$         & $I$ \\
    Set algebra     & $A\cup B$ & $A\cap B$   & $\overline A$ & $\emptyset$ & $S$
  \end{tabular}

  \smallskip
  \raggedright
  In a Boolean algebra $\vee$, $\wedge$, $'$ behave exactly like $\cup$,
  $\cap$, complement. So any set identity carries over --- useful when
  simplifying Boolean expressions.
\end{rulebox}

\begin{example}[6.4 Ex.~6]
  For $x,y\in L$, a Boolean algebra:\quad
  \textbf{involution} $(x')'=x$;\quad
  \textbf{De Morgan} $(x\wedge y)'=x'\vee y'$ and $(x\vee y)'=x'\wedge y'$.\\
  Proof: the matching set laws $\overline{\overline A}=A$,
  $\overline{A\cap B}=\overline A\cup\overline B$,
  $\overline{A\cup B}=\overline A\cap\overline B$ hold, and Theorem 4
  translates them.
\end{example}

\noindent\textbf{Boolean properties} (from the slides; $x,y,z\in L$, and
the primed column is the set version in $\mathcal P(S)$):\\[3pt]
\begin{rulepar}
  {\small
    \begin{tabular}{@{}r@{\ }l@{\qquad}r@{\ }l@{}}
      \toprule
      1.  & $x\le y\iff x\vee y=y$                        & 1$'$.  & $A\subseteq B\iff A\cup B=B$                                             \\
      2.  & $x\le y\iff x\wedge y=x$                      & 2$'$.  & $A\subseteq B\iff A\cap B=A$                                             \\
      3.  & $x\vee x=x$,\ \ $x\wedge x=x$                 & 3$'$.  & $A\cup A=A$,\ \ $A\cap A=A$                                              \\
      4.  & $x\vee y=y\vee x$,\ \ $x\wedge y=y\wedge x$   & 4$'$.  & $A\cup B=B\cup A$,\ \ $A\cap B=B\cap A$                                  \\
      5.  & $x\vee(y\vee z)=(x\vee y)\vee z$              & 5$'$.  & $A\cup(B\cup C)=(A\cup B)\cup C$                                         \\
          & $x\wedge(y\wedge z)=(x\wedge y)\wedge z$      &        & $A\cap(B\cap C)=(A\cap B)\cap C$                                         \\
      6.  & $x\vee(x\wedge y)=x$,\ \ $x\wedge(x\vee y)=x$ & 6$'$.  & $A\cup(A\cap B)=A$,\ \ $A\cap(A\cup B)=A$                                \\
      7.  & $0\le x\le I$ for all $x$                     & 7$'$.  & $\emptyset\subseteq A\subseteq S$                                        \\
      8.  & $x\vee0=x$,\ \ $x\wedge0=0$                   & 8$'$.  & $A\cup\emptyset=A$,\ \ $A\cap\emptyset=\emptyset$                        \\
      9.  & $x\vee I=I$,\ \ $x\wedge I=x$                 & 9$'$.  & $A\cup S=S$,\ \ $A\cap S=A$                                              \\
      10. & $x\wedge(y\vee z)=(x\wedge y)\vee(x\wedge z)$ & 10$'$. & $A\cap(B\cup C)=(A\cap B)\cup(A\cap C)$                                  \\
          & $x\vee(y\wedge z)=(x\vee y)\wedge(x\vee z)$   &        & $A\cup(B\cap C)=(A\cup B)\cap(A\cup C)$                                  \\
      11. & unique $x'$: $x\vee x'=I$, $x\wedge x'=0$     & 11$'$. & unique $\overline A$: $A\cup\overline A=S$, $A\cap\overline A=\emptyset$ \\
      \bottomrule
    \end{tabular}}

  \medskip
  \noindent 3 idempotent, 4 commutative, 5 associative, 6 absorption,
  10 distributive, 11 complement --- plus involution and De Morgan from
  Ex.~6.
\end{rulepar}

\index{De Morgan's laws}\index{duality}%
\begin{rulepar}
  \textbf{Substitution rule} (Kolman 6.4, Theorem 3): any identity for sets using $\cup$, $\cap$ and complement holds in every Boolean algebra with $\vee$ for $\cup$ and $\wedge$ for $\cap$. They are also the laws of propositional logic and of \texttt{||}, \texttt{\&\&}, \texttt{!} on \texttt{bool} ($=B_1$).

  Every law comes in a pair: swap $\vee \leftrightarrow \wedge$ and
  $0 \leftrightarrow I$ to get the other one (\emph{duality}).
  Note that \emph{both} distributive laws hold --- unlike ordinary
  arithmetic, where $a + bc \ne (a+b)(a+c)$.
\end{rulepar}

\begin{rulepar}
  \paragraph{Absorption laws (exam method).}\index{absorption law}
  \[
    x \vee (x \wedge y) = x \qquad\qquad x \wedge (x \vee y) = x
  \]
  \textbf{Recognise it:} the same variable stands \emph{alone} outside the
  brackets and again \emph{inside}, with the \emph{other} operator in the
  brackets. The whole bracket vanishes --- $y$ can be anything, even a
  long expression.
\end{rulepar}

\textbf{Why it works} --- three ways to see it:
\begin{itemize}[nosep]
  \item \emph{Sets:} $A \cap B \subseteq A$, so adding it to $A$ adds
        nothing: $A \cup (A \cap B) = A$. Dually, $A \subseteq A \cup B$,
        so $A \cap (A \cup B) = A$.
  \item \emph{Order:} $x \wedge y \le x$, and the order law
        ``$a \le b \iff a \vee b = b$'' gives $x \vee (x \wedge y) = x$.
  \item \emph{Code:} \texttt{x || (x \&\& y)}: if \texttt{x} is true the
        result is true; if \texttt{x} is false, \texttt{x \&\& y} is false
        too, so the result is false. Either way it equals \texttt{x}.
\end{itemize}

\textbf{Check by truth table} (the proof that always works):
\begin{center}
  \begin{tabular}{cc|c|c|c|c}
    $x$ & $y$ & $x \wedge y$ & $x \vee (x \wedge y)$ & $x \vee y$ & $x \wedge (x \vee y)$ \\ \hline
    0   & 0   & 0            & 0                     & 0          & 0                     \\
    0   & 1   & 0            & 0                     & 1          & 0                     \\
    1   & 0   & 0            & 1                     & 1          & 1                     \\
    1   & 1   & 1            & 1                     & 1          & 1
  \end{tabular}
\end{center}
Both result columns equal the $x$ column.

\begin{example}[Absorption with bigger pieces]
  \begin{itemize}[nosep]
    \item $(x \wedge y) \vee (x \wedge y \wedge z) = x \wedge y$ ---
          take $x \wedge y$ as ``the variable'': it appears alone, and
          again inside $(x \wedge y) \wedge z$.
    \item $z' \wedge (z' \vee x \vee y) = z'$ --- the alone part is
          $z'$; the bracket is $z' \vee$ (anything).
    \item $(x \vee y) \wedge \bigl(x \vee y \vee (w \wedge z)\bigr) = x \vee y$.
  \end{itemize}
\end{example}

\textbf{Trap --- the look-alike with a complement:} $x \vee (x' \wedge y)$
has $x'$ inside, not $x$, so absorption does \emph{not} apply. It
simplifies to $x \vee y$, not $x$ --- see the next example. Similarly
$x \wedge (x' \vee y) = x \wedge y$.

\begin{example}[Simplify $x \vee (x' \wedge y)$]
  \[
    x \vee (x' \wedge y)
    \overset{\text{distr.}}{=} (x \vee x') \wedge (x \vee y)
    \overset{\text{compl.}}{=} I \wedge (x \vee y)
    \overset{\text{bounds}}{=} x \vee y.
  \]
  Name the law at each step, as in a proof. The usual trick is the
  distributive law that turns a complement pair $x \vee x'$ into $I$.
\end{example}

\index{atom}\textbf{Atoms:} the elements directly above $0$ (in $B_n$
the strings with exactly one $1$; in $D_n$ the primes). $B_n$ has $n$
atoms, and every element is the join of the atoms below it ---
like every subset being the union of its one-element subsets.

\subsection{Exercises covered}

\noindent\textbf{6.4 Ex.~1--6: is the poset a Boolean algebra?}\\[2pt]
\begin{minipage}[b]{0.16\textwidth}\centering
  \begin{tikzpicture}[hasse,boxed,x=0.6cm,y=0.6cm]
    \node[pt,label=below:$a$] (a) at (0,0) {};
    \node[pt,label=left:$b$]  (b) at (-0.8,0.9) {};
    \node[pt,label=right:$c$] (c) at (0.8,0.9) {};
    \node[pt,label=left:$d$]  (d) at (-0.8,1.9) {};
    \node[pt,label=right:$e$] (e) at (0.8,1.9) {};
    \node[pt,label=above:$f$] (f) at (0,2.8) {};
    \draw (a)--(b)--(d)--(f)--(e)--(c)--(a) (b)--(e);
  \end{tikzpicture}\\ \small 1.
\end{minipage}%
\begin{minipage}[b]{0.16\textwidth}\centering
  \begin{tikzpicture}[hasse,boxed,x=0.6cm,y=0.6cm]
    \node[pt,label=below:$a$] (a) at (0,0) {};
    \node[pt,label=left:$b$]  (b) at (-0.7,0.7) {};
    \node[pt,label=right:$c$] (c) at (0.7,0.7) {};
    \node[pt,label=right:$d$] (d) at (0,1.4) {};
    \node[pt,label=left:$e$]  (e) at (-0.7,2.1) {};
    \node[pt,label=right:$f$] (f) at (0.7,2.1) {};
    \node[pt,label=above:$g$] (g) at (0,2.8) {};
    \draw (a)--(b)--(d)--(c)--(a) (d)--(e)--(g)--(f)--(d);
  \end{tikzpicture}\\ \small 2.
\end{minipage}%
\begin{minipage}[b]{0.2\textwidth}\centering
  \begin{tikzpicture}[hasse,boxed,x=0.6cm,y=0.6cm]
    \node[pt,label=below:$a$] (a) at (0,0) {};
    \node[pt,label=left:$b$]  (b) at (-0.8,0.8) {};
    \node[pt,label=left:$c$]  (c) at (-1.6,1.6) {};
    \node[pt,label=left:$d$]  (d) at (0.8,0.8) {};
    \node[pt,label=right:$e$] (e) at (1.6,1.6) {};
    \node[pt,label=right:$f$] (f) at (1.6,0.1) {};
    \draw (c)--(b)--(a)--(d)--(e) (d)--(f);
  \end{tikzpicture}\\ \small 3.
\end{minipage}%
\begin{minipage}[b]{0.1\textwidth}\centering
  \begin{tikzpicture}[hasse,boxed,x=0.6cm,y=0.6cm]
    \node[pt,label=right:$a$] (a) at (0,0) {};
    \node[pt,label=right:$c$] (c) at (0,0.7) {};
    \node[pt,label=right:$b$] (b) at (0,1.4) {};
    \node[pt,label=right:$d$] (d) at (0,2.1) {};
    \node[pt,label=right:$e$] (e) at (0,2.8) {};
    \draw (a)--(c)--(b)--(d)--(e);
  \end{tikzpicture}\\ \small 4.
\end{minipage}%
\begin{minipage}[b]{0.2\textwidth}\centering
  \begin{tikzpicture}[hasse,boxed,x=0.6cm,y=0.6cm]
    \node[pt,label=below:$a$] (a) at (0,0) {};
    \node[pt,label=left:$b$]  (b) at (-1.1,0.9) {};
    \node[pt,label=right:$d$] (d) at (0,0.9) {};
    \node[pt,label=right:$c$] (c) at (1.1,0.9) {};
    \node[pt,label=left:$f$]  (f) at (-1.1,1.9) {};
    \node[pt,label=right:$e$] (e) at (0,1.9) {};
    \node[pt,label=right:$g$] (g) at (1.1,1.9) {};
    \node[pt,label=above:$h$] (h) at (0,2.8) {};
    \draw (a)--(b) (a)--(d) (a)--(c) (b)--(f) (b)--(e) (d)--(f) (d)--(g)
    (c)--(e) (c)--(g) (f)--(h) (e)--(h) (g)--(h);
  \end{tikzpicture}\\ \small 5.
\end{minipage}%
\begin{minipage}[b]{0.18\textwidth}\centering
  \begin{tikzpicture}[hasse,boxed,x=0.6cm,y=0.6cm]
    \node[pt,label=below:$a$] (a) at (0,0) {};
    \node[pt,label=left:$b$]  (b) at (-0.9,0.8) {};
    \node[pt,label=right:$c$] (c) at (0,0.8) {};
    \node[pt,label=right:$d$] (d) at (0.9,0.8) {};
    \node[pt,label=left:$e$]  (e) at (0,1.6) {};
    \node[pt,label=left:$f$]  (f) at (-0.8,2.4) {};
    \node[pt,label=right:$g$] (g) at (0.8,2.4) {};
    \node[pt,label=above:$h$] (h) at (0,3.2) {};
    \draw (a)--(b)--(e) (a)--(c)--(e) (a)--(d)--(e) (e)--(f)--(h)--(g)--(e);
  \end{tikzpicture}\\ \small 6.
\end{minipage}

\medskip
\begin{enumerate}[label=\arabic*.,nosep]
  \item \textbf{No} --- $6$ elements, $6\ne2^n$.
  \item \textbf{No} --- $7$ elements, $7\ne2^n$.
  \item \textbf{No} --- not even a lattice: $c$ and $e$ have no common
        upper bound, so $c\vee e$ does not exist.
  \item \textbf{No} --- a lattice (every chain is), but $5$ elements,
        $5\ne2^n$.
  \item \textbf{Yes} --- $8=2^3$ elements, and the diagram has the same
        structure as $B_3$, so $L\cong B_3$.
  \item \textbf{No} --- also $8$ elements, which shows that having $2^n$
        elements is not enough on its own.
\end{enumerate}
\erratum{For Ex.~6 the slide argues that $b$ has two complements,
  $c$ and $d$, ``since $b\wedge c=a$, $b\vee c=e$''. But the top of this
  lattice is $h$, not $e$, so $b\vee c=e$ makes $c$ \emph{no} complement
  of $b$. The right argument: $b$ has \textbf{no complement at all} ---
  every element is either $\le e$ (then $b\vee x\le e\ne h$) or $\ge e$
  (then $b\wedge x=b\ne a$). So the lattice isn't complemented and can't be
  a Boolean algebra. The slide's calculation \emph{is} valid inside the
  lower diamond $\{a,b,c,d,e\}$, where $e$ is the top: there $b$ has two
  complements, $c$ and $d$. That diamond is $M_3$, so the lattice isn't
  distributive either --- a second valid reason.}

% =====================================================================
\kolman{5}
\section{Functions on Boolean algebras}\src{Kolman 6.5}
\index{Boolean function}\index{truth table!of a Boolean function}\index{Boolean polynomial}%

\begin{defn}
  With $B=\{0,1\}$, a \textbf{Boolean function} is a function
  $f:B^n\to B$: the input is a bit string of length $n$, the output a
  single bit. E.g.\ $f:B^3\to B$ takes $(x_1,x_2,x_3)$ and returns $0$ or
  $1$.
\end{defn}
A Boolean function is completely described by its \textbf{truth table},
one row per input, so $2^n$ rows. Many different polynomials give the same function; simplifying with the laws means fewer gates in the circuit.
\notslides{The slide breaks off at ``There are''. Each of the $2^n$
  rows can be $0$ or $1$, so there are $2^{2^n}$ Boolean functions of $n$
  variables --- $16$ for $n=2$, $256$ for $n=3$.}

\begin{defn}
  \textbf{Boolean polynomials} (Boolean expressions) in $x_1,\dots,x_n$
  are built by these rules:
  (1) each variable $x_i$ is one;
  (2) the constants $0$ and $1$ are;
  (3) if $p$ and $q$ are, so are $p\vee q$ and $p\wedge q$;
  (4) if $p$ is, so is $p'$.
\end{defn}
Example: from $x,y,z$ build $x\wedge y$, then $(x\wedge y)'$, then
$(x\wedge y)'\vee z$ --- each step uses one rule on polynomials already
built.

\begin{rulebox}[title={Simplification rules}]
  \index{Boolean polynomial!simplification rules}%
  \begin{tabular}{@{}l@{\qquad}l@{}}
    $x\vee0=x$,\quad $x\wedge1=x$   & $x\vee1=1$,\quad $x\wedge0=0$                           \\
    $x\vee x=x$,\quad $x\wedge x=x$ & $x\vee x'=1$,\quad $x\wedge x'=0$                       \\
    \textbf{Absorption}:            & $x\vee(x\wedge y)=x$,\quad $x\wedge(x\vee y)=x$         \\
    \textbf{De Morgan}:             & $(x\vee y)'=x'\wedge y'$,\quad $(x\wedge y)'=x'\vee y'$
  \end{tabular}
\end{rulebox}

\textbf{Truth table method:} $2^n$ rows, inputs listed by counting in
binary ($000, 001, \ldots, 111$); one helper column per subexpression,
innermost first; the last column is $f$.

\begin{example}[Truth table of $p(x, y, z) = (x' \vee y) \wedge z$]
  \begin{center}
    \begin{tabular}{ccc|cc|c}
      $x$ & $y$ & $z$ & $x'$ & $x' \vee y$ & $f$ \\ \hline
      0   & 0   & 0   & 1    & 1           & 0   \\
      0   & 0   & 1   & 1    & 1           & 1   \\
      0   & 1   & 0   & 1    & 1           & 0   \\
      0   & 1   & 1   & 1    & 1           & 1   \\
      1   & 0   & 0   & 0    & 0           & 0   \\
      1   & 0   & 1   & 0    & 0           & 0   \\
      1   & 1   & 0   & 0    & 1           & 0   \\
      1   & 1   & 1   & 0    & 1           & 1
    \end{tabular}
  \end{center}
  Shortcut for checking: $\wedge z$ means every row with $z = 0$ gives
  $0$ --- fill those first, then work out only the $z = 1$ rows.
\end{example}

\begin{example}[6.5 Ex.~1]\leavevmode\par\noindent
  \begin{minipage}[t]{0.3\textwidth}
    \textbf{1a.} $p_1=(x\vee y)\wedge z$:\\[2pt]
    \small
    \begin{tabular}{@{}ccc|c@{}}
      $x$ & $y$ & $z$ & $p_1$ \\\hline
      0   & 0   & 0   & 0     \\ 0&0&1&0\\ 0&1&0&0\\ 0&1&1&1\\
      1   & 0   & 0   & 0     \\ 1&0&1&1\\ 1&1&0&0\\ 1&1&1&1
    \end{tabular}
  \end{minipage}\hfill
  \begin{minipage}[t]{0.66\textwidth}
    \textbf{1b.} $p_2=(x\vee y')\vee(y\wedge1)$. Since $y\wedge1=y$,
    $p_2=x\vee y'\vee y=x\vee1=\mathbf1$ --- always true.

    \medskip
    \textbf{1c.} $p_3=x\vee(y'\wedge z)\vee(x\wedge y)
      =\bigl(x\vee(x\wedge y)\bigr)\vee(y'\wedge z)=\mathbf{x\vee(y'\wedge z)}$
    by absorption.

    \medskip
    \textbf{1d.} $p=(x_1\wedge x_2)\vee\bigl(x_1\vee(x_2'\wedge x_3)\bigr)$:
    build it up one column at a time (table below).
  \end{minipage}

  \medskip
  \centering\small
  \begin{tabular}{@{}ccc|cccc@{}}
    $x_1$ & $x_2$ & $x_3$ & $x_1\wedge x_2$ & $x_2'\wedge x_3$ & $x_1\vee(x_2'\wedge x_3)$ & $p$ \\\hline
    0     & 0     & 0     & 0               & 0                & 0                         & 0   \\ 0&0&1&0&1&1&1\\ 0&1&0&0&0&0&0\\ 0&1&1&0&0&0&0\\
    1     & 0     & 0     & 0               & 0                & 1                         & 1   \\ 1&0&1&0&1&1&1\\ 1&1&0&1&0&1&1\\ 1&1&1&1&0&1&1
  \end{tabular}

  \raggedright
  \notslides{The last two columns are equal: $x_1\wedge x_2$ is absorbed
    by the $x_1$ inside the second term, so $p=x_1\vee(x_2'\wedge x_3)$.}
\end{example}

\subsection{Logic diagrams}
\index{logic diagram}\index{AND gate}\index{OR gate}\index{NOT gate}\index{NAND gate}\index{NOR gate}%
Each operation is a \textbf{gate}: $x\vee y$ is an \textbf{OR} gate (1
if at least one input is 1), $x\wedge y$ an \textbf{AND} gate (1 only if
both are), $x'$ a \textbf{NOT} gate (flips $0\leftrightarrow1$). Two more
are common: \textbf{NAND}$(A,B)=(A\wedge B)'$ and
\textbf{NOR}$(A,B)=(A\vee B)'$ --- an AND or OR with a NOT bubble on the
output.

\begin{center}
  \begin{tikzpicture}[gates]
    \foreach \g/\lbl [count=\i] in {or gate/OR, and gate/AND, nand gate/NAND, nor gate/NOR}{
        \node[\g,inputs=nn] (G\i) at (2.4*\i,0) {};
        \draw (G\i.input 1) -- ++(-0.35,0);
        \draw (G\i.input 2) -- ++(-0.35,0);
        \draw (G\i.output) -- ++(0.35,0);
        \node at (2.4*\i,-0.85) {\lbl};
      }
    \node[not gate] (N) at (12,0) {};
    \draw (N.input) -- ++(-0.35,0);
    \draw (N.output) -- ++(0.35,0);
    \node at (12,-0.85) {NOT};
  \end{tikzpicture}
\end{center}

\begin{example}[Truth table and circuit of $p(x,y,z)=(x\wedge y)\vee(y\wedge z')$]\leavevmode\par\noindent
  \begin{minipage}[c]{0.3\textwidth}
    \small
    \begin{tabular}{@{}ccc|c@{}}
      $x$ & $y$ & $z$ & $p$ \\\hline
      0   & 0   & 0   & 0   \\ 0&0&1&0\\ 0&1&0&1\\ 0&1&1&0\\
      1   & 0   & 0   & 0   \\ 1&0&1&0\\ 1&1&0&1\\ 1&1&1&1
    \end{tabular}
  \end{minipage}\hfill
  \begin{minipage}[c]{0.66\textwidth}
    \centering
    \begin{tikzpicture}[gates]
      \coordinate (L) at (0,0);
      \coordinate (Y) at (0.6,0);
      \node[and gate,inputs=nn] (A1) at (2.6,0.8) {};
      \node[and gate,inputs=nn] (A2) at (2.6,-0.8) {};
      \node[not gate,anchor=output] (N) at ([xshift=-0.25cm]A2.input 2) {};
      \node[or gate,inputs=nn] (O) at (4.6,0) {};
      \draw (A1.input 1) -- (A1.input 1 -| L) node[left] {$x$};
      \draw (A1.input 2) -- (A1.input 2 -| L) node[left] {$y$};
      \fill (A1.input 2 -| Y) circle (1.5pt);
      \draw (A1.input 2 -| Y) |- (A2.input 1);
      \draw (N.input) -- (N.input -| L) node[left] {$z$};
      \draw (N.output) -- (A2.input 2);
      \draw (A1.output) -- ++(0.3,0) |- (O.input 1);
      \draw (A2.output) -- ++(0.3,0) |- (O.input 2);
      \draw (O.output) -- ++(0.4,0) node[right] {$p$};
      \node[above] at (A1.north) {$x\wedge y$};
      \node[below] at (A2.south) {$y\wedge z'$};
    \end{tikzpicture}
  \end{minipage}

  \smallskip
  Reading the diagram back: the two AND gates give $x\wedge y$ and
  $y\wedge z'$, and the OR gate joins them, which gives back
  $(x\wedge y)\vee(y\wedge z')$.
\end{example}

\subsection{Exercises covered}

\noindent\textbf{6.5 Ex.~3.} $p(x,y,z)=(x\wedge y')\vee\bigl(y\wedge(x'\vee y)\bigr)$.
Simplify first:
\[
  p=(x\wedge y')\vee\bigl((y\wedge x')\vee(y\wedge y)\bigr)
  =(x\wedge y')\vee\bigl((x'\wedge y)\vee y\bigr)
  =(x\wedge y')\vee y
  =(x\vee y)\wedge(y'\vee y)
  =(x\vee y)\wedge1
  =x\vee y.
\]
So $f(x,y,z)=x\vee y$: the truth table is $0$ exactly on the rows with
$x=y=0$ ($000$ and $001$) and $1$ on the other six. $z$ doesn't matter.
\notslides{A shortcut: $y\wedge(x'\vee y)=y$ straight away by
  absorption.}

\medskip
\noindent\begin{minipage}[c]{0.42\textwidth}
  \textbf{6.5 Ex.~18} (Kolman Fig.~75). Read the circuit left to right:
  \begin{enumerate}[nosep]
    \item AND: $y\wedge z$
    \item OR: $x\vee(y\wedge z)$
    \item NOT: $\bigl(x\vee(y\wedge z)\bigr)'$
    \item lower branch NOT: $z'$
    \item final OR:
  \end{enumerate}
  \[
    f(x,y,z)=\bigl(x\vee(y\wedge z)\bigr)'\vee z'.
  \]
\end{minipage}\hfill
\begin{minipage}[c]{0.56\textwidth}
  \centering
  \begin{tikzpicture}[gates]
    \coordinate (L) at (0,0);
    \coordinate (B) at (0.45,0);
    \node[and gate,inputs=nn] (A) at (1.4,-0.15) {};
    \node[or gate,inputs=nn] (O1) at (3.0,0.55) {};
    \node[not gate] (N1) at (4.3,0.55) {};
    \node[not gate] (N2) at (3.0,-1.1) {};
    \node[or gate,inputs=nn] (O2) at (6.0,0) {};
    \coordinate (M1) at (5.35,0);
    \coordinate (M2) at (5.15,0);
    \draw (O1.input 1) -- (O1.input 1 -| L) node[left] {$x$};
    \draw (A.input 1) -- (A.input 1 -| L) node[left] {$y$};
    \draw (A.input 2) -- (A.input 2 -| L) node[left] {$z$};
    \fill (A.input 2 -| B) circle (1.5pt);
    \draw (A.input 2 -| B) |- (N2.input);
    \draw (A.output) -- ++(0.25,0) |- (O1.input 2);
    \draw (O1.output) -- (N1.input);
    \draw (N1.output) -- (N1.output -| M1) |- (O2.input 1);
    \draw (N2.output) -- (N2.output -| M2) |- (O2.input 2);
    \draw (O2.output) -- ++(0.4,0) node[right] {$f$};
  \end{tikzpicture}
\end{minipage}
\notslides{De Morgan simplifies it further:
  $\bigl(x\vee(y\wedge z)\bigr)'=x'\wedge(y'\vee z')$, so
  $f=(x'\wedge y')\vee(x'\wedge z')\vee z'=(x'\wedge y')\vee z'$, since
  $z'$ absorbs $x'\wedge z'$.}

\begin{keybox}[title=Key points (6.4--6.5)]
  \textbf{Boolean algebra} = finite lattice $\cong B_n$ (subsets of an
  $n$-set, or bit strings of length $n$): bounded, complemented,
  distributive, $2^n$ elements. \textbf{$2^n$ elements is necessary, not
    sufficient.} \textbf{$D_n$} is a Boolean algebra $\iff$ $n$ is a product
  of distinct primes. $\vee,\wedge,'$ behave like $\cup,\cap$, complement,
  so every set law carries over --- especially absorption and De Morgan.
  A \textbf{Boolean function} $B^n\to B$ is a truth table with $2^n$ rows;
  a \textbf{Boolean polynomial} is built from variables, $0$, $1$,
  $\vee$, $\wedge$, $'$. \textbf{Logic diagrams}: OR $=\vee$, AND
  $=\wedge$, NOT $='$; NAND $=(\wedge)'$, NOR $=(\vee)'$. Read circuits
  left to right, one gate at a time.
\end{keybox}

% ---------------------------------------------------------------------
\kolman{6}
\section{Circuit design}\src{Kolman 6.6, pp.~284--292}
\index{circuit design}\index{minterm}\index{Boolean function!expression from a truth table}\index{truth table!to Boolean expression}\index{S(f)@$S(f)$}%

6.5 goes from an expression to a truth table. This section goes the other
way: given only the truth table of $f:B^n\to B$, find an expression ---
and so a circuit --- that produces it. For $f:B^n\to B$ let
\[
  S(f)=\{\,b\in B^n \mid f(b)=1\,\}\qquad\text{(the rows where $f$ is $1$).}
\]

\begin{rulebox}[title={6.6 Theorem 1}]
  Let $f,f_1,f_2:B^n\to B$.
  (a) If $S(f)=S(f_1)\cup S(f_2)$, then $f=f_1\vee f_2$.
  (b) If $S(f)=S(f_1)\cap S(f_2)$, then $f=f_1\wedge f_2$.
\end{rulebox}

\begin{defn}
  For an input $b=(c_1,\dots,c_n)\in B^n$, the \textbf{minterm} $E_b$ is
  $\bar x_1\wedge\bar x_2\wedge\dots\wedge\bar x_n$, where
  $\bar x_k=x_k$ if $c_k=1$ and $\bar x_k=x_k'$ if $c_k=0$. It is $1$ on
  the single input $b$ and $0$ everywhere else.
\end{defn}
For two variables: $E_{00}=x'\wedge y'$, $E_{01}=x'\wedge y$,
$E_{10}=x\wedge y'$, $E_{11}=x\wedge y$. For three:
$E_{011}=x'\wedge y\wedge z$, and so on --- a $0$ in the input becomes a
primed variable.

\begin{rulebox}[title={6.6 Theorem 2: every Boolean function has an expression}]
  Any $f:B^n\to B$ is produced by a Boolean expression: if
  $S(f)=\{b_1,\dots,b_k\}$, then
  \[
    f=E_{b_1}\vee E_{b_2}\vee\dots\vee E_{b_k}
  \]
  --- one minterm per row where $f=1$, joined by $\vee$. (Proof: each
  $E_{b_i}$ has $S=\{b_i\}$; apply Theorem 1(a).) If $f$ is never $1$,
  $f=0$.
\end{rulebox}

\begin{rulepar}
  \textbf{Method: truth table $\to$ circuit.}
  \begin{enumerate}[nosep]
    \item Mark the rows with $f=1$; write the minterm of each.
    \item Join the minterms with $\vee$ (Theorem 2). This is already a
          correct answer.
    \item Simplify with the laws (distributive, $x\vee x'=1$,
          $x\wedge1=x$, absorption, De Morgan) or a Karnaugh map (below).
          Fewer terms means fewer gates.
    \item Draw it: one NOT per primed variable, one AND per term, one OR
          joining the terms.
    \item \textbf{Verify}: a truth table of the circuit, one helper
          column per gate output, and check the last column against the
          given table row by row.
  \end{enumerate}
\end{rulepar}
\textbf{Allowed gates.} If a task says ``only NOT, OR and AND'', a NAND
or NOR gate is not allowed --- draw an AND or OR followed by a NOT
instead. An AND (or OR) of three terms can be drawn as two two-input
gates in a row (associativity).

\begin{example}[6.6 Ex.~3: simplifying a sum of minterms]
  \leavevmode\par\noindent
  \begin{minipage}[t]{0.25\textwidth}
    \small
    \begin{tabular}{@{}ccc|c@{}}
      $x$ & $y$ & $z$ & $f$ \\\hline
      0   & 0   & 0   & 0   \\ 0&0&1&0\\ 0&1&0&0\\ 0&1&1&1\\
      1   & 0   & 0   & 0   \\ 1&0&1&0\\ 1&1&0&0\\ 1&1&1&1
    \end{tabular}
  \end{minipage}\hfill
  \begin{minipage}[t]{0.72\textwidth}
    $S(f)=\{011,111\}$, so by Theorem 2
    \begin{align*}
      f & =(x'\wedge y\wedge z)\vee(x\wedge y\wedge z) \\
        & =(x'\vee x)\wedge(y\wedge z)       & & \text{distributive}   \\
        & =1\wedge(y\wedge z)=y\wedge z.     & & x\vee x'=1,\ 1\wedge p=p
    \end{align*}
    The circuit is a single AND gate.
  \end{minipage}
\end{example}

\begin{example}[6.6 Fig.~83: a full worked circuit, with the check]
  \leavevmode\par\noindent
  \begin{minipage}[c]{0.4\textwidth}
    $f(x,y)$ is $1$ on $00$, $01$, $10$ and $0$ on $11$.\\[2pt]
    Minterms: $(x'\wedge y')\vee(x'\wedge y)\vee(x\wedge y')$.\\[2pt]
    First two: $x'\wedge(y'\vee y)=x'$. Then
    $x'\vee(x\wedge y')=(x'\vee x)\wedge(x'\vee y')=x'\vee y'$.\\[2pt]
    So $f=\mathbf{x'\vee y'}$.
  \end{minipage}\hfill
  \begin{minipage}[c]{0.27\textwidth}
    \centering
    \begin{tikzpicture}[gates]
      \node[not gate] (N1) at (1,0.5) {};
      \node[not gate] (N2) at (1,-0.5) {};
      \node[or gate,inputs=nn] (O) at (2.6,0) {};
      \draw (N1.input) -- ++(-0.5,0) node[left] {$x$};
      \draw (N2.input) -- ++(-0.5,0) node[left] {$y$};
      \draw (N1.output) -- ++(0.25,0) |- (O.input 1);
      \draw (N2.output) -- ++(0.25,0) |- (O.input 2);
      \draw (O.output) -- ++(0.4,0) node[right] {$f$};
    \end{tikzpicture}
  \end{minipage}\hfill
  \begin{minipage}[c]{0.28\textwidth}
    \small
    \begin{tabular}{@{}cc|cc|c@{}}
      $x$ & $y$ & $x'$ & $y'$ & $x'\vee y'$ \\\hline
      0   & 0   & 1    & 1    & 1           \\
      0   & 1   & 1    & 0    & 1           \\
      1   & 0   & 0    & 1    & 1           \\
      1   & 1   & 0    & 0    & 0
    \end{tabular}
  \end{minipage}

  \smallskip
  The last column matches $f$ on every row, so the circuit is correct.
  (It is a NAND, but drawn with NOT and OR only.)
\end{example}

\subsection{Karnaugh maps}
\index{Karnaugh map}%
A \textbf{Karnaugh map} writes the values of $f$ in a grid where
\emph{neighbouring squares differ in exactly one variable}; then a block
of $1$'s collapses to a short term, because $x\vee x'=1$ cancels the
variable that changes inside the block.

\begin{center}
  \small
  \begin{tabular}{@{}c@{\qquad\qquad}c@{}}
    two variables & three variables (Kolman's layout) \\[2pt]
    \begin{tabular}{c|cc}
           & $y'$ & $y$ \\\hline
      $x'$ & 00   & 01  \\
      $x$  & 10   & 11
    \end{tabular}
    &
    \begin{tabular}{c|cccc}
           & $y'z'$ & $y'z$ & $yz$ & $yz'$ \\\hline
      $x'$ & 000    & 001   & 011  & 010   \\
      $x$  & 100    & 101   & 111  & 110
    \end{tabular}
  \end{tabular}
\end{center}
The column order $00,01,11,10$ (Gray code, not counting order) is what
makes neighbours differ in one bit. The left and right edges are
\textbf{neighbours too} (glue them into a cylinder).

\begin{rulepar}
  \textbf{Reading a Karnaugh map.} Cover every $1$ with rectangles of
  $1$'s of size $1,2,4$ or $8$ (powers of two; rectangles may overlap and
  may wrap around the edge). Each rectangle gives one term: the
  variables that are \emph{constant} on it (primed if constant $0$).
  Join the terms with $\vee$. \textbf{Bigger rectangles give shorter
    terms}: in the 3-variable map, a single square is a 3-variable
  minterm, two squares give 2 variables, a full row or two adjacent
  columns give 1 variable.
\end{rulepar}

\begin{example}[6.6 Ex.~6]
  \leavevmode\par\noindent
  \begin{minipage}[c]{0.45\textwidth}
    \centering\small
    $S(f)=\{000,011,100,101,111\}$\\[4pt]
    \begin{tabular}{c|cccc}
           & $y'z'$       & $y'z$        & $yz$         & $yz'$ \\\hline
      $x'$ & \textbf{1}   & 0            & \textbf{1}   & 0     \\
      $x$  & \textbf{1}   & \textbf{1}   & \textbf{1}   & 0
    \end{tabular}
  \end{minipage}\hfill
  \begin{minipage}[c]{0.52\textwidth}
    Rectangles: the column $y'z'$ (both rows) gives $y'\wedge z'$; the
    column $yz$ gives $y\wedge z$; the pair $100,101$ in row $x$ gives
    $x\wedge y'$. So
    \[
      f=(y'\wedge z')\vee(x\wedge y')\vee(y\wedge z).
    \]
    (The pair $101,111$, giving $x\wedge z$, would do instead of
    $x\wedge y'$ --- the cover is not unique.)
  \end{minipage}
\end{example}

\subsection{Grouping rules and slide examples}
\index{Karnaugh map!grouping rules}%
\begin{rulepar}
  \textbf{Grouping rules.} (1) Group adjacent cells containing $1$.
  (2) A group has $1,2,4,8,\dots$ cells (a power of $2$).
  (3) Groups are rectangles (squares included).
  (4) Make groups as large as possible.
  (5) Groups may overlap.
  (6) Opposite edges are adjacent --- the map wraps around.
  (7) Every $1$ must be covered by at least one group.
  Reading a group: variables that stay constant stay (value $1$: plain,
  value $0$: primed); the variable that changes is dropped.
\end{rulepar}
The slides draw the three-variable map with $AB$ across the top
($00,01,11,10$) and $C$ down the side ($0,1$) --- Kolman's layout above
is the same grid turned a quarter-turn; use whichever you like, but keep the
Gray order on the two-variable axis.
\begin{itemize}[nosep]
  \item \textbf{Two variables.} $f(0,0)=f(0,1)=1$, $f(1,0)=f(1,1)=0$: one
        group of two in row $x'$; $y$ changes (dropped), $x=0$ stays, so
        $f=x'$. (Check by laws: $x'y'\vee x'y=x'(y'\vee y)=x'$.)
  \item \textbf{Three variables.} $F(A,B,C)=1$ exactly on $010,100,101,110$.
        Horizontal group $\{010,110\}$: $A$ changes, $B=1$, $C=0$ gives
        $BC'$. Vertical group $\{100,101\}$: $A=1$, $B=0$, $C$ changes gives
        $AB'$. So $F=BC'+AB'$.
  \item \textbf{Exercise 1.} Rows $x=0,1$, columns $yz=00,01,11,10$,
        both rows $1\ 1\ 0\ 0$: one group of four (columns $00,01$), only
        $y=0$ is constant, so $f=y'$.
  \item \textbf{Exercise 2.} Row $x'$: $1\ 0\ 0\ 1$; row $x$: $1\ 0\ 1\ 1$.
        Columns $00$ and $10$ are neighbours across the edge, so the block
        $\{x',x\}\times\{00,10\}$ has $z=0$ constant: $z'$. The leftover
        $1$ at $x\,yz=111$ pairs with $110$ (row $x$, columns $11,10$):
        $x\wedge y$. So $f=z'\vee(x\wedge y)$.
\end{itemize}

\notslides{Four variables $w,x,y,z$ (Kolman pp.~291--292; on the 2024 MC2
  paper): a $4\times4$ grid, rows $wx$ and columns $yz$ both in Gray
  order $00,01,11,10$; left--right \emph{and} top--bottom edges wrap.
  Rectangles of $1,2,4,8$ squares give terms of $4,3,2,1$ variables.
  To check an MC option quickly, evaluate it on a few squares (a $1$ and
  a $0$) rather than deriving the expression from scratch.}

\begin{keybox}[title=Key points (6.6)]
  $S(f)$ = the inputs where $f=1$. \textbf{Minterm} $E_b$: AND of all
  variables, primed where $b$ has a $0$; it is $1$ only at $b$.
  \textbf{Theorem 2}: $f$ = OR of the minterms of its $1$-rows --- every
  truth table has an expression, hence a circuit (NOT/AND/OR suffice).
  Simplify with the laws or a \textbf{Karnaugh map} (Gray-order columns,
  edges wrap, biggest power-of-two rectangles). Always \textbf{verify}
  the circuit with a truth table, one column per gate.
\end{keybox}

% =====================================================================
\newpage
\part{Trees}

\kolman{1}
\section{Rooted trees}\src{Kolman 7.1}
\index{tree}\index{rooted tree}\index{root}\index{parent}\index{child}\index{sibling}\index{leaf}\index{level}\index{height of a tree}\index{ordered tree}\index{n-tree@$n$-tree}\index{complete n-tree@complete $n$-tree}\index{binary tree}\index{descendant}\index{subtree}\index{tree!is this relation a tree?}\index{in-degree}%

\begin{defn}
  Let $T$ be a relation on a set $A$. $T$ is a \textbf{tree} if there is a
  vertex $v_0\in A$ such that there is a \emph{unique path} from $v_0$ to
  every other vertex, and there is \emph{no path} from $v_0$ to itself.
  $v_0$ is the \textbf{root}, and $(T,v_0)$ is a \textbf{rooted tree}.
  Elements of $A$ are \textbf{vertices}.
\end{defn}
\erratum{The slide says ``a unique path from $v_0$ to \emph{every}
  vertex $v\in A$'', which would include $v_0$ itself and contradict the
  second bullet. Kolman says every \emph{other} vertex, $v\neq v_0$.}

\begin{rulebox}[title={7.1 Theorem 1}]
  In a rooted tree $(T,v_0)$: (a) $T$ has no cycles; (b) $v_0$ is the
  unique root; (c) $\operatorname{indeg}(v_0)=0$ and every other vertex has
  $\operatorname{indeg}(v)=1$.
\end{rulebox}

\textbf{Terminology} (draw the root at the top, edges point down):
\begin{itemize}[nosep]
  \item \textbf{parent}: the vertex directly above; \textbf{children}:
        the vertices directly below; \textbf{siblings}: same parent;
  \item \textbf{leaf}: a vertex with no children;
  \item \textbf{level}: the root is level $0$, its children level $1$, and
        so on; the \textbf{height} of the tree is its highest level.
\end{itemize}

\begin{center}
  \begin{tikzpicture}[level distance=0.9cm,
      level 1/.style={sibling distance=22mm},level 2/.style={sibling distance=8mm}]
    \node[tn]{$v_0$}
    child{node[tn]{$v_1$} child{node[tn]{$v_4$}}}
    child{node[tn]{$v_2$} child{node[tn]{$v_5$}} child{node[tn]{$v_6$}}}
    child{node[tn]{$v_3$} child{node[tn]{$v_7$}} child{node[tn]{$v_8$}} child{node[tn]{$v_9$}}};
    \node[right] at (3.4,0) {level $0$};
    \node[right] at (3.4,-0.9) {level $1$};
    \node[right] at (3.4,-1.8) {level $2$};
  \end{tikzpicture}
\end{center}
Here the leaves are $v_4,\dots,v_9$ and the height is $2$.

\begin{example}[7.1 Ex.~2 (answer not on slides)]
  $A=\{v_1,\dots,v_{10}\}$ and
  \[
    T=\{(v_2,v_3),(v_2,v_1),(v_4,v_5),(v_4,v_6),(v_5,v_8),(v_6,v_7),
    (v_4,v_2),(v_7,v_9),(v_7,v_{10})\}.
  \]
  Show $T$ is a rooted tree and find the root.\\
  Only $v_4$ is never the second coordinate of a pair (in-degree $0$), so
  it is the only candidate root. Unique paths from $v_4$:
  $v_4\to v_5\to v_8$; $v_4\to v_6\to v_7\to v_9$ and $v_{10}$;
  $v_4\to v_2\to v_3$ and $v_1$. All ten vertices are reached exactly
  once, so $T$ is a rooted tree with \textbf{root $v_4$}. (Check: $10$
  vertices, $9$ edges, every other vertex has in-degree $1$.)
\end{example}

\textbf{Ordered tree}: a rooted tree whose children of each vertex have a
specified order ($v_1<v_2<v_3$). Swapping the left-to-right order of the
children gives a \emph{different} ordered tree.

\begin{rulebox}[title={7.1 Theorem 2: a tree is like a parent relation}]
  Read $(a,b)\in T$ as ``$a$ is the parent of $b$''. Then $T$ is
  \textbf{irreflexive} (no one is their own parent), \textbf{asymmetric}
  (if $a$ is $b$'s parent, $b$ is not $a$'s) and, for $(a,b),(b,c)\in T$,
  $(a,c)\notin T$: the grandparent is not the parent --- so $T$ is
  \textbf{not transitive}.
\end{rulebox}

\begin{defn}
  A rooted tree is an \textbf{$n$-tree} if every vertex has \emph{at most}
  $n$ children, and a \textbf{complete $n$-tree} if every non-leaf has
  \emph{exactly} $n$ children (so $0$ or $n$ children each). For $n=2$:
  a \textbf{binary tree}, whose children are told apart as \emph{left} and
  \emph{right}.
\end{defn}
A vertex $v_2$ is a \textbf{descendant} of $v_1$ if there is a path
$v_1\rightsquigarrow v_2$. \textbf{7.1 Theorem 3:} $v$ together with all
its descendants is again a tree, the \textbf{subtree} $T(v)$ rooted at
$v$.

\Needspace{9\baselineskip}
\begin{example}[Subtrees (7.1 Ex.~3 on the slides)]\leavevmode\par\noindent
  \begin{minipage}[c]{0.36\textwidth}
    \centering
    \begin{tikzpicture}[level distance=0.9cm,
        level 1/.style={sibling distance=16mm},level 2/.style={sibling distance=8mm},
        level 3/.style={sibling distance=8mm}]
      \node[tn]{$v_4$}
      child{node[tn]{$v_6$} child{node[tn]{$v_7$} child{node[tn]{$v_9$}} child{node[tn]{$v_{10}$}}}}
      child{node[tn]{$v_2$} child{node[tn]{$v_1$}} child{node[tn]{$v_3$}}}
      child{node[tn]{$v_5$} child{node[tn]{$v_8$}}};
    \end{tikzpicture}
  \end{minipage}\hfill
  \begin{minipage}[c]{0.6\textwidth}
    \begin{itemize}[nosep]
      \item $T(v_6)=\{v_6,v_7,v_9,v_{10}\}$
      \item $T(v_7)=\{v_7,v_9,v_{10}\}$
      \item $T(v_2)=\{v_2,v_1,v_3\}$
      \item $T(v_8)=\{v_8\}$ (a leaf is a one-vertex subtree)
    \end{itemize}
  \end{minipage}
\end{example}

\begin{example}[7.1 Ex.~3--4]
  Is $R$ a tree? If so, find the root.\\
  \textbf{3.} $A=\{a,b,c,d,e,f\}$, $R=\{(a,b),(c,e),(f,a),(f,c),(f,d)\}$.
  Only $f$ has in-degree $0$; paths $f\to a\to b$, $f\to c\to e$, $f\to d$
  reach every vertex once. Tree, \textbf{root $f$}.\\
  \textbf{4.} $A=\{1,\dots,6\}$, $R=\{(2,1),(3,4),(5,2),(6,5),(6,3)\}$.
  Only $6$ has in-degree $0$; paths $6\to5\to2\to1$ and $6\to3\to4$.
  Tree, \textbf{root $6$}.
\end{example}

\begin{example}[7.1 Ex.~9--13]\leavevmode\par\noindent
  \begin{center}
    \begin{tikzpicture}[level distance=0.8cm,
        level 1/.style={sibling distance=32mm},level 2/.style={sibling distance=9mm},
        level 3/.style={sibling distance=9mm},level 4/.style={sibling distance=9mm}]
      \node[tn]{$v_0$}
      child{node[tn]{$v_1$} child{node[tn]{$v_4$} child{node[tn]{$v_{10}$}} child{node[tn]{$v_{11}$}}}}
      child{node[tn]{$v_2$} child{node[tn]{$v_5$}} child{node[tn]{$v_6$} child{node[tn]{$v_{12}$}}}}
      child{node[tn]{$v_3$} child{node[tn]{$v_7$}}
          child{node[tn]{$v_8$} child{node[tn]{$v_{13}$} child{node[tn]{$v_{15}$}}} child{node[tn]{$v_{14}$}}}
          child{node[tn]{$v_9$}}};
    \end{tikzpicture}
  \end{center}
  \begin{itemize}[nosep]
    \item \textbf{9.} Level-$3$ vertices: $v_{10},v_{11},v_{12},v_{13},v_{14}$.
          Leaves: $v_5,v_7,v_9,v_{10},v_{11},v_{12},v_{14},v_{15}$.
    \item \textbf{10.} Siblings of $v_8$ (same parent $v_3$): $v_7,v_9$.
          Descendants of $v_8$: $v_{13},v_{14},v_{15}$.
    \item \textbf{11.} $T(v_2)=\{v_2,v_5,v_6,v_{12}\}$;
          $T(v_3)=\{v_3,v_7,v_8,v_9,v_{13},v_{14},v_{15}\}$.
    \item \textbf{12.} Height of $(T,v_0)$: longest path
          $v_0-v_3-v_8-v_{13}-v_{15}$ has $4$ edges, so $4$. Height of $T(v_3)$:
          $v_3-v_8-v_{13}-v_{15}$, so $3$.
    \item \textbf{13.} The most children any vertex has is $3$, so it is a
          \textbf{$3$-tree}; but $v_1$ has only one child, so it is
          \textbf{not complete}.
  \end{itemize}
\end{example}

\subsection{Exam method: which relation (or matrix) is a (directed) tree?}
\index{tree!which relation is a tree?}\index{tree!which matrix is a tree?}\index{in-degree!column sums}\index{directed tree}%
A ``directed tree'' or ``rooted directed tree'' in an MC question is just
Kolman's \emph{tree} (7.1): a relation $T$ on $A$, read as a digraph with an
arrow $a\to b$ for each $(a,b)\in T$, with one root and a unique path from it
to every other vertex. It must be a tree \emph{on all of $A$}.

Facts to use: no cycles (so no loops $(x,x)$ and no symmetric pair $(x,y),(y,x)$); the root has in-degree $0$ and every other vertex in-degree exactly $1$; so $n$ vertices means exactly $n-1$ edges.

\textbf{Method for ``which relation is a tree?'':}
\index{tree!which relation is a tree?}%
\begin{enumerate}[nosep]
  \item Cross out any relation with a loop $(x,x)$ or a pair $(x,y),
          (y,x)$.
  \item Count edges: must be (number of vertices) $- 1$.
  \item Tally the \emph{second} coordinates (in-degrees). Exactly one
        vertex of $A$ may be missing --- the root; none may appear twice.
        A vertex of $A$ that appears in no pair at all is unreachable.
  \item From the root, follow the arrows and confirm every vertex is
        reached (rules out a cycle sitting apart from the root).
\end{enumerate}

\begin{example}[On $A = \{1,2,3,4,5\}$: tree or not?]
  \begin{itemize}[nosep]
    \item $\{(1,2), (2,3), (3,1), (4,5)\}$: cycle $1 \to 2 \to 3 \to 1$,
          and $4$ and $5$ cannot be reached from it. Not a tree.
    \item $\{(4,1), (4,2), (2,3), (1,3)\}$: $3$ appears twice as a second
          coordinate (two parents), $5$ never appears. Not a tree.
    \item $\{(3,1), (3,5), (5,2), (5,4)\}$: $4 = 5 - 1$ edges; second
          coordinates $1, 5, 2, 4$ --- each once, only $3$ missing, so
          the root is $3$. $3 \to 1$, $3 \to 5 \to 2$, $5 \to 4$ reaches
          everything. \textbf{A tree}, root $3$.
  \end{itemize}
\end{example}

\textbf{Same question with a matrix} (``which matrix defines a tree?''):
\index{tree!which matrix is a tree?}\index{tree!from a matrix}\index{in-degree!column sums}%
row $i$, column $j$ is a $1$ for the pair $(i,j)$, i.e.\ an arrow
$i \to j$. The method above becomes:
\begin{enumerate}[nosep]
  \item \textbf{Diagonal all $0$} --- any $1$ on the diagonal is a loop.
  \item \textbf{Count the $1$'s:} must be (number of vertices) $- 1$.
  \item \textbf{Column sums are the in-degrees:} exactly one column of
        zeros (the root), every other column sums to $1$.
  \item \textbf{Follow the arrows from the root} (read each row for
        where an arrow leads) and check every vertex is reached.
\end{enumerate}
Steps 1--3 can all pass and the relation still fail: a cycle can sit
apart from the root and satisfy the in-degree test by itself.

\begin{example}[Which $5 \times 5$ matrix is a tree?]
  \begin{itemize}[nosep]
    \item \textbf{Pairs} $(1,2), (2,1), (3,4), (5,3)$: $4$ ones, column
          sums $1,1,1,1,0$, so the root would be $5$. But from $5$ you
          only reach $3$ and $4$; $1 \leftrightarrow 2$ is a cycle cut
          off from the root. \textbf{Not a tree.}
    \item \textbf{Pairs} $(1,2), (3,1), (3,4), (4,5)$: $4$ ones, column
          sums $1,1,0,1,1$, root $3$. $3 \to 1 \to 2$ and
          $3 \to 4 \to 5$ reach everything. \textbf{A tree.}
  \end{itemize}
\end{example}

% ---------------------------------------------------------------------
\kolman{2}
\section{Labeled trees}\src{Kolman 7.2}
\index{labeled tree}\index{expression tree}\index{positional tree}\index{Huffman code}\index{array representation of a tree}%

Vertices can carry labels (numbers, words, operators). The \emph{structure}
says how the objects are related, the \emph{labels} say what they are:
categories, decisions, expressions, programs, data structures.

\textbf{Expression trees.} Each operator is the \emph{central operator} of
its subexpression: it sits at the root of that subtree, with the left and
right operands as its children. For $(3-(2\times x))+((x-2)-(3+x))$:
\begin{center}
  \begin{tikzpicture}[level distance=0.85cm,
      level 1/.style={sibling distance=42mm},level 2/.style={sibling distance=22mm},
      level 3/.style={sibling distance=10mm}]
    \node[ex]{$+$}
    child{node[ex]{$-$} child{node[ex]{$3$}} child{node[ex]{$\times$} child{node[ex]{$2$}} child{node[ex]{$x$}}}}
    child{node[ex]{$-$}
        child{node[ex]{$-$} child{node[ex]{$x$}} child{node[ex]{$2$}}}
        child{node[ex]{$+$} child{node[ex]{$3$}} child{node[ex]{$x$}}}};
  \end{tikzpicture}
\end{center}

\textbf{Positional trees.} A \emph{positional} tree is an ordered tree
where each child has a fixed position: $1,2,\dots,n$ in an $n$-tree (left,
middle, right for $n=3$). The position stays fixed \emph{even if some
  children are missing}: a vertex may have only its position-$3$ child. A
\textbf{positional binary tree} has positions $L$ and $R$; a vertex may
have both, only $L$, only $R$, or neither.

\textbf{Storing a binary tree.} One cell per vertex: \texttt{LEFT |
  DATA | RIGHT}, where LEFT and RIGHT are the \emph{indices} of the child
cells, and $0$ means ``no child'' (null). The first cell (\emph{Start})
holds the index of the root in its LEFT field.

\begin{example}[7.2 Ex.~1: the expression tree above, as an array]
  \leavevmode\par\noindent
  \begin{center}\small
    \begin{tabular}{@{}c|ccc@{}}
      Index & L  & Data     & R  \\\hline
      1     & 2  & Start    & 0  \\ 2 & 3 & $+$ & 8\\ 3 & 4 & $-$ & 5\\ 4 & 0 & $3$ & 0\\
      5     & 6  & $\times$ & 7  \\ 6 & 0 & $2$ & 0\\ 7 & 0 & $x$ & 0\\ 8 & 9 & $-$ & 12\\
      9     & 10 & $-$      & 11 \\ 10 & 0 & $x$ & 0\\ 11 & 0 & $2$ & 0\\ 12 & 13 & $+$ & 14\\
      13    & 0  & $3$      & 0  \\ 14 & 0 & $x$ & 0
    \end{tabular}
  \end{center}
  Read it: Start points to index $2$ (the root $+$); its LEFT is $3$
  (the left $-$), its RIGHT is $8$, and so on.
\end{example}
\begin{example}[7.2 Ex.~17 (the table on the slide)]
  \leavevmode\par\noindent
  \begin{center}\small
    \begin{tabular}{@{}c|ccc@{}}
      Index & LEFT & DATA  & RIGHT \\\hline
      1     & 8    & Start & 0     \\ 2 & 5 & D & 7\\ 3 & 9 & E & 0\\ 4 & 2 & C & 3\\
      5     & 0    & F     & 0     \\ 6 & 0 & B & 4\\ 7 & 0 & G & 0\\ 8 & 6 & A & 0\\ 9 & 0 & H & 0
    \end{tabular}
  \end{center}
  Draw the digraph by the same reading rule; the slides give no solution.
\end{example}

\textbf{Huffman codes.} In a fixed-length code every character uses the
same number of bits; in a \textbf{Huffman code} frequent characters get
\emph{shorter} codes. A positional binary tree stores the code: $0$ = go
left, $1$ = go right, and the characters sit at the leaves.

\begin{example}[7.2 Ex.~4: decoding]\leavevmode\par\noindent
  \begin{minipage}[c]{0.4\textwidth}
    \centering
    \begin{tikzpicture}[x=1cm,y=0.8cm,every node/.style={font=\small},
        lab/.style={font=\scriptsize,midway,fill=white,inner sep=1pt}]
      \node (r) at (0,0) {$\bullet$};
      \node (E) at (-1,-1) {E};
      \node (n1) at (1,-1) {$\bullet$};
      \node (A) at (0,-2) {A};
      \node (n2) at (2,-2) {$\bullet$};
      \node (S) at (1,-3) {S};
      \node (n3) at (3,-3) {$\bullet$};
      \node (R) at (2,-4) {R};
      \node (C) at (4,-4) {C};
      \draw (r) -- node[lab]{0} (E);  \draw (r) -- node[lab]{1} (n1);
      \draw (n1) -- node[lab]{0} (A); \draw (n1) -- node[lab]{1} (n2);
      \draw (n2) -- node[lab]{0} (S); \draw (n2) -- node[lab]{1} (n3);
      \draw (n3) -- node[lab]{0} (R); \draw (n3) -- node[lab]{1} (C);
    \end{tikzpicture}
  \end{minipage}\hfill
  \begin{minipage}[c]{0.56\textwidth}
    Codes: E $=0$, A $=10$, S $=110$, R $=1110$, C $=1111$. Decode by
    walking from the root; at a leaf, write the letter and return to the
    root.
    \[
      0101100=\underbrace{0}_{E}\;\underbrace{10}_{A}\;\underbrace{110}_{S}\;
      \underbrace{0}_{E}\;\Rightarrow\;\text{EASE}.
    \]
  \end{minipage}
  \notslides{The question lists a second string, $0101110$, but the slides
    only solve the first: $0\,|\,10\,|\,1110$ gives \textbf{EAR}. One
    flipped bit turns EASE into EAR --- the weakness of a variable-length
    code (Kolman Example 4).}
\end{example}

% ---------------------------------------------------------------------
\kolman{3}
\section{Tree searching}\src{Kolman 7.3}
\index{tree traversal}\index{preorder}\index{inorder}\index{postorder}\index{Polish notation}\index{prefix form}\index{prefix form!evaluating}%

To do something with \emph{every} vertex (compute, print, count, search)
we need a systematic order: \textbf{tree traversal}. For a vertex $v$ with
left child $v_L$ and right child $v_R$ there are three actions: $V$ (visit
$v$), $L$ (traverse the left subtree), $R$ (traverse the right subtree).

\begin{rulepar}
  \begin{center}
    \begin{tabular}{l|ccc|l}
                         & first & second & third & the root is visited         \\\hline
      \textbf{Preorder}  & $V$   & $L$    & $R$   & \emph{before} the subtrees  \\
      \textbf{Inorder}   & $L$   & $V$    & $R$   & \emph{between} the subtrees \\
      \textbf{Postorder} & $L$   & $R$    & $V$   & \emph{after} the subtrees
    \end{tabular}
  \end{center}
\end{rulepar}
Each step is recursive: ``traverse the left subtree'' means apply the same
rule to the subtree rooted at $v_L$, if it exists.

\begin{example}[Slide 44: all three traversals]\leavevmode\par\noindent
  \begin{minipage}[c]{0.3\textwidth}
    \centering
    \begin{tikzpicture}[level distance=0.9cm]
      \node[tn]{1}
      child{node[tn]{2} child{node[tn]{4}} child{node[tn]{5}}}
      child{node[tn]{3} child[missing] child{node[tn]{6}}};
    \end{tikzpicture}
  \end{minipage}\hfill
  \begin{minipage}[c]{0.66\textwidth}
    \begin{itemize}[nosep]
      \item Preorder: $1,2,4,5,3,6$
      \item Inorder: $4,2,5,1,3,6$
      \item Postorder: $4,5,2,6,3,1$
    \end{itemize}
  \end{minipage}
\end{example}

\textbf{Polish (prefix) form} writes arithmetic without parentheses by
putting the operator \emph{before} its operands: infix $6-4$ becomes
prefix $-\,6\,4$.
\notslides{Traversing an expression tree in \emph{preorder} gives its
  Polish (prefix) form, \emph{inorder} gives the ordinary infix form (without
  parentheses it can be ambiguous), and \emph{postorder} gives the postfix
  (reverse Polish) form.}

\begin{rulepar}
  \textbf{Evaluating a prefix expression:} find an operator that is
  immediately followed by two numbers, replace the three symbols by the
  result, and repeat.
\end{rulepar}

\begin{example}[Evaluate $\times\,-\,6\,4\,+\,5\,\div\,2\,2$]
  $-\,6\,4=2$ and $\div\,2\,2=1$, giving $\times\,2\,+\,5\,1$; then
  $+\,5\,1=6$, giving $\times\,2\,6=\mathbf{12}$.
\end{example}
\begin{example}[7.3 Ex.~19: $\times\,-\,+\,3\,4\,-\,7\,2\,\div\,12\,\times\,3\,-\,6\,4$]
  \leavevmode\par\noindent
  \begin{tabbing}
    $\times\,-\,+\,3\,4\,-\,7\,2\,\div\,12\,\times\,3\,-\,6\,4$\quad\=($+\,3\,4=7$)\\
    $\times\,-\,7\,-\,7\,2\,\div\,12\,\times\,3\,-\,6\,4$\>($-\,7\,2=5$)\\
    $\times\,-\,7\,5\,\div\,12\,\times\,3\,-\,6\,4$\>($-\,7\,5=2$)\\
    $\times\,2\,\div\,12\,\times\,3\,-\,6\,4$\>($-\,6\,4=2$)\\
    $\times\,2\,\div\,12\,\times\,3\,2$\>($\times\,3\,2=6$)\\
    $\times\,2\,\div\,12\,6$\>($\div\,12\,6=2$)\\
    $\times\,2\,2=\mathbf{4}$
  \end{tabbing}
\end{example}

\subsection{Exam method: traversal and Polish notation}
\index{tree!traversal}\index{Polish notation}\index{postfix form}\index{infix form}\index{prefix form}\index{reverse Polish notation|see{postfix}}%
An algebraic expression is a \textbf{positional binary tree}: operators
at the inner vertices, numbers and variables at the leaves, left child =
left operand. The three traversals differ only in \emph{when the root is
  written} --- same recursion, the \texttt{visit(v)} line moved:
\begin{rulepar}
  \begin{center}
    \renewcommand{\arraystretch}{1.3}
    \begin{tabular}{@{}llll@{}}
      \toprule
      Algorithm & Order             & Gives                                  & Operator sits \\
      \midrule
      PREORDER  & root, left, right & prefix / \textbf{Polish} form          & first         \\
      INORDER   & left, root, right & infix (ordinary, needs parentheses)    & middle        \\
      POSTORDER & left, right, root & postfix / \textbf{reverse Polish} form & last          \\
      \bottomrule
    \end{tabular}
  \end{center}
\end{rulepar}
Postfix is what a stack machine (an HP calculator, the JVM) executes:
push operands, and each operator pops two and pushes the result.

\textbf{Method:} write the expression with full parentheses, one pair per
operator, then recurse: for postfix, ``left part, right part, operator''.
\textbf{Sanity checks on MC options:} postfix must \emph{start with a leaf}
and \emph{end with the root operator}; prefix the reverse. Every leaf
appears exactly once, and the leaves keep their left-to-right order in
all three forms --- only the operators move.

\begin{example}[Postfix form of $\bigl(1 - a \times a\bigr) \div \bigl(1 + a\bigr)$]
  \leavevmode\par\noindent
  \begin{minipage}[c]{0.3\textwidth}
    \centering
    \begin{tikzpicture}[every node/.style={inner sep=1.5pt}, line width=0.5pt,
        level distance=0.9cm,
        level 1/.style={sibling distance=1.8cm},
        level 2/.style={sibling distance=0.9cm}]
      \node {$\div$}
      child {node {$-$}
          child {node {$1$}}
          child {node {$\times$} child {node {$a$}} child {node {$a$}}}}
      child {node {$+$}
          child {node {$1$}}
          child {node {$a$}}};
    \end{tikzpicture}
  \end{minipage}\hfill
  \begin{minipage}[c]{0.66\textwidth}
    \begin{itemize}[nosep]
      \item left subtree $1 - (a \times a)$: $\;1\ \ a\,a\,\times\ \ -$
      \item right subtree $1 + a$: $\;1\,a\,+$
      \item root last: $\;\mathbf{1\,a\,a \times - \;1\,a + \;\div}$
    \end{itemize}
    Prefix (PREORDER) would be $\div\, -\, 1 \times a\,a\, +\, 1\,a$.
    \textbf{Traps:} the prefix form is offered as a distractor; the
    infix $1 - a \times a \div 1 + a$ without parentheses means
    something else; and the evaluated value $\frac{1-a^2}{1+a}$ (or
    $1 - a$ simplified) is not a ``form'' at all.
  \end{minipage}
\end{example}

\begin{keybox}[title=Key points (7.1--7.3)]
  A \textbf{rooted tree}: a unique path from the root to every other
  vertex, in-degree $0$ at the root and $1$ everywhere else, no cycles.
  Terms: parent, child, sibling, leaf, level, \textbf{height} (the highest
  level), descendants, subtree $T(v)$. \textbf{$n$-tree}: at most $n$
  children; \textbf{complete}: exactly $0$ or $n$. A tree relation is
  irreflexive, asymmetric and \emph{not} transitive. \textbf{Labeled}
  trees: expression trees, positional trees, Huffman codes ($0$ = left,
  $1$ = right, letters at leaves), and the \texttt{LEFT|DATA|RIGHT} array
  (index $0$ = null). \textbf{Traversals}: preorder $VLR$, inorder $LVR$,
  postorder $LRV$ --- ask ``when is the root visited?'' Prefix expressions:
  evaluate an operator followed by two numbers, repeat.
\end{keybox}

% ---------------------------------------------------------------------
\kolman{4}
\section{Undirected trees}\src{Kolman 7.4}
\index{undirected tree}\index{tree!undirected}\index{symmetric closure!undirected tree}\index{undirected edge}\index{adjacent vertices}\index{neighbor}\index{simple path}\index{simple cycle}\index{acyclic}\index{connected}%

Take a rooted tree and forget the hierarchy: no parent/child, only the
underlying shape. Formally, an \textbf{undirected tree} is the
\textbf{symmetric closure} of a tree: for every $(a,b)\in T$ add $(b,a)$
as well. The pair $a\to b$, $b\to a$ is drawn as one line $a - b$, the
\textbf{undirected edge} $\{a,b\}=\{b,a\}$. So each line stands for two
ordered pairs, and $(x,y)\in T\iff(y,x)\in T$.

\begin{defn}
  If $\{a,b\}$ is an undirected edge, $a$ and $b$ are \textbf{adjacent
    vertices} (\textbf{neighbors}). A path $v_0,v_1,\dots,v_n$ is
  \textbf{simple} if no undirected edge is used more than once (going
  $a\to b$ and later $b\to a$ uses the edge $\{a,b\}$ twice). A simple
  path with $v_0=v_n$ is a \textbf{simple cycle}. A symmetric relation
  with no simple cycles is \textbf{acyclic}.
\end{defn}
\erratum{Slide 10 (7.4 Ex.~2) switches to ``simple = no \emph{vertex}
  repeated'', contradicting the definition two slides earlier. Use the
  definition above (no \emph{edge} repeated). The example's verdict
  survives either way: $f,e,d,c,d,a$ walks $d\to c\to d$, which uses the
  edge $\{c,d\}$ twice.}

\begin{rulebox}[title={7.4 Theorems 1--3: what makes an undirected tree}]
  Let $R$ be a symmetric relation on a set $A$ with $n$ elements.
  \begin{enumerate}[nosep,label=(\arabic*)]
    \item \textbf{Theorem 1.} $R$ is an undirected tree $\iff$ $R$ is
          \textbf{connected and acyclic}. \quad TREE $=$ CONNECTED $+$
          ACYCLIC.
    \item \textbf{Theorem 2.} $R$ is an undirected tree $\iff$ $R$ is
          acyclic, but adding \emph{any} new edge creates a cycle
          (maximally acyclic) $\iff$ $R$ is connected, but removing
          \emph{any} edge disconnects it (minimally connected).
    \item \textbf{Theorem 3.} A tree with $n$ vertices has exactly
          $\mathbf{n-1}$ \textbf{edges} ($10$ vertices: $9$ edges; $100$
          vertices: $99$).
  \end{enumerate}
\end{rulebox}
\erratum{Slide 13 states Theorem 2 as ``tree iff (1) and (2)''. Each
  condition \emph{on its own} already characterises a tree, so read it as
  ``tree iff (1), and tree iff (2)''.}

\noindent The four cases (slide 12): connected $+$ no cycle $=$ tree;
connected $+$ cycle, disconnected $+$ no cycle (a \emph{forest}) and
disconnected $+$ cycle are not trees. Example: $A=\{1,2,3,4\}$,
$R=\{\{1,2\},\{2,3\},\{3,4\}\}$ is the path $1-2-3-4$: connected, acyclic,
$3=4-1$ edges, so a tree.

\subsection{Spanning trees}
\index{spanning tree}\index{tree!spanning}%
\begin{defn}
  Let $R$ be symmetric and connected. A \textbf{spanning tree} $T$ of $R$
  is a tree with the same vertices that uses only edges of $R$:
  $V(T)=V(R)$ and $E(T)\subseteq E(R)$.
\end{defn}
A graph usually has \emph{many} spanning trees, but every one of them
contains all vertices, is connected, has no cycles and has exactly
$n-1$ edges. \textbf{By deletion:} while there is a cycle, remove one of
its edges (removing an edge \emph{on a cycle} never disconnects the
graph). Count: you delete $|E(R)|-(n-1)$ edges in total.

\begin{example}[7.4 Ex.~3]\leavevmode\par\noindent
  \begin{minipage}[c]{0.3\textwidth}
    \centering
    \begin{tikzpicture}[ug,x=0.85cm,y=0.85cm]
      \node[pt,label=above:$a$] (a) at (0,2) {};
      \node[pt,label=left:$b$] (b) at (-1.2,1) {};
      \node[pt,label=right:$c$] (c) at (1.2,1) {};
      \node[pt,label=left:$d$] (d) at (0,0) {};
      \node[pt,label=right:$e$] (e) at (1.2,0) {};
      \node[pt,label=below:$f$] (f) at (0,-1) {};
      \draw (a)--(b) (a)--(c) (a)--(d) (b)--(c) (b)--(d) (c)--(d)
            (c)--(e) (d)--(e) (d)--(f) (e)--(f);
      \node at (0,-1.8) {\small $R$};
    \end{tikzpicture}
  \end{minipage}\hfill
  \begin{minipage}[c]{0.3\textwidth}
    \centering
    \begin{tikzpicture}[ug,x=0.85cm,y=0.85cm]
      \node[pt,label=above:$a$] (a) at (0,2) {};
      \node[pt,label=left:$b$] (b) at (-1.2,1) {};
      \node[pt,label=right:$c$] (c) at (1.2,1) {};
      \node[pt,label=left:$d$] (d) at (0,0) {};
      \node[pt,label=right:$e$] (e) at (1.2,0) {};
      \node[pt,label=below:$f$] (f) at (0,-1) {};
      \draw (a)--(b) (a)--(c) (b)--(d) (c)--(e) (d)--(f);
      \node at (0,-1.8) {\small $T'$};
    \end{tikzpicture}
  \end{minipage}\hfill
  \begin{minipage}[c]{0.3\textwidth}
    \centering
    \begin{tikzpicture}[ug,x=0.85cm,y=0.85cm]
      \node[pt,label=above:$a$] (a) at (0,2) {};
      \node[pt,label=left:$b$] (b) at (-1.2,1) {};
      \node[pt,label=right:$c$] (c) at (1.2,1) {};
      \node[pt,label=left:$d$] (d) at (0,0) {};
      \node[pt,label=right:$e$] (e) at (1.2,0) {};
      \node[pt,label=below:$f$] (f) at (0,-1) {};
      \draw (a)--(d) (b)--(d) (c)--(d) (d)--(f) (e)--(f);
      \node at (0,-1.8) {\small $T''$};
    \end{tikzpicture}
  \end{minipage}
  \par\smallskip
  $R$ has $6$ vertices and $10$ edges. Every spanning tree has
  $6-1=\mathbf{5}$ edges, so $5$ edges must go. $T'$ and $T''$ are both
  spanning trees: \textbf{not unique}.
\end{example}

\begin{example}[7.4 Exercise, slide 30 (answer not on slides)]
  \leavevmode\par\noindent
  \begin{minipage}[c]{0.42\textwidth}
    \centering
    \begin{tikzpicture}[ug,x=0.75cm,y=0.75cm]
      \node[pt,label=left:$a$] (a) at (0,0) {};
      \node[pt,label=above:$b$] (b) at (1.5,0.8) {};
      \node[pt,label=below:$c$] (c) at (1.5,-0.8) {};
      \node[pt,label=below:$d$] (d) at (3,0) {};
      \node[pt,label=above:$e$] (e) at (4.5,0.8) {};
      \node[pt,label=below:$g$] (g) at (4.5,-0.8) {};
      \node[pt,label=right:$f$] (f) at (6,0) {};
      \draw (a)--(b)--(d)--(c)--(a) (d)--(e)--(f)--(g)--(d);
    \end{tikzpicture}
  \end{minipage}\hfill
  \begin{minipage}[c]{0.55\textwidth}
    $7$ vertices, $8$ edges, so delete $8-6=2$ edges: one from each of the
    two $4$-cycles $a\,b\,d\,c$ and $d\,e\,f\,g$. E.g.\ delete
    $\{a,b\},\{e,f\}$, or delete $\{c,d\},\{d,g\}$. Since any one of $4$
    edges can go from each cycle, there are $4\cdot4=16$ spanning trees.
  \end{minipage}
\end{example}

\subsection{Spanning trees by vertex merging}
\index{spanning tree!vertex merging}\index{merging vertices}\index{vertex merging}%
Instead of deleting edges, \emph{build} the tree. \textbf{Merging} $a$
and $b$ replaces them by one new vertex $a'=\{a,b\}$ that is adjacent to
everything $a$ or $b$ was adjacent to (two edges to the same vertex
become one). \textbf{Record} the original edge $\{a,b\}$ at each merge ---
the recorded edges \emph{are} the spanning tree.

\begin{rulepar}
  \textbf{Spanning tree by merging} ($R$ symmetric, connected):
  \begin{enumerate}[nosep]
    \item Pick a start vertex $v_1$ and an adjacent vertex $v_2$,
          $\{v_1,v_2\}\in R$.
    \item Record the original edge $\{v_1,v_2\}$; merge $v_1,v_2$ into one
          vertex.
    \item Repeat (merged vertex $+$ any neighbour) until one vertex is
          left. The $n-1$ recorded edges form a spanning tree.
  \end{enumerate}
  \textbf{In the matrix:} merging rows/columns $i$ and $j$ = OR the two
  rows into one, OR the two columns into one, set the diagonal entry to
  $0$, delete the leftover row and column. The matrix shrinks by one each
  step: $n\times n\to\dots\to1\times1$.
\end{rulepar}
Why it cannot make a cycle: each merge joins a vertex \emph{outside} the
merged blob to the blob, and an edge between two different groups never
closes a cycle. (It is Prim's algorithm, below, without weights.)

\begin{example}[7.4 Ex.~7: merging in the matrix]
  The $4$-cycle $a-c-b-d-a$ (edges $\{a,c\},\{a,d\},\{b,c\},\{b,d\}$):
  \[
    \begin{array}{c|cccc}
        & a & b & c & d \\\hline
      a & 0 & 0 & 1 & 1 \\
      b & 0 & 0 & 1 & 1 \\
      c & 1 & 1 & 0 & 0 \\
      d & 1 & 1 & 0 & 0
    \end{array}
    \;\xrightarrow[\text{record }\{a,c\}]{a'=\{a,c\}}\;
    \begin{array}{c|ccc}
         & a' & b & d \\\hline
      a' & 0  & 1 & 1 \\
      b  & 1  & 0 & 1 \\
      d  & 1  & 1 & 0
    \end{array}
    \;\xrightarrow[\text{record }\{c,b\}]{a''=\{a,b,c\}}\;
    \begin{array}{c|cc}
          & a'' & d \\\hline
      a'' & 0   & 1 \\
      d   & 1   & 0
    \end{array}
    \;\xrightarrow[\text{record }\{a,d\}]{a'''}\;
    \begin{array}{c|c}
           & a''' \\\hline
      a''' & 0
    \end{array}
  \]
  Row $a'$ = row $a$ OR row $c$ $=(0,1,1)$ for $(a',b,d)$. Spanning tree
  $\{a,c\},\{c,b\},\{a,d\}$: $3=4-1$ edges. When merging $a'$ with $b$,
  record an \emph{original} edge between the two groups: $\{c,b\}$ (the
  only one, since $a$--$b$ is not an edge). At the last step $\{a,d\}$ or
  $\{b,d\}$ both work. The slides draw the merged vertex without saying
  which original edge to record.
\end{example}
\noindent The slides' second example (slides 24--25) merges $v_0,v_1$,
then $v_2$, in a $7$-vertex graph where $v_0$ and $v_2$ are both adjacent
to $v_5$: after the second merge the two edges to $v_5$ become \emph{one}
undirected edge, which is exactly how the cycle $v_0-v_2-v_5$ is broken.

% ---------------------------------------------------------------------
\kolman{5}
\section{Minimal spanning trees}\src{Kolman 7.5}
\index{weighted graph}\index{weight of an edge}\index{minimal spanning tree}\index{minimum spanning tree|see{minimal spanning tree}}\index{MST|see{minimal spanning tree}}\index{nearest neighbor}\index{Prim's algorithm}\index{greedy algorithm}%

A \textbf{weighted graph} has a number, the \textbf{weight}, on each edge
(distance, cost, travel time). For a vertex $u$, a \textbf{nearest
  neighbor} is a vertex joined to $u$ by a minimum-weight edge (there may
be ties). A \textbf{minimal spanning tree} (MST) is a spanning tree whose
total weight
\[
  w(T)=\sum_{e\in E(T)} w(e)
\]
is as small as possible. Trying every spanning tree is hopeless on a big
graph; the algorithms below build an MST directly.

\begin{rulebox}[title={Prim's algorithm (7.5 Theorem 1: it always gives an MST)}]
  $R$ symmetric, connected, weighted, $n$ vertices.
  \begin{enumerate}[nosep]
    \item Start: $V=\{v_1\}$, $E=\emptyset$.
    \item Choose a minimum-weight edge $\{v_i,v_j\}$ with
          $v_i\in V$ and $v_j\notin V$ (one end \emph{in} the tree, one
          \emph{out}). Ties: pick any.
    \item Add $v_j$ to $V$ and $\{v_i,v_j\}$ to $E$.
    \item Repeat until $|E|=n-1$.
  \end{enumerate}
  The rule is ``the cheapest edge from \emph{my current tree} to a new
  vertex'' --- not the globally cheapest remaining edge. It is
  \textbf{greedy}: the cheapest valid choice now, never looking ahead,
  and Theorem~1 says that is still optimal.
\end{rulebox}

\begin{rulebox}[title={Kruskal's algorithm}]
  Go through \emph{all} edges from lightest to heaviest. Keep an edge if
  it does not close a cycle with the edges kept so far; otherwise skip
  it. Stop at $n-1$ edges. Unlike Prim, the kept edges may form
  \emph{several} separate trees (a forest) along the way; they join up
  at the end.
\end{rulebox}
\noindent\textbf{Both give the same total weight.} The MST itself need
not be unique when weights tie, but its weight is.

\begin{example}[Prim vs.\ Kruskal (slide 47)]\leavevmode\par\noindent
  \begin{minipage}[c]{0.4\textwidth}
    \centering
    \begin{tikzpicture}[ug,x=1cm,y=1cm]
      \node[pt,label=left:$A$] (A) at (0,1) {};
      \node[pt,label=above:$B$] (B) at (1,2) {};
      \node[pt,label=above:$C$] (C) at (2.6,2) {};
      \node[pt,label=right:$F$] (F) at (3.6,1) {};
      \node[pt,label=below:$E$] (E) at (2.6,0) {};
      \node[pt,label=below:$D$] (D) at (1,0) {};
      \draw[mst] (A)--node[w]{4}(B);
      \draw[mst] (A)--node[w]{2}(D);
      \draw[mst] (B)--node[w]{1}(C);
      \draw (B)--node[w,pos=0.35]{5}(D);
      \draw (B)--node[w,pos=0.35]{4}(E);
      \draw[mst] (C)--node[w]{3}(E);
      \draw (C)--node[w]{6}(F);
      \draw[mst] (E)--node[w]{2}(F);
      \draw (D)--node[w]{7}(E);
    \end{tikzpicture}
  \end{minipage}\hfill
  \begin{minipage}[c]{0.57\textwidth}
    \textbf{Prim from $A$:} $A-D$ (2); from $\{A,D\}$ the cheapest way out
    is $A-B$ (4) (not $B-C$ (1): neither end is in the tree yet); then
    $B-C$ (1), $C-E$ (3), $E-F$ (2).\\
    \textbf{Kruskal:} $B-C$ (1), $A-D$ (2), $E-F$ (2), $C-E$ (3), $A-B$
    (4) --- five edges, done; $B-E$ (4) would have closed a cycle.\\
    Same tree (thick), total $\mathbf{12}$.
  \end{minipage}
\end{example}

\begin{example}[7.5 Ex.~1 and Ex.~5: walking trails (answer not on slides)]
  \leavevmode\par\noindent
  \begin{center}
    \begin{tikzpicture}[ug,x=1cm,y=1cm]
      \node[pt,label=left:$A$] (A) at (0,0.9) {};
      \node[pt,label=above:$B$] (B) at (1.5,1.6) {};
      \node[pt,label=below left:$C$] (C) at (2.2,1.1) {};
      \node[pt,label=above:$D$] (D) at (5,1.6) {};
      \node[pt,label=above:$H$] (H) at (6,1.6) {};
      \node[pt,label=above:$E$] (E) at (3.4,0.6) {};
      \node[pt,label=below:$F$] (F) at (2.8,-0.2) {};
      \node[pt,label=below:$G$] (G) at (4.6,-0.2) {};
      \draw[mst] (A)--node[w]{3}(B);
      \draw[mst] (A)--node[w]{2}(C);
      \draw (B)--node[w]{3}(C);
      \draw (B)--node[w]{6}(D);
      \draw[mst] (C)--node[w]{5}(E);
      \draw (C)--node[w]{5}(F);
      \draw[mst] (E)--node[w]{2}(D);
      \draw[mst] (D)--node[w]{2}(H);
      \draw[mst] (E)--node[w]{4}(F);
      \draw[mst] (E)--node[w]{3}(G);
      \draw (F)--node[w]{4}(G);
      \draw (G)--node[w]{6}(H);
    \end{tikzpicture}
  \end{center}
  $8$ vertices, so the MST has $7$ edges.
  \textbf{Prim from $A$:} $A-C$ 2; $A-B$ 3 (tie with $B-C$); $C-E$ 5 (tie
  with $C-F$); $E-D$ 2; $D-H$ 2; $E-G$ 3; $E-F$ 4 (tie with $F-G$).\\
  \textbf{Prim from $E$:} $E-D$ 2; $D-H$ 2; $E-G$ 3; $E-F$ 4; $C-E$ 5;
  $A-C$ 2; $A-B$ 3.\\
  \textbf{Kruskal:} the 2's ($A-C$, $D-E$, $D-H$), the 3's ($A-B$, $E-G$;
  skip $B-C$), $E-F$ 4 (skip $F-G$), $C-E$ 5.\\
  Minimum total distance $2+3+5+2+2+3+4=\mathbf{21}$\,km. The ties mean
  several different MSTs, all of weight $21$.
\end{example}

\subsection{Prim in the matrix}
\index{Prim's algorithm!matrix version}\index{weighted matrix}%
A weighted graph's matrix holds the weight $w$ at $(i,j)$ and $0$ for
``no edge''; $M=M^T$. Prim becomes the merging of 7.4 with weights:
\begin{rulepar}
  Pick the smallest non-zero entry in the merged vertex's row; record that
  edge. To merge rows $i$ and $j$: entry by entry keep the
  \textbf{minimum non-zero} weight (a $0$ means no edge, so it does not
  count as smaller), do the same for the columns, set the diagonal to $0$
  and delete the extra row and column.
\end{rulepar}
\begin{example}[Slide 43 (answer not on slides)]
  Vertices $1$--$4$:
  \[
    M=\begin{bmatrix}0&4&3&0\\4&0&5&3\\3&5&0&2\\0&3&2&0\end{bmatrix}.
  \]
  Start at $1$: row $1$'s smallest is $3$ at column $3$, record $\{1,3\}$.
  Merge rows $1,3$ into $1'$: $(\,\cdot\,,\min(4,5),\cdot\,,\min(\text{--},2))
  =(\,\cdot\,,4,\cdot\,,2)$; smallest $2$ (it came from row $3$), record
  $\{3,4\}$. Merge with $4$: weight to $2$ is $\min(4,3)=3$ (from row
  $4$), record $\{2,4\}$. MST $\{1,3\},\{3,4\},\{2,4\}$, weight
  $3+2+3=\mathbf{8}$.
\end{example}

\subsection{Exam method: spanning-tree and MST questions}
\index{spanning tree!exam method}\index{minimal spanning tree!exam method}%
\begin{itemize}[nosep]
  \item \textbf{How many edges?} Always $n-1$. An option with the wrong
        count is out at once.
  \item \textbf{Is this a spanning tree of $G$?} $n-1$ edges, all taken
        from $G$, every vertex touched, no cycle (with $n-1$ edges,
        ``connected'' and ``no cycle'' imply each other --- check the
        easier one).
  \item \textbf{MST weight:} run Kruskal on paper --- list the edges by
        weight and tick them off, skipping any that join two vertices
        already connected. Fastest by hand, and the total is the same as
        Prim's.
  \item \textbf{``Which edge does Prim pick next?''} Only edges with
        exactly one end in the current tree count. The trap option is the
        cheapest edge \emph{anywhere} (that is Kruskal's choice).
  \item \textbf{Is the MST unique?} If all weights are different, yes.
        With ties it can differ, the weight cannot.
\end{itemize}

\begin{keybox}[title=Key points (7.4--7.5)]
  \textbf{Undirected tree} $=$ symmetric closure of a tree $=$
  \textbf{connected $+$ acyclic}; equivalently maximally acyclic, or
  minimally connected; $n$ vertices $\Rightarrow$ $n-1$ edges.
  \textbf{Simple} path: no edge used twice. \textbf{Spanning tree}: all
  vertices, edges from $R$, connected, acyclic, $n-1$ edges, usually not
  unique; build it by deleting cycle edges or by \textbf{merging}
  vertices (OR the rows, record the edge). \textbf{MST}: spanning tree of
  least total weight. \textbf{Prim} (greedy): grow one tree, always the
  cheapest edge from the tree to a new vertex. \textbf{Kruskal}: cheapest
  edge overall that closes no cycle, may be a forest on the way. Same
  total weight.
\end{keybox}

\newpage
\phantomsection
\addcontentsline{toc}{part}{Exam review and 2023/2024 practice}
\section*{Exam review and past-paper practice}
\index{exam review}\index{past MC questions!2023}\index{past MC questions!2024}%
The last lecture reviews the course and works questions from the 2023 and
2024 MC2 papers. The answers below are the ones highlighted on the slides;
the reasoning is added here.

\begin{rulepar}
  \begin{tabular}{@{}p{0.2\textwidth}p{0.74\textwidth}@{}}
    \textbf{Relations} & Equivalence $=$ R+S+T; partial order $=$ R+A+T. \\
    \textbf{Linear order} & $\forall a,b:\ a\le b\vee b\le a$. $(\Z,\le)$ is linear; $(\Z^+,\mid)$ and $(\mathcal P(A),\subseteq)$ are partial only. \\
    \textbf{Hasse diagram} & lower $=$ smaller; edges are covering pairs; no loops, no shortcuts. Minimal/maximal: nothing strictly below/above. Least/greatest: $\le$ / $\ge$ \emph{every} element. \\
    \textbf{Lattice} & every pair has a join (lub, $\vee$) and a meet (glb, $\wedge$). \textbf{One pair without a join or meet and it is not a lattice.} Product order on $\Z\times\Z$: join $=$ coordinatewise max, meet $=$ coordinatewise min. \\
    \textbf{Boolean algebra} & bounded (least $0$, greatest $1$), distributive, complemented ($a\vee a'=1$, $a\wedge a'=0$). \\
    \textbf{Functions} & injective $f(a)=f(b)\Rightarrow a=b$; surjective every $y$ hit; bijective both; invertible iff bijective. Boolean matrix: \textbf{total} (everywhere defined) $=$ exactly one $1$ per \emph{row}; \textbf{injective} $=$ at most one $1$ per column; \textbf{surjective} $=$ at least one $1$ per column. \\
    \textbf{Trees} & undirected tree: connected, no cycles, $n$ vertices $\Rightarrow$ $n-1$ edges. Rooted directed tree: exactly one root, every other vertex exactly one parent, all reachable from the root. Preorder root--left--right; inorder left--root--right; postorder left--right--root. \\
    \textbf{Growth} & $1\ll\log n\ll n\ll n\log n\ll n^2\ll n^3\ll2^n$; $O$ is an upper bound, $\Theta$ a tight one. \\
  \end{tabular}
\end{rulepar}

\begin{center}
  \small\renewcommand{\arraystretch}{1.25}
  \begin{tabular}{@{}llp{0.62\textwidth}@{}}
    \toprule
    Paper & Question & Answer and why \\
    \midrule
    2023 & Q3 Hasse   & (a) $g$ is a minimal element. $f,g,h$ are the bottom elements; there is no single least element, $a$ is not above $b$, and $g\not\le h$. \\
    2023 & Q4 $D_n$   & (f) $D_{105}$. Boolean iff $n$ is square-free: $105=3\cdot5\cdot7$; the others have a square factor ($63=3^2\cdot7$, $75=3\cdot5^2$, $84=2^2\cdots$, $90=2\cdot3^2\cdot5$, $92=2^2\cdot23$). \\
    2023 & Q5 linear  & (e) $(\{1,2,4\},\mid)$ is the chain $1\mid2\mid4$; the others contain incomparable elements. \\
    2023 & Q7 growth  & (b) $n^3\lg(n^2)+n=2n^3\lg n+n$ has the same order as $n^3\lg n$. \\
    2023 & Q8 join    & (d) $(-1,1)\vee(2,-3)=(\max(-1,2),\max(1,-3))=(2,1)$. \\
    2023 & Q9 functions & (d) $f(2)=f(5)=1$. Not everywhere defined ($6$ has no image), not onto, not one-to-one, $f^{-1}$ not a function. \\
    2023 & Q12 MST    & (f) $14$ (Kruskal/Prim on the weighted graph; the slide's marked answer). \\
    2024 & Q2 tree    & (a) $R_1$: root $1$, every other vertex one parent, all reachable. $R_2$: $3$ has two parents; $R_3$: $5,6,7$ unreachable; $R_4$: $5$ has two parents; $R_5$: loop $(3,3)$; $R_6$: cycle $1\to2\to3\to4\to1$. \\
    2024 & Q9 growth  & (c) $4n^2+n\lg n$ is $\Theta(n^2)$, like $n^2+n+\lg n$. \\
    2024 & Q10 meet   & (b) $(-1,1)\wedge(2,-3)=(\min,\min)=(-1,-3)$. \\
    2024 & Q11 function & (d) $\{(1,3),(2,1),(3,2)\}$: a function with $f(2)=1<f(3)=2$; the others have an element with two images or $f(2)\not<f(3)$. \\
    \bottomrule
  \end{tabular}
\end{center}

\clearpage
\newgeometry{margin=1cm,top=1.8cm,bottom=1.4cm,headheight=14pt}
\phantomsection
\addcontentsline{toc}{part}{Multiplication Table and Primes}
\index{multiplication table}\index{prime numbers!list of primes up to 2500}\index{primes!list}%
\begingroup
\pagestyle{empty}\thispagestyle{empty}
\centering
{\large\textbf{Multiplication table $1$--$30$}}
\quad{\small (grey diagonal = squares)}\par\smallskip
\setlength{\arrayrulewidth}{0.2pt}
\newlength{\mulw}\setlength{\mulw}{\dimexpr(\textwidth-2\arrayrulewidth*31)/31\relax}
\setlength{\tabcolsep}{0pt}\renewcommand{\arraystretch}{1.5}
\newcommand{\mulh}[1]{\cellcolor{black!28}\textbf{#1}}
\fontsize{9}{11}\selectfont
\begin{tabular}{|*{31}{>{\centering\arraybackslash}p{\mulw}|}}
  \hline
  \mulh{$\times$} & \mulh{1}                       & \mulh{2}                       & \mulh{3}                       & \mulh{4}                        & \mulh{5}                        & \mulh{6}                        & \mulh{7}                        & \mulh{8}                        & \mulh{9}                        & \mulh{10}                        & \mulh{11}                        & \mulh{12}                        & \mulh{13}                        & \mulh{14}                        & \mulh{15}                        & \mulh{16}                        & \mulh{17}                        & \mulh{18}                        & \mulh{19}                        & \mulh{20}                        & \mulh{21}                        & \mulh{22}                        & \mulh{23}                        & \mulh{24}                        & \mulh{25}                        & \mulh{26}                        & \mulh{27}                        & \mulh{28}                        & \mulh{29}                        & \mulh{30}                        \\ \hline
  \mulh{1}        & \cellcolor{black!14}\textbf{1} & 2                              & 3                              & 4                               & 5                               & 6                               & 7                               & 8                               & 9                               & 10                               & 11                               & 12                               & 13                               & 14                               & 15                               & 16                               & 17                               & 18                               & 19                               & 20                               & 21                               & 22                               & 23                               & 24                               & 25                               & 26                               & 27                               & 28                               & 29                               & 30                               \\ \hline
  \mulh{2}        & 2                              & \cellcolor{black!14}\textbf{4} & 6                              & 8                               & 10                              & 12                              & 14                              & 16                              & 18                              & 20                               & 22                               & 24                               & 26                               & 28                               & 30                               & 32                               & 34                               & 36                               & 38                               & 40                               & 42                               & 44                               & 46                               & 48                               & 50                               & 52                               & 54                               & 56                               & 58                               & 60                               \\ \hline
  \mulh{3}        & 3                              & 6                              & \cellcolor{black!14}\textbf{9} & 12                              & 15                              & 18                              & 21                              & 24                              & 27                              & 30                               & 33                               & 36                               & 39                               & 42                               & 45                               & 48                               & 51                               & 54                               & 57                               & 60                               & 63                               & 66                               & 69                               & 72                               & 75                               & 78                               & 81                               & 84                               & 87                               & 90                               \\ \hline
  \mulh{4}        & 4                              & 8                              & 12                             & \cellcolor{black!14}\textbf{16} & 20                              & 24                              & 28                              & 32                              & 36                              & 40                               & 44                               & 48                               & 52                               & 56                               & 60                               & 64                               & 68                               & 72                               & 76                               & 80                               & 84                               & 88                               & 92                               & 96                               & 100                              & 104                              & 108                              & 112                              & 116                              & 120                              \\ \hline
  \mulh{5}        & 5                              & 10                             & 15                             & 20                              & \cellcolor{black!14}\textbf{25} & 30                              & 35                              & 40                              & 45                              & 50                               & 55                               & 60                               & 65                               & 70                               & 75                               & 80                               & 85                               & 90                               & 95                               & 100                              & 105                              & 110                              & 115                              & 120                              & 125                              & 130                              & 135                              & 140                              & 145                              & 150                              \\ \hline
  \mulh{6}        & 6                              & 12                             & 18                             & 24                              & 30                              & \cellcolor{black!14}\textbf{36} & 42                              & 48                              & 54                              & 60                               & 66                               & 72                               & 78                               & 84                               & 90                               & 96                               & 102                              & 108                              & 114                              & 120                              & 126                              & 132                              & 138                              & 144                              & 150                              & 156                              & 162                              & 168                              & 174                              & 180                              \\ \hline
  \mulh{7}        & 7                              & 14                             & 21                             & 28                              & 35                              & 42                              & \cellcolor{black!14}\textbf{49} & 56                              & 63                              & 70                               & 77                               & 84                               & 91                               & 98                               & 105                              & 112                              & 119                              & 126                              & 133                              & 140                              & 147                              & 154                              & 161                              & 168                              & 175                              & 182                              & 189                              & 196                              & 203                              & 210                              \\ \hline
  \mulh{8}        & 8                              & 16                             & 24                             & 32                              & 40                              & 48                              & 56                              & \cellcolor{black!14}\textbf{64} & 72                              & 80                               & 88                               & 96                               & 104                              & 112                              & 120                              & 128                              & 136                              & 144                              & 152                              & 160                              & 168                              & 176                              & 184                              & 192                              & 200                              & 208                              & 216                              & 224                              & 232                              & 240                              \\ \hline
  \mulh{9}        & 9                              & 18                             & 27                             & 36                              & 45                              & 54                              & 63                              & 72                              & \cellcolor{black!14}\textbf{81} & 90                               & 99                               & 108                              & 117                              & 126                              & 135                              & 144                              & 153                              & 162                              & 171                              & 180                              & 189                              & 198                              & 207                              & 216                              & 225                              & 234                              & 243                              & 252                              & 261                              & 270                              \\ \hline
  \mulh{10}       & 10                             & 20                             & 30                             & 40                              & 50                              & 60                              & 70                              & 80                              & 90                              & \cellcolor{black!14}\textbf{100} & 110                              & 120                              & 130                              & 140                              & 150                              & 160                              & 170                              & 180                              & 190                              & 200                              & 210                              & 220                              & 230                              & 240                              & 250                              & 260                              & 270                              & 280                              & 290                              & 300                              \\ \hline
  \mulh{11}       & 11                             & 22                             & 33                             & 44                              & 55                              & 66                              & 77                              & 88                              & 99                              & 110                              & \cellcolor{black!14}\textbf{121} & 132                              & 143                              & 154                              & 165                              & 176                              & 187                              & 198                              & 209                              & 220                              & 231                              & 242                              & 253                              & 264                              & 275                              & 286                              & 297                              & 308                              & 319                              & 330                              \\ \hline
  \mulh{12}       & 12                             & 24                             & 36                             & 48                              & 60                              & 72                              & 84                              & 96                              & 108                             & 120                              & 132                              & \cellcolor{black!14}\textbf{144} & 156                              & 168                              & 180                              & 192                              & 204                              & 216                              & 228                              & 240                              & 252                              & 264                              & 276                              & 288                              & 300                              & 312                              & 324                              & 336                              & 348                              & 360                              \\ \hline
  \mulh{13}       & 13                             & 26                             & 39                             & 52                              & 65                              & 78                              & 91                              & 104                             & 117                             & 130                              & 143                              & 156                              & \cellcolor{black!14}\textbf{169} & 182                              & 195                              & 208                              & 221                              & 234                              & 247                              & 260                              & 273                              & 286                              & 299                              & 312                              & 325                              & 338                              & 351                              & 364                              & 377                              & 390                              \\ \hline
  \mulh{14}       & 14                             & 28                             & 42                             & 56                              & 70                              & 84                              & 98                              & 112                             & 126                             & 140                              & 154                              & 168                              & 182                              & \cellcolor{black!14}\textbf{196} & 210                              & 224                              & 238                              & 252                              & 266                              & 280                              & 294                              & 308                              & 322                              & 336                              & 350                              & 364                              & 378                              & 392                              & 406                              & 420                              \\ \hline
  \mulh{15}       & 15                             & 30                             & 45                             & 60                              & 75                              & 90                              & 105                             & 120                             & 135                             & 150                              & 165                              & 180                              & 195                              & 210                              & \cellcolor{black!14}\textbf{225} & 240                              & 255                              & 270                              & 285                              & 300                              & 315                              & 330                              & 345                              & 360                              & 375                              & 390                              & 405                              & 420                              & 435                              & 450                              \\ \hline
  \mulh{16}       & 16                             & 32                             & 48                             & 64                              & 80                              & 96                              & 112                             & 128                             & 144                             & 160                              & 176                              & 192                              & 208                              & 224                              & 240                              & \cellcolor{black!14}\textbf{256} & 272                              & 288                              & 304                              & 320                              & 336                              & 352                              & 368                              & 384                              & 400                              & 416                              & 432                              & 448                              & 464                              & 480                              \\ \hline
  \mulh{17}       & 17                             & 34                             & 51                             & 68                              & 85                              & 102                             & 119                             & 136                             & 153                             & 170                              & 187                              & 204                              & 221                              & 238                              & 255                              & 272                              & \cellcolor{black!14}\textbf{289} & 306                              & 323                              & 340                              & 357                              & 374                              & 391                              & 408                              & 425                              & 442                              & 459                              & 476                              & 493                              & 510                              \\ \hline
  \mulh{18}       & 18                             & 36                             & 54                             & 72                              & 90                              & 108                             & 126                             & 144                             & 162                             & 180                              & 198                              & 216                              & 234                              & 252                              & 270                              & 288                              & 306                              & \cellcolor{black!14}\textbf{324} & 342                              & 360                              & 378                              & 396                              & 414                              & 432                              & 450                              & 468                              & 486                              & 504                              & 522                              & 540                              \\ \hline
  \mulh{19}       & 19                             & 38                             & 57                             & 76                              & 95                              & 114                             & 133                             & 152                             & 171                             & 190                              & 209                              & 228                              & 247                              & 266                              & 285                              & 304                              & 323                              & 342                              & \cellcolor{black!14}\textbf{361} & 380                              & 399                              & 418                              & 437                              & 456                              & 475                              & 494                              & 513                              & 532                              & 551                              & 570                              \\ \hline
  \mulh{20}       & 20                             & 40                             & 60                             & 80                              & 100                             & 120                             & 140                             & 160                             & 180                             & 200                              & 220                              & 240                              & 260                              & 280                              & 300                              & 320                              & 340                              & 360                              & 380                              & \cellcolor{black!14}\textbf{400} & 420                              & 440                              & 460                              & 480                              & 500                              & 520                              & 540                              & 560                              & 580                              & 600                              \\ \hline
  \mulh{21}       & 21                             & 42                             & 63                             & 84                              & 105                             & 126                             & 147                             & 168                             & 189                             & 210                              & 231                              & 252                              & 273                              & 294                              & 315                              & 336                              & 357                              & 378                              & 399                              & 420                              & \cellcolor{black!14}\textbf{441} & 462                              & 483                              & 504                              & 525                              & 546                              & 567                              & 588                              & 609                              & 630                              \\ \hline
  \mulh{22}       & 22                             & 44                             & 66                             & 88                              & 110                             & 132                             & 154                             & 176                             & 198                             & 220                              & 242                              & 264                              & 286                              & 308                              & 330                              & 352                              & 374                              & 396                              & 418                              & 440                              & 462                              & \cellcolor{black!14}\textbf{484} & 506                              & 528                              & 550                              & 572                              & 594                              & 616                              & 638                              & 660                              \\ \hline
  \mulh{23}       & 23                             & 46                             & 69                             & 92                              & 115                             & 138                             & 161                             & 184                             & 207                             & 230                              & 253                              & 276                              & 299                              & 322                              & 345                              & 368                              & 391                              & 414                              & 437                              & 460                              & 483                              & 506                              & \cellcolor{black!14}\textbf{529} & 552                              & 575                              & 598                              & 621                              & 644                              & 667                              & 690                              \\ \hline
  \mulh{24}       & 24                             & 48                             & 72                             & 96                              & 120                             & 144                             & 168                             & 192                             & 216                             & 240                              & 264                              & 288                              & 312                              & 336                              & 360                              & 384                              & 408                              & 432                              & 456                              & 480                              & 504                              & 528                              & 552                              & \cellcolor{black!14}\textbf{576} & 600                              & 624                              & 648                              & 672                              & 696                              & 720                              \\ \hline
  \mulh{25}       & 25                             & 50                             & 75                             & 100                             & 125                             & 150                             & 175                             & 200                             & 225                             & 250                              & 275                              & 300                              & 325                              & 350                              & 375                              & 400                              & 425                              & 450                              & 475                              & 500                              & 525                              & 550                              & 575                              & 600                              & \cellcolor{black!14}\textbf{625} & 650                              & 675                              & 700                              & 725                              & 750                              \\ \hline
  \mulh{26}       & 26                             & 52                             & 78                             & 104                             & 130                             & 156                             & 182                             & 208                             & 234                             & 260                              & 286                              & 312                              & 338                              & 364                              & 390                              & 416                              & 442                              & 468                              & 494                              & 520                              & 546                              & 572                              & 598                              & 624                              & 650                              & \cellcolor{black!14}\textbf{676} & 702                              & 728                              & 754                              & 780                              \\ \hline
  \mulh{27}       & 27                             & 54                             & 81                             & 108                             & 135                             & 162                             & 189                             & 216                             & 243                             & 270                              & 297                              & 324                              & 351                              & 378                              & 405                              & 432                              & 459                              & 486                              & 513                              & 540                              & 567                              & 594                              & 621                              & 648                              & 675                              & 702                              & \cellcolor{black!14}\textbf{729} & 756                              & 783                              & 810                              \\ \hline
  \mulh{28}       & 28                             & 56                             & 84                             & 112                             & 140                             & 168                             & 196                             & 224                             & 252                             & 280                              & 308                              & 336                              & 364                              & 392                              & 420                              & 448                              & 476                              & 504                              & 532                              & 560                              & 588                              & 616                              & 644                              & 672                              & 700                              & 728                              & 756                              & \cellcolor{black!14}\textbf{784} & 812                              & 840                              \\ \hline
  \mulh{29}       & 29                             & 58                             & 87                             & 116                             & 145                             & 174                             & 203                             & 232                             & 261                             & 290                              & 319                              & 348                              & 377                              & 406                              & 435                              & 464                              & 493                              & 522                              & 551                              & 580                              & 609                              & 638                              & 667                              & 696                              & 725                              & 754                              & 783                              & 812                              & \cellcolor{black!14}\textbf{841} & 870                              \\ \hline
  \mulh{30}       & 30                             & 60                             & 90                             & 120                             & 150                             & 180                             & 210                             & 240                             & 270                             & 300                              & 330                              & 360                              & 390                              & 420                              & 450                              & 480                              & 510                              & 540                              & 570                              & 600                              & 630                              & 660                              & 690                              & 720                              & 750                              & 780                              & 810                              & 840                              & 870                              & \cellcolor{black!14}\textbf{900} \\ \hline
\end{tabular}
\par\medskip
{\raggedright\fontsize{8.6}{10.4}\selectfont
  \textbf{Primes up to 2500} (367 of them): 2 \ \ 3 \ \ 5 \ \ 7 \ \ 11 \ \ 13 \ \ 17 \ \ 19 \ \ 23 \ \ 29 \ \ 31 \ \ 37 \ \ 41 \ \ 43 \ \ 47 \ \ 53 \ \ 59 \ \ 61 \ \ 67 \ \ 71 \ \ 73 \ \ 79 \ \ 83 \ \ 89 \ \ 97 \ \ 101 \ \ 103 \ \ 107 \ \ 109 \ \ 113 \ \ 127 \ \ 131 \ \ 137 \ \ 139 \ \ 149 \ \ 151 \ \ 157 \ \ 163 \ \ 167 \ \ 173 \ \ 179 \ \ 181 \ \ 191 \ \ 193 \ \ 197 \ \ 199 \ \ 211 \ \ 223 \ \ 227 \ \ 229 \ \ 233 \ \ 239 \ \ 241 \ \ 251 \ \ 257 \ \ 263 \ \ 269 \ \ 271 \ \ 277 \ \ 281 \ \ 283 \ \ 293 \ \ 307 \ \ 311 \ \ 313 \ \ 317 \ \ 331 \ \ 337 \ \ 347 \ \ 349 \ \ 353 \ \ 359 \ \ 367 \ \ 373 \ \ 379 \ \ 383 \ \ 389 \ \ 397 \ \ 401 \ \ 409 \ \ 419 \ \ 421 \ \ 431 \ \ 433 \ \ 439 \ \ 443 \ \ 449 \ \ 457 \ \ 461 \ \ 463 \ \ 467 \ \ 479 \ \ 487 \ \ 491 \ \ 499 \ \ 503 \ \ 509 \ \ 521 \ \ 523 \ \ 541 \ \ 547 \ \ 557 \ \ 563 \ \ 569 \ \ 571 \ \ 577 \ \ 587 \ \ 593 \ \ 599 \ \ 601 \ \ 607 \ \ 613 \ \ 617 \ \ 619 \ \ 631 \ \ 641 \ \ 643 \ \ 647 \ \ 653 \ \ 659 \ \ 661 \ \ 673 \ \ 677 \ \ 683 \ \ 691 \ \ 701 \ \ 709 \ \ 719 \ \ 727 \ \ 733 \ \ 739 \ \ 743 \ \ 751 \ \ 757 \ \ 761 \ \ 769 \ \ 773 \ \ 787 \ \ 797 \ \ 809 \ \ 811 \ \ 821 \ \ 823 \ \ 827 \ \ 829 \ \ 839 \ \ 853 \ \ 857 \ \ 859 \ \ 863 \ \ 877 \ \ 881 \ \ 883 \ \ 887 \ \ 907 \ \ 911 \ \ 919 \ \ 929 \ \ 937 \ \ 941 \ \ 947 \ \ 953 \ \ 967 \ \ 971 \ \ 977 \ \ 983 \ \ 991 \ \ 997 \ \ 1009 \ \ 1013 \ \ 1019 \ \ 1021 \ \ 1031 \ \ 1033 \ \ 1039 \ \ 1049 \ \ 1051 \ \ 1061 \ \ 1063 \ \ 1069 \ \ 1087 \ \ 1091 \ \ 1093 \ \ 1097 \ \ 1103 \ \ 1109 \ \ 1117 \ \ 1123 \ \ 1129 \ \ 1151 \ \ 1153 \ \ 1163 \ \ 1171 \ \ 1181 \ \ 1187 \ \ 1193 \ \ 1201 \ \ 1213 \ \ 1217 \ \ 1223 \ \ 1229 \ \ 1231 \ \ 1237 \ \ 1249 \ \ 1259 \ \ 1277 \ \ 1279 \ \ 1283 \ \ 1289 \ \ 1291 \ \ 1297 \ \ 1301 \ \ 1303 \ \ 1307 \ \ 1319 \ \ 1321 \ \ 1327 \ \ 1361 \ \ 1367 \ \ 1373 \ \ 1381 \ \ 1399 \ \ 1409 \ \ 1423 \ \ 1427 \ \ 1429 \ \ 1433 \ \ 1439 \ \ 1447 \ \ 1451 \ \ 1453 \ \ 1459 \ \ 1471 \ \ 1481 \ \ 1483 \ \ 1487 \ \ 1489 \ \ 1493 \ \ 1499 \ \ 1511 \ \ 1523 \ \ 1531 \ \ 1543 \ \ 1549 \ \ 1553 \ \ 1559 \ \ 1567 \ \ 1571 \ \ 1579 \ \ 1583 \ \ 1597 \ \ 1601 \ \ 1607 \ \ 1609 \ \ 1613 \ \ 1619 \ \ 1621 \ \ 1627 \ \ 1637 \ \ 1657 \ \ 1663 \ \ 1667 \ \ 1669 \ \ 1693 \ \ 1697 \ \ 1699 \ \ 1709 \ \ 1721 \ \ 1723 \ \ 1733 \ \ 1741 \ \ 1747 \ \ 1753 \ \ 1759 \ \ 1777 \ \ 1783 \ \ 1787 \ \ 1789 \ \ 1801 \ \ 1811 \ \ 1823 \ \ 1831 \ \ 1847 \ \ 1861 \ \ 1867 \ \ 1871 \ \ 1873 \ \ 1877 \ \ 1879 \ \ 1889 \ \ 1901 \ \ 1907 \ \ 1913 \ \ 1931 \ \ 1933 \ \ 1949 \ \ 1951 \ \ 1973 \ \ 1979 \ \ 1987 \ \ 1993 \ \ 1997 \ \ 1999 \ \ 2003 \ \ 2011 \ \ 2017 \ \ 2027 \ \ 2029 \ \ 2039 \ \ 2053 \ \ 2063 \ \ 2069 \ \ 2081 \ \ 2083 \ \ 2087 \ \ 2089 \ \ 2099 \ \ 2111 \ \ 2113 \ \ 2129 \ \ 2131 \ \ 2137 \ \ 2141 \ \ 2143 \ \ 2153 \ \ 2161 \ \ 2179 \ \ 2203 \ \ 2207 \ \ 2213 \ \ 2221 \ \ 2237 \ \ 2239 \ \ 2243 \ \ 2251 \ \ 2267 \ \ 2269 \ \ 2273 \ \ 2281 \ \ 2287 \ \ 2293 \ \ 2297 \ \ 2309 \ \ 2311 \ \ 2333 \ \ 2339 \ \ 2341 \ \ 2347 \ \ 2351 \ \ 2357 \ \ 2371 \ \ 2377 \ \ 2381 \ \ 2383 \ \ 2389 \ \ 2393 \ \ 2399 \ \ 2411 \ \ 2417 \ \ 2423 \ \ 2437 \ \ 2441 \ \ 2447 \ \ 2459 \ \ 2467 \ \ 2473 \ \ 2477\par}
\endgroup
\clearpage
\loadgeometry{refbase}

\clearpage
\phantomsection
\addcontentsline{toc}{part}{Index}
\printindex

\end{document}
